% Il corpo della tesi è in italiano, anche per la sillabazione.
\selectlanguage{italian}

% Introduzione generale.
\chapter*{Introduzione}
\label{chap:introduzione}
\addcontentsline{toc}{chapter}{Introduzione}
\markboth{Introduzione}{Introduzione}

% Stesura del 2026-11-17, testo unico senza suddivisioni (richiesta dell'utente), riallineata il 2026-11-19 alla nuova
% struttura dei capitoli 5 e 6 (diagnosi dei due colli di bottiglia, distillazione, compressione, stato dell'arte). I numeri
% richiamano soltanto risultati conclusi dei capitoli 3-6; aggiornarli se cambieranno con le valutazioni in corso.

Per riconoscere una parola guardando le labbra di chi la pronuncia non basta osservare la forma della bocca in un istante qualsiasi. Conta come quella forma cambia, in quale ordine si susseguono i movimenti e quanto a lungo durano. Parole come \emph{spend} e \emph{spent} si distinguono per un dettaglio brevissimo alla fine del movimento, e una mano che ruota in senso orario attraversa le stesse posizioni di una che ruota in senso antiorario, soltanto nell'ordine inverso. In compiti come questi l'informazione non sta nelle singole immagini, ma nel modo in cui esse si succedono.

Le event camera sono sensori costruiti attorno a questa idea, in cui al posto di fotogrammi scattati a intervalli regolari ogni pixel segnala per conto proprio i cambiamenti di luminosità nel momento in cui avvengono, con un evento che ne riporta la posizione, l'istante e il verso~\cite{lichtsteiner2008dvs}. Una scena immobile produce pochissimi dati, mentre un movimento genera un flusso di eventi con una risoluzione temporale dell'ordine del microsecondo, in un'ampia gamma di condizioni di illuminazione e con consumi ridotti. Sono proprietà che rendono questi sensori interessanti in robotica e, più in generale, ovunque latenza e consumo contino quanto l'accuratezza~\cite{gallego2022event}. Il prezzo di questa ricchezza è un segnale sparso e irregolare, che una rete neurale, abituata a ricevere tensori di dimensioni fisse, non può elaborare così com'è. Prima di ogni elaborazione il flusso va suddiviso in intervalli di tempo e raccolto in una sequenza di passi, e questa operazione, che ha l'aria di un dettaglio di preparazione dei dati, decide già che cosa la rete potrà vedere. Se un passo dura cinquanta millisecondi in una registrazione e duecento in un'altra, lo stesso ritardo di un passo indica intervalli diversi, e la rete non ha modo di accorgersene.

Candidate naturali per elaborare questi segnali sono le reti neurali spiking, i cui neuroni accumulano l'ingresso nel tempo in un potenziale di membrana e comunicano con impulsi binari. Su un hardware capace di sfruttare la sparsità, l'assenza di attività si traduce in calcolo che non viene eseguito~\cite{li2024sparseacc,du2024spike}. Negli ultimi anni gli Spiking Transformer hanno portato in questo contesto la self-attention, che permette a regioni lontane dell'immagine di influenzarsi a vicenda, e architetture come Spikformer, Spike-driven Transformer e QKFormer hanno ottenuto risultati competitivi sui principali benchmark neuromorfici~\cite{zhou2023spikformer,yao2023spikedriven,zhou2024qkformer}. QKFormer, da cui muove questa tesi, organizza l'elaborazione in una gerarchia di stadi e ricorre a forme di attenzione più economiche di quella completa. Il tempo, però, vi è trattato in modo sorprendentemente semplice, perché l'attenzione opera su un passo alla volta e il passato raggiunge il presente soltanto attraverso il potenziale dei neuroni, uno stato che decade e riassume l'intera storia in un unico valore. Misure recenti confermano che, nelle attenzioni spiking più diffuse, l'uscita di un passo dipende quasi soltanto dall'ingresso dello stesso passo~\cite{shi2026multidelay}.

L'accesso al passato si può rendere esplicito aprendo l'attenzione a blocchi di passi consecutivi~\cite{lee2025spiking}, oppure mescolando le feature con loro copie ritardate~\cite{shi2026multidelay}, ma in entrambi i casi la rete deve conservare informazione da un passo all'altro, e questa memoria occupa spazio e richiede accessi che il semplice conteggio dei parametri non lascia vedere. Il problema passa in secondo piano in un modello grande, ma diventa centrale in uno piccolo, pensato per dispositivi con risorse limitate come le FPGA~\cite{karamimanesh2025fpga}. Con circa mezzo milione di parametri ogni scelta va giustificata due volte, per ciò che aggiunge alle prestazioni e per ciò che costa in moltiplicazioni, in stato da conservare e in dati da spostare. Anche il numero dei parametri, del resto, racconta soltanto una parte della storia, perché nella baseline da cui parte questa tesi il primo patch embedding contiene meno di un quinto dei pesi ma concentra la maggior parte dello stato da conservare, lavorando su mappe ad alta risoluzione, e i due patch embedding insieme producono tre quarti delle moltiplicazioni dell'intera rete. Un modulo con pochi pesi può rivelarsi costoso, se costringe a ricordare molte feature, e uno con più pesi può costare meno, se opera dove le posizioni sono poche.

È in questo regime che si colloca la tesi, guidata dalla domanda di come spendere un budget compatto perché uno Spiking Transformer rappresenti le relazioni temporali che contano nei dati a eventi, senza perdere di vista il costo della memoria. Non si tratta di costruire il modello più accurato a qualunque prezzo, né quello con il minor numero di operazioni, ma di capire dove l'informazione temporale valga le risorse necessarie a rappresentarla. Posta così, la domanda si scompone in questioni più concrete, nell'ordine in cui si presentano a chi progetta una rete di questo tipo. Occorre capire anzitutto dove convenga collocare la poca capacità disponibile, e poi che cosa debba rappresentare un passo, cioè quanto pesi il modo in cui il flusso di eventi viene suddiviso nel tempo. Con un passo dal significato chiaro ci si può chiedere se un confronto esplicito tra passato e presente renda la memoria più utile di un semplice richiamo, e a quale prezzo. Resta infine da vedere come cambino questi benefici quando il budget cresce o quando cambia il modo di addestrare la rete.

Il metodo segue da vicino questa scomposizione, e ogni passo del lavoro modifica un solo fattore rispetto a un riferimento fissato, mentre dati, ricetta di addestramento e regole di valutazione restano gli stessi, così che ogni differenza abbia un'origine riconoscibile. I modelli sono addestrati su quattro dataset neuromorfici dalle caratteristiche molto diverse, DVS-Lip per la lettura labiale~\cite{tan2022dvslip}, DVS128 Gesture per il riconoscimento di gesti~\cite{amir2017dvsgesture}, CIFAR10-DVS per la classificazione di oggetti~\cite{li2017cifar10dvs} e THU-EACT-50-CHL per il riconoscimento di azioni~\cite{gao2023action}. Le prestazioni sono misurate con la Macro-F1 su una validazione di sviluppo, e due modelli vengono confrontati campione per campione, per separare le differenze reali da quelle che il caso può produrre. Il test ufficiale, valutato soltanto a posteriori sui modelli già scelti, serve a verificarli e a confrontare i modelli finali con la letteratura, senza entrare in alcuna scelta. Il costo è descritto da conteggi analitici di parametri, operazioni e stato persistente, accompagnati da una stima aritmetica dell'energia. Il lavoro non esegue i modelli su hardware dedicato, e nessun risparmio di energia o di latenza viene presentato come misurato.

I quattro dataset non svolgono tutti lo stesso ruolo, perché su DVS128 Gesture e CIFAR10-DVS la struttura temporale conta meno, o viene già catturata da una buona rappresentazione dell'ingresso, e questi compiti servono soprattutto a verificare che una modifica non danneggi ciò che già funziona. Su DVS-Lip e THU-EACT-50-CHL, invece, l'ordine e la durata dei movimenti sono parte essenziale del riconoscimento, ed è su questi due dataset che la memoria temporale viene studiata più a fondo. Nemmeno gli esperimenti rispondono tutti alla stessa domanda, visto che riaddestrare una rete senza un componente dice se quel componente serva, spegnerlo in una rete già addestrata dice soltanto quanto la soluzione appresa vi si appoggi, e un controllo di pari capacità dice se un guadagno venga dalla struttura o dai parametri aggiunti. Tenere distinte queste letture è ciò che permette di non attribuire a un modulo più di quanto i dati consentano.

Il percorso comincia dalla baseline, la versione compatta di QKFormer descritta nel capitolo~\ref{chap:baseline_metodologia} insieme al protocollo sperimentale. Il capitolo~\ref{chap:riallocazione_budget} ne parte con una diagnosi, che mostra come quasi tutto il costo si concentri nei primi strati mentre la memoria temporale resta affidata ai soli neuroni, e ridistribuisce la capacità di conseguenza. Un frontend piramidale, un mixer causale che legge il passato a più ritardi e una convoluzione depthwise al posto dell'attenzione del primo stadio portano la Macro-F1 di DVS-Lip da 44 a 55 punti, con appena 320 parametri in più e con meno stato da conservare. Ne nasce il modello di riferimento per il resto della tesi, insieme a un'indicazione che il lavoro ritroverà più volte, cioè che con un budget così piccolo conta il luogo in cui la capacità viene spesa, prima ancora della sua quantità.

Il capitolo~\ref{chap:clock_input} tiene ferma questa architettura e cambia soltanto il clock dell'input. Il confronto tra tre modi di suddividere gli eventi mostra che la scelta pesa sulle prestazioni quanto una modifica della rete, e che una rappresentazione a lunghezza variabile, in cui ogni passo ha la stessa durata fisica in tutto il dataset, migliora nettamente DVS128 Gesture e THU-EACT-50-CHL. Un controllo con un numero fisso di passi indica che eliminare i tratti vuoti della registrazione non basta, e che dove le durate variano molto conta anche che un passo abbia lo stesso significato fisico in tutte le registrazioni. Con un passo di durata nota, i ritardi della rete diventano distanze precise nel passato.

Su questa base il capitolo~\ref{chap:memoria_relazionale} sviluppa il contributo centrale, partendo dall'osservazione che il mixer a ritardi del modello di riferimento richiama il passato e lo somma al presente, ma non li mette mai a confronto. La memoria proposta, un \emph{interval-relational mixer} (IRM) nella sua forma antisimmetrica, riassume il passato in bande causali a distanze fisiche fisse, tiene conto di quanta storia esista davvero all'inizio di ogni sequenza e confronta il passato con il presente attraverso un termine bilineare a basso rango. Il termine cambia segno quando i due istanti si scambiano, come nel rilevatore di correlazione con cui Hassenstein e Reichardt spiegarono la percezione del movimento negli insetti~\cite{hassenstein1956systemtheoretische}, e risponde quindi al verso del cambiamento più che alla sua semplice presenza. Su THU-EACT-50-CHL l'IRM migliora il modello di riferimento di quasi otto punti di Macro-F1 in validazione e di quasi undici sul test ufficiale. I controlli mostrano che il guadagno non viene dai parametri aggiunti, ma dal confronto applicato a un passato collocato a distanze precise, che né tracce esponenziali né un passato più lungo o un orologio appreso riescono a sostituire. Su DVS-Lip, invece, lo stesso modulo sposta il risultato di poco. Con mezzo milione di parametri la rete per THU-EACT-50-CHL non riesce ad adattarsi nemmeno ai propri dati di addestramento, e la memoria relazionale le fornisce la capacità temporale che le manca, mentre quella per DVS-Lip impara quasi tutto il training ma generalizza poco, e modelli di accuratezza quasi identica sbagliano clip in buona parte diverse. Il limite di questo secondo compito non è la memoria, ma la varianza della soluzione.

Seguendo questa diagnosi, nel capitolo~\ref{chap:scalabilita} la distillazione raccoglie su DVS-Lip in un solo modello compatto ciò che varianti diverse sanno separatamente, e ne aumenta la Macro-F1 di oltre cinque punti senza alcun costo in inferenza, mentre su THU-EACT-50-CHL la stessa ricetta peggiora il risultato e la memoria relazionale resta la risposta. Entrambi i compiti migliorano con la capacità, e il margine dell'IRM si riduce man mano che il budget cresce, a conferma che la memoria relazionale conta soprattutto dove le risorse sono scarse. Una compressione eseguita dopo l'addestramento, senza riaddestrare la rete, rende poi l'IRM più economico del mixer che sostituisce a ogni larghezza. Su THU-EACT-50-CHL l'intera catena di scelte porta oltre lo stato dell'arte, e con meno di mezzo milione di parametri la rete supera già di quasi diciannove punti la migliore rete spiking pubblicata, e con uno o due milioni supera tutti i metodi che non usano dati esterni. Lo stesso capitolo mette in luce anche ciò che nessuna delle leve provate riesce a cambiare, a cominciare dall'incapacità della rete di riconoscere un movimento eseguito a una velocità molto diversa da quelle viste in addestramento.

Prodotti bilineari a basso rango, correlatori di movimento, decomposizioni ai valori singolari e distillazione sono strumenti noti, e il lavoro non rivendica un nuovo tipo di operatore. Il suo contributo sta nella combinazione di questi strumenti in una soluzione specifica per uno Spiking Transformer compatto, nella catena di confronti controllati che porta a quella soluzione e nell'evidenza che ne misura l'utilità, il costo e i confini.

Prima di questo percorso, il capitolo~\ref{chap:fondamenti_stato_arte} raccoglie i fondamenti necessari a seguirlo, dalla rappresentazione degli eventi alla dinamica dei neuroni spiking, dagli Spiking Transformer alle strategie di modellazione temporale esplicita, fino alle grandezze con cui se ne misura il costo. Alla fine, le conclusioni riprendono le domande iniziali alla luce dei risultati, ne tracciano i limiti e indicano le strade che restano aperte, a partire da quella di una realizzazione su FPGA.


% Il background diventa il primo capitolo numerato.
% chapter1.tex conserva il precedente capitolo di contesto, escluso dalla compilazione.
\chapter{Fondamenti teorici e stato dell'arte}
\label{chap:fondamenti_stato_arte}

Questo capitolo raccoglie i fondamenti necessari a comprendere la progettazione di uno Spiking Transformer compatto per dati neuromorfici. L'esposizione segue il cammino dell'informazione, dal sensore alla sua rappresentazione discreta, dalla memoria del singolo neurone alle relazioni tra feature, fino al costo computazionale e di memoria che ogni scelta comporta.

La sezione~\ref{sec:rappresentazione_event_camera} affronta il primo passaggio obbligato, la conversione del flusso asincrono di eventi in tensori regolari, e mostra come questa operazione, apparentemente preliminare, costituisca già una forma di modellazione. La sezione~\ref{sec:dinamica_neuronale_memoria} descrive la memoria intrinseca dei neuroni spiking e i limiti di uno stato che riassume ricorsivamente il passato. Le sezioni~\ref{sec:snn_spiking_transformer} e~\ref{sec:qkformer_riferimento} ricostruiscono l'incontro tra attenzione e computazione spiking fino a QKFormer, l'architettura da cui deriva la baseline compatta, distinguendo la memoria dei neuroni dalle interazioni esplicite tra istanti diversi. La sezione~\ref{sec:strategie_modellazione_temporale} esamina due modi di rendere esplicito l'accesso alla storia, l'attenzione spazio-temporale e il mixing mediante ritardi, mentre la sezione~\ref{sec:efficienza_neuromorfica_vincoli} introduce le grandezze con cui se ne misura il costo.

Più che una rassegna esaustiva, il capitolo raccoglie contributi selezionati perché orientano una scelta progettuale compiuta nel seguito del lavoro, e li presenta in funzione di essa.


% ==============================================================================
%  Rappresentazione dell'informazione nelle event camera
% ==============================================================================

\section{Event camera e rappresentazione degli eventi}
\label{sec:rappresentazione_event_camera}

Una fotocamera convenzionale scatta fotografie a intervalli regolari, per esempio trenta al secondo, e ogni fotografia registra la luminosità di tutti i pixel, anche di quelli in cui non è cambiato nulla. Una event camera funziona in modo opposto, e ogni pixel resta in silenzio finché la luce che riceve non cambia abbastanza, e in quel momento segnala il cambiamento con un \emph{evento}. Una scena immobile non produce quasi nulla, mentre una mano che si muove produce una scia di eventi lungo i propri contorni, nell'istante esatto in cui si sposta.

Per elaborare questo segnale con una rete neurale, che lavora su tensori di dimensioni fisse, occorre raccogliere gli eventi in una struttura regolare. Questa sezione descrive prima il segnale prodotto dal sensore e poi le scelte necessarie per trasformarlo in un tensore. Sono scelte tutt'altro che neutre, perché decidono quali informazioni arrivano alla rete e quali vanno perse prima ancora che l'elaborazione cominci.

\subsection{Dal contrasto luminoso al flusso di eventi}
\label{subsec:contrasto_eventi}

Nel modello ideale di un Dynamic Vision Sensor (DVS) ogni pixel lavora indipendentemente dagli altri e osserva il logaritmo dell'intensità luminosa che lo colpisce, $L(x,y,t)=\log I(x,y,t)$~\cite{lichtsteiner2008dvs}. Il pixel conserva un valore di riferimento, cioè il livello di $L$ al momento del suo ultimo evento, e lo confronta di continuo con il valore attuale. Quando la differenza raggiunge una soglia $C>0$, verso l'alto o verso il basso, il pixel emette un evento e aggiorna il riferimento. In formule, l'evento $i$-esimo viene emesso nell'istante $t_i$ in cui
\begin{equation}
\label{eq:soglia_contrasto}
L(x_i,y_i,t_i)-L(x_i,y_i,t_i^{\mathrm{ref}})=p_i\,C,
\end{equation}
dove $t_i^{\mathrm{ref}}$ è l'istante a cui risale il riferimento, cioè l'evento precedente dello stesso pixel o la sua inizializzazione, e la polarità $p_i$ vale $+1$ se la luce è aumentata e $-1$ se è diminuita~\cite{gallego2022event}. Ogni evento è quindi descritto da quattro numeri,
\begin{equation}
\label{eq:evento_quadrupla}
e_i=(x_i,y_i,t_i,p_i),
\end{equation}
la colonna e la riga del pixel, l'istante e la polarità. Un evento con $x_i=64$, $y_i=30$, $t_i=1\,250\,400\,\mu\mathrm{s}$ e $p_i=+1$, per esempio, dice che il pixel della colonna 64 e della riga 30 è diventato più luminoso $1{,}25$ secondi dopo l'inizio della registrazione.

La figura~\ref{fig:pixel_event_generation} mostra il meccanismo su un singolo pixel, attraversato prima da un oggetto chiaro e poi da uno scuro. Finché la luce cresce, ogni incremento pari a $C$ produce un evento positivo, quando cala ogni decremento pari a $C$ ne produce uno negativo. La soglia agisce sulla variazione accumulata dall'ultimo evento, non sulla velocità del cambiamento, sicché nel modello ideale una variazione lenta e una rapida della stessa entità producono lo stesso numero di eventi, soltanto distribuiti in tempi diversi. Le oscillazioni più piccole della soglia, invece, non producono nulla.

\begin{figure}[tbp]
    \centering
    \includegraphics[width=\textwidth]{images/cap1/cap1_generazione_eventi.pdf}
    \caption[Generazione degli eventi in un pixel ideale]{Generazione degli eventi in un pixel ideale, su un segnale sintetico con soglia $C=0{,}2$. In alto, il logaritmo dell'intensità (linea continua), il valore di riferimento del pixel (tratteggio) e la fascia di ampiezza $\pm C$ attorno a esso. Ogni volta che la curva esce dalla fascia il pixel emette un evento, positivo (triangolo rosso) o negativo (triangolo blu), e il riferimento si sposta di $C$. In basso, la sequenza di eventi che il pixel trasmette. L'oscillazione centrale, più piccola della soglia, non produce alcun evento.}
    \label{fig:pixel_event_generation}
\end{figure}

Perché proprio il logaritmo? Perché la differenza tra due logaritmi è il logaritmo del rapporto,
\begin{equation}
\label{eq:rapporto_luminanza}
\Delta L=\log\frac{I(x_i,y_i,t_i)}{I(x_i,y_i,t_i^{\mathrm{ref}})},
\end{equation}
e il pixel reagisce quindi alle variazioni relative della luce, non a quelle assolute. Con una soglia $C=0{,}2$, scelta qui a titolo illustrativo, un evento positivo richiede un aumento della luce di circa il $22\%$, poiché $e^{0{,}2}\approx1{,}22$. Vale allo stesso modo per un pixel in piena luce che passa da $1000$ a $1220$ unità e per uno in penombra che passa da $10$ a $12{,}2$. La stessa scena produce così eventi simili in condizioni di illuminazione molto diverse, ed è una delle ragioni dell'ampia gamma dinamica di questi sensori~\cite{gallego2022event}.

La polarità, va precisato, dice soltanto se in quel pixel la luce è aumentata o diminuita, non in quale direzione si muova qualcosa. Si pensi a una barra scura che scorre verso destra su uno sfondo chiaro. I pixel raggiunti dal suo bordo anteriore si scuriscono ed emettono eventi negativi, quelli lasciati dal bordo posteriore tornano chiari ed emettono eventi positivi. Se la barra scorresse verso sinistra, ogni pixel emetterebbe esattamente la stessa sequenza, e la direzione si potrebbe leggere soltanto nell'ordine con cui si attivano pixel vicini. Allo stesso modo, un evento non implica necessariamente un movimento, perché anche una lampada che sfarfalla fa cambiare la luce di molti pixel insieme. Il sensore reale, infine, si allontana dal modello ideale in vari modi. Le soglie positiva e negativa possono differire tra loro e variare da pixel a pixel, e il rumore, la banda limitata del fotorecettore e un breve periodo refrattario dopo ogni evento aggiungono altri scostamenti~\cite{gallego2022event}.

Una registrazione è dunque una lunga sequenza di eventi in ordine di tempo,
\begin{equation}
\label{eq:stream_completo}
\mathcal E=\{e_i\}_{i=1}^{N},\qquad t_i\leq t_{i+1}.
\end{equation}
Il numero $N$ degli eventi non dipende dalla durata ma da quanto accade nella scena. Nei dataset usati in questo lavoro, per esempio, una parola pronunciata in circa un secondo produce in media poco più di diecimila eventi, mentre un'immagine di CIFAR10-DVS fatta muovere davanti al sensore per un tempo simile ne produce circa duecentomila (sezione~\ref{sec:dataset}). Anche all'interno della stessa registrazione possono alternarsi lunghi silenzi e raffiche improvvise, e chi elabora il segnale deve saper gestire gli uni e le altre.

\subsection{Discretizzazione e significato del clock}
\label{subsec:discretizzazione_clock}

Una rete neurale spiking elabora l'ingresso a passi discreti. Per presentarle un flusso di eventi bisogna quindi decidere che cosa entri in ciascun passo, cioè come tagliare il flusso in pezzi. La regola che associa ogni passo a un intervallo di tempo, o più in generale a un gruppo di eventi, prende il nome di \emph{clock} dell'input. Le regole possibili sono diverse, e la figura~\ref{fig:partizioni_flusso} le confronta sullo stesso flusso.

La più semplice divide il tempo in intervalli consecutivi della stessa durata $\Delta t$, detti bin, a partire da un'origine $t_0$. Il bin $k$-esimo raccoglie gli eventi caduti nel suo intervallo,
\begin{equation}
\label{eq:finestre_fisse}
\mathcal E_k=\bigl\{e_i\in\mathcal E\ \big|\ t_0+k\Delta t\leq t_i<t_0+(k+1)\Delta t\bigr\}.
\end{equation}
Con $t_0=0$ e $\Delta t=50\,\mathrm{ms}$, un evento all'istante $t_i=123\,\mathrm{ms}$ finisce nel bin $k=2$, quello che va da $100$ a $150\,\mathrm{ms}$. Il bordo sinistro è incluso e quello destro escluso, così che un evento che cade esattamente a $150\,\mathrm{ms}$ appartenga al bin successivo e mai a due bin insieme. In questo schema ogni passo dura lo stesso tempo, mentre il numero di eventi che contiene cambia liberamente da un passo all'altro.

Fissata la durata dei bin, resta da decidere quanti passi elaborare. La soluzione più comune è una finestra di durata prefissata, $D_{\mathrm{win}}=T\Delta t$, che produce sempre $T$ passi. Con bin da $50\,\mathrm{ms}$ e una finestra di $2\,\mathrm{s}$ i passi sono 40. Una registrazione di $1{,}2\,\mathrm{s}$ ne riempie i primi 24 e lascia vuoti gli ultimi 16, mentre di una registrazione di $2{,}5\,\mathrm{s}$ vanno persi gli ultimi $0{,}5\,\mathrm{s}$.

In alternativa si può conservare l'intera registrazione, dal primo evento, che fa da origine $t_0$, all'ultimo, che cade all'istante $t_{\mathrm{last}}$, mantenendo bin della stessa durata. Il numero di passi diventa allora
\begin{equation}
\label{eq:lunghezza_clock_fisico}
T=\left\lfloor\frac{t_{\mathrm{last}}-t_0}{\Delta t}\right\rfloor+1.
\end{equation}
La parte intera conta i bin completamente trascorsi prima dell'ultimo evento, e l'unità aggiunta fa posto al bin che lo contiene. Con bin da $50\,\mathrm{ms}$ e un ultimo evento a $1199\,\mathrm{ms}$ si ottengono 24 passi. Se l'ultimo evento cadesse esattamente a $1200\,\mathrm{ms}$, sul bordo sinistro di un nuovo bin, i passi sarebbero 25, ed è proprio l'unità aggiunta a garantire che quell'evento trovi posto. In questo schema la durata dei bin è la stessa per tutte le registrazioni, e a cambiare è la lunghezza della sequenza.

Una terza possibilità fa il contrario, fissa il numero di passi $T$ e adatta la durata dei bin a ciascuna registrazione $n$, di durata $D_n$,
\begin{equation}
\label{eq:delta_t_rapporto}
\Delta t_n=\frac{D_n}{T}.
\end{equation}
Con $T=40$, un gesto di $4\,\mathrm{s}$ viene diviso in bin da $100\,\mathrm{ms}$ e uno di $8\,\mathrm{s}$ in bin da $200\,\mathrm{ms}$. Il clock cambia così da un campione all'altro. Una distanza di cinque passi corrisponde a mezzo secondo nel primo gesto e a un secondo intero nel secondo, e la rete non ha modo di saperlo. Fissare il numero dei bin, in altre parole, significa non fissare la risoluzione temporale.

Si può infine rinunciare del tutto a misurare il tempo e contare gli eventi. Una prima variante chiude un gruppo ogni $M$ eventi, nell'ordine in cui arrivano. Una seconda divide l'intera registrazione in $T$ gruppi che contengono lo stesso numero di eventi, a meno di uno,
\begin{equation}
\label{eq:gruppi_eventi_bilanciati}
n_k\in\left\{\left\lfloor\frac NT\right\rfloor,
                 \left\lceil\frac NT\right\rceil\right\},
\qquad \sum_{k=1}^{T}n_k=N,
\end{equation}
dove $n_k$ è il numero di eventi del gruppo $k$. Con $N=10\,003$ eventi e $T=40$ gruppi, per esempio, tre gruppi ricevono 251 eventi e gli altri 37 ne ricevono 250. In questo schema ogni passo contiene la stessa quantità di attività, ma la sua durata dipende da quanto è intensa la scena. Una raffica concentrata in un decimo di secondo può riempire molti gruppi, mentre un secondo di quiete può starne in uno solo. La seconda variante richiede di conoscere in anticipo il numero totale $N$ degli eventi, e si può applicare soltanto a registrazioni già concluse, cioè offline. La prima non ha questo limite, perché ogni gruppo si chiude non appena raggiunge $M$ eventi. In entrambi i casi, però, una distanza misurata in gruppi non corrisponde a un intervallo di tempo fisso.

\begin{figure}[tbp]
    \centering
    \includegraphics[width=\textwidth]{images/cap1/cap1_partizioni.pdf}
    \caption[Quattro suddivisioni di uno stesso flusso di eventi]{Quattro modi di suddividere lo stesso flusso sintetico, che contiene una raffica di attività tra $0{,}35$ e $0{,}55\,\mathrm{s}$ e termina a circa $1{,}17\,\mathrm{s}$. In alto, la frequenza degli eventi. Sotto, i passi prodotti da ciascuna regola, colorati in base al numero di eventi che contengono. La finestra fissa lascia una coda di passi vuoti (tratteggiati), la registrazione intera termina con l'ultimo evento, il numero di passi fisso allarga i bin per coprire la durata, la partizione a eventi costanti produce passi con lo stesso contenuto ma di durata molto diversa, stretti durante la raffica e larghi altrove.}
    \label{fig:partizioni_flusso}
\end{figure}

Quando i passi seguono il numero di eventi, la loro durata non è più fissata in anticipo ma diventa una proprietà di ciascun gruppo, misurabile come intervallo (\emph{span}) tra il suo primo e il suo ultimo evento. Se la rete riceve soltanto i conteggi, però, questa durata le resta invisibile. Un gruppo di 250 eventi raccolti in $5\,\mathrm{ms}$ e uno raccolto in mezzo secondo occupano entrambi esattamente un passo.

La scelta del clock non è dunque un dettaglio di preprocessing. Da essa dipende il significato fisico di ogni ritardo espresso in passi e di ogni costante di decadimento della rete. Dimezzare $\Delta t$ su una stessa registrazione raddoppia il numero dei passi, dimezza in media i conteggi di ciascun bin e fa sì che un neurone che perde una frazione fissa del proprio potenziale a ogni passo dimentichi due volte più in fretta in tempo reale. I protocolli adottati in questo lavoro sono descritti nel capitolo~\ref{chap:baseline_metodologia} e messi a confronto nel capitolo~\ref{chap:clock_input}.

\subsection{Conteggi e informazione intra-bin}
\label{subsec:conteggi_intrabin}

Una volta suddiviso il flusso, bisogna decidere che cosa conservare di ciascun bin. La rappresentazione più semplice conta, per ogni pixel e per ogni bin, quanti eventi di ciascuna polarità vi sono caduti,
\begin{equation}
\label{eq:event_frame_accumulo}
F_k^p(x,y)=\sum_{e_i\in\mathcal E_k}
\mathbf 1[x_i=x,\ y_i=y,\ p_i=p],
\qquad F_k\in\mathbb N_0^{2\times H\times W}.
\end{equation}
L'indicatore $\mathbf 1[\cdot]$ vale 1 quando la condizione tra parentesi è vera e 0 altrimenti, quindi la somma conta semplicemente gli eventi del bin $k$ caduti nel pixel $(x,y)$ con polarità $p$. Il risultato è un'immagine a due canali, uno per gli eventi positivi e uno per quelli negativi, spesso chiamata \emph{event frame} (figura~\ref{fig:confronto_rappresentazioni_discrete}a). Se in un bin un pixel riceve tre eventi positivi e uno negativo, il canale positivo vale 3 e quello negativo 1.

Esistono anche versioni più povere della stessa idea, come quella binaria, $\mathbf 1[F_k^p(x,y)>0]$, che registra soltanto se il pixel è stato attivo, e nell'esempio darebbe 1 in entrambi i canali. Una versione a un solo canale somma gli eventi con il loro segno, e darebbe $3-1=2$. Quest'ultima ha un difetto evidente, perché due eventi positivi e due negativi nello stesso pixel si annullano e producono zero, come se non fosse accaduto nulla. Tenere separate le polarità evita il problema.

Tutte queste versioni condividono però un limite più sottile, perché una somma non dipende dall'ordine degli addendi e i conteggi di un bin non cambiano se gli eventi al suo interno vengono rimescolati. Si pensi a un oggetto chiaro che attraversa rapidamente un pixel scuro. Il pixel prima si illumina e poi torna scuro, ed emette un evento positivo seguito da uno negativo. Un oggetto scuro che attraversa un pixel chiaro produce la sequenza opposta. Se entrambi gli eventi cadono nello stesso bin, i due casi danno esattamente gli stessi conteggi, uno per canale, e nessuna rete, per quanto sofisticata, potrà distinguerli sulla base di quel solo bin. La rappresentazione fissa un limite all'informazione disponibile, che va tenuto distinto dalla capacità della rete di sfruttare quella che rimane.

La \emph{voxel grid} recupera parte dell'informazione temporale perduta~\cite{zhu2019unsupervised}. Si considera un pacchetto di eventi consecutivi, elaborato come un'unica unità, e se ne riscalano i timestamp in modo che il primo evento valga 0 e l'ultimo $B-1$,
\begin{equation}
\label{eq:timestamp_normalizzato}
t_i^*=(B-1)\frac{t_i-t_{\min}}{t_{\max}-t_{\min}},
\end{equation}
dove $B$ è il numero di bin temporali desiderati e il pacchetto deve avere durata non nulla, $t_{\max}>t_{\min}$. Ogni evento viene poi distribuito tra i due bin più vicini alla sua posizione $t_i^*$,
\begin{equation}
\label{eq:voxel_grid_interpolazione}
\mathcal V(b,x,y)=\sum_i p_i
\max\bigl(0,\,1-|b-t_i^*|\bigr)
\mathbf 1[x_i=x,\ y_i=y],
\qquad b\in\{0,\ldots,B-1\}.
\end{equation}
Il peso $\max(0, 1-|b-t_i^*|)$ vale 1 quando l'evento cade esattamente sul bin $b$, decresce in modo lineare con la distanza e si annulla a un bin di distanza. Con $B=5$, un evento positivo con $t_i^*=2{,}3$ aggiunge $0{,}7$ al bin 2 e $0{,}3$ al bin 3, e niente agli altri. L'evento non viene attribuito tutto a un bin o tutto all'altro, e la sua posizione temporale sopravvive, sfumata, nella ripartizione dei pesi (figura~\ref{fig:confronto_rappresentazioni_discrete}b). La formulazione completa interpola anche nello spazio, quando la correzione della distorsione ottica produce coordinate non intere~\cite{zhu2019unsupervised}.

La somma usa la polarità con il suo segno, e può quindi produrre le stesse cancellazioni dell'event frame a un canale, che volumi separati per eventi positivi e negativi eviterebbero. La scala temporale, inoltre, è relativa a ciascun pacchetto, e con $B=5$ un intervallo di $10\,\mathrm{ms}$ occupa meno di mezzo bin in un pacchetto lungo $100\,\mathrm{ms}$ e ben due bin in uno lungo $20\,\mathrm{ms}$, e non va confuso con un clock fisico comune a tutti i campioni.

La \emph{time surface} adotta un punto di vista diverso. Invece di contare gli eventi, descrive quanto sia recente l'attività di ciascun pixel. Nella forma più semplice,
\begin{equation}
\label{eq:time_surface}
S_p(x,y,t)=
\exp\left[-\frac{t-t_{\mathrm{last}}^p(x,y)}{\tau_s}\right],
\end{equation}
dove $t_{\mathrm{last}}^p(x,y)$ è l'istante dell'ultimo evento di polarità $p$ arrivato al pixel prima di $t$, e il valore è zero dove non è ancora arrivato nulla. Un pixel appena attivato vale circa 1, e il suo valore decade poi in modo esponenziale con costante $\tau_s$. Con $\tau_s=50\,\mathrm{ms}$, un pixel attivato $10\,\mathrm{ms}$ prima vale circa $0{,}82$, uno attivato $100\,\mathrm{ms}$ prima circa $0{,}14$. Il risultato somiglia alla scia di una cometa, intensa dove il movimento è appena passato e sempre più tenue alle sue spalle (figura~\ref{fig:confronto_rappresentazioni_discrete}c). È la variante elementare di una famiglia più ampia. HATS, per esempio, somma i contributi esponenziali di più eventi in un intorno spazio-temporale e li aggrega in descrittori~\cite{sironi2018hats}.

\begin{figure}[tbp]
    \centering
    \includegraphics[width=\textwidth]{images/cap1/cap1_rappresentazioni.pdf}
    \caption[Event frame, voxel grid e time surface a confronto]{Tre rappresentazioni degli stessi $100\,\mathrm{ms}$ di una registrazione di DVS128 Gesture, durante una rotazione del braccio sinistro (circa $9\,600$ eventi). (a)~Event frame, con i conteggi degli eventi positivi e negativi in due canali separati. (b)~Voxel grid con $B=5$ bin temporali, in cui la distribuzione dell'attività cambia da un bin all'altro, con i valori positivi in rosso e quelli negativi in blu. (c)~Time surface dell'ultimo evento di entrambe le polarità, valutata alla fine dell'intervallo, in cui la posizione più recente del movimento appare come una scia.}
    \label{fig:confronto_rappresentazioni_discrete}
\end{figure}

Event frame, voxel grid e time surface riassumono dunque aspetti diversi della stessa storia, la quantità dell'attività, la sua collocazione nel tempo e la sua recenza. La tabella~\ref{tab:confronto_rappresentazioni} ne riassume i compromessi. Una rappresentazione più ricca costa di più da costruire e da conservare, e la sua utilità dipende in ultima analisi dalle operazioni che il modello è in grado di applicarvi.

\begin{table}[tbp]
\centering
\small
\renewcommand{\arraystretch}{1.2}
\caption[Rappresentazioni degli eventi a confronto]{Informazione conservata e principali limiti delle rappresentazioni considerate.}
\label{tab:confronto_rappresentazioni}
\begin{tabularx}{\textwidth}{@{}>{\raggedright\arraybackslash}p{2.25cm} >{\raggedright\arraybackslash}X >{\raggedright\arraybackslash}X >{\raggedright\arraybackslash}X@{}}
\hline
\textbf{Formato} & \textbf{Informazione conservata} & \textbf{Descrizione del tempo} & \textbf{Limite principale} \\
\hline
Stream raw & Posizione, polarità e timestamp dei singoli eventi & Istanti originari & Sequenza irregolare e carico variabile \\
Conteggi ON/OFF & Molteplicità per pixel e polarità & Appartenenza al bin & Ordine intra-bin eliminato \\
Voxel grid & Accumuli su più bin temporali & Posizione temporale discretizzata o interpolata & Aggregazione e possibile cancellazione nella versione firmata \\
Time surface dell'ultimo evento & Recenza per pixel e polarità & Tempo trascorso dall'ultima attivazione & Molteplicità e storia anteriore non rappresentate separatamente \\
\hline
\end{tabularx}
\end{table}

\subsection{Tempo valido, silenzi e padding}
\label{subsec:tempo_valido_padding}

Quando le sequenze hanno lunghezze diverse sorge un problema pratico. Le reti vengono addestrate su batch, cioè gruppi di campioni elaborati insieme in un unico tensore, e un tensore richiede che tutte le sequenze abbiano la stessa lunghezza. Le sequenze più brevi vengono quindi allungate con passi fittizi, il cosiddetto \emph{padding}, fino alla lunghezza della più lunga, $T_{\max}$. Per tenere traccia di quali passi siano veri, a ogni campione $n$ di lunghezza $T_n$ si associa una maschera
\begin{equation}
\label{eq:maschera_passi_validi}
m_{k,n}=\mathbf 1[k<T_n],
\qquad k=0,\ldots,T_{\max}-1,
\end{equation}
che vale 1 sui passi osservati e 0 su quelli aggiunti. In un batch di tre sequenze lunghe 25, 33 e 40 passi, per esempio, tutte vengono portate a 40, e la maschera della prima vale 1 sui primi 25 passi e 0 sugli ultimi 15.

Un passo di padding e un bin valido senza eventi possono contenere entrambi soltanto zeri, ma significano cose opposte. Il bin vuoto dice che è trascorso del tempo senza attività, ed è a tutti gli effetti un'osservazione. Il padding non osserva nulla. Confonderli altera il campione, perché se la rete calcolasse la media delle proprie feature su tutti i 40 passi anziché sui 25 validi, il risultato della prima sequenza verrebbe ridotto di oltre un terzo senza alcuna ragione legata ai dati. Lo stesso accade se il padding entra nelle statistiche della normalizzazione, e un modulo dotato di stato interno deve far avanzare il proprio tempo soltanto sui passi validi.

C'è infine una terza situazione, diversa da entrambe, in cui la registrazione finisce prima della finestra prevista, come nell'esempio della finestra di $2\,\mathrm{s}$, e i passi successivi sono validi ma privi di eventi. La rete, però, non resta necessariamente inattiva, perché i potenziali di membrana e le memorie interne continuano a evolvere anche senza ingresso. Questa coda vuota e il padding di un batch vanno quindi dichiarati separatamente.

L'intera catena di elaborazione si può riassumere nello schema
\begin{equation}
\label{eq:decomposizione_flusso}
\mathcal E
\overset{\mathcal R}{\longrightarrow}X_{1:T}
\overset{\mathcal M}{\longrightarrow}Z_{1:T}
\overset{\mathcal D}{\longrightarrow}\hat y,
\end{equation}
dove $\mathcal R$ costruisce la rappresentazione, $\mathcal M$ ne elabora la dinamica e $\mathcal D$ prende la decisione. Ognuno dei tre stadi può limitare l'uso dell'informazione temporale. Per questo, quando si confrontano architetture diverse, rappresentazione e clock vanno tenuti fissi, e quando è l'ingresso stesso l'oggetto dell'esperimento ogni sua modifica va dichiarata come tale.


% ==============================================================================
%  Dinamica neuronale e memoria nelle Spiking Neural Networks
% ==============================================================================

\section{Dinamica neuronale e memoria nelle Spiking Neural Networks}
\label{sec:dinamica_neuronale_memoria}

Un neurone artificiale convenzionale, come quello di un percettrone multistrato con attivazione ReLU, calcola la propria uscita a partire dal solo ingresso corrente. Se riceve due volte lo stesso ingresso risponde due volte allo stesso modo, e non conserva nulla da un'elaborazione all'altra. Il neurone di una Spiking Neural Network (SNN) si comporta diversamente. Possiede una variabile interna, il potenziale di membrana, che accumula gli ingressi ricevuti e sopravvive da un passo al successivo, e la sua risposta dipende quindi anche da ciò che è accaduto prima. È la struttura di una rete ricorrente, con una differenza importante. In una rete ricorrente classica il passato rientra attraverso connessioni esplicite che riportano l'uscita di uno strato al suo ingresso, mentre in una SNN la ricorrenza è già dentro ogni neurone, e la rete ha memoria anche quando i suoi strati sono collegati soltanto in avanti~\cite{neftci2019surrogate}.

Questa sezione descrive come un neurone accumuli e dimentichi gli ingressi, come la velocità con cui dimentica dipenda dal clock dell'input, come si addestri una rete che comunica con impulsi discontinui e che cosa una memoria di questo tipo permetta di fare.

\subsection{Integrazione, soglia e reset}
\label{subsec:integrazione_soglia_reset}

Il modello più semplice è l'Integrate-and-Fire (IF), che si può immaginare come un accumulatore. A ogni passo il neurone somma al proprio potenziale il contributo dell'ingresso,
\begin{equation}
\label{eq:if_discreto}
H_k=V_{k-1}+X_k,
\end{equation}
e quando il potenziale raggiunge una soglia $V_{\mathrm{th}}$ emette uno spike,
\begin{equation}
\label{eq:if_spike}
S_k=\Theta(H_k-V_{\mathrm{th}}),
\qquad
\Theta(z)=
\begin{cases}
1,&z\geq0,\\
0,&z<0.
\end{cases}
\end{equation}
La funzione a gradino $\Theta$ traduce il confronto con la soglia in una risposta tutto o niente. Lo spike vale 1 se il potenziale ha raggiunto la soglia e 0 altrimenti, senza valori intermedi, e un potenziale appena sopra la soglia produce lo stesso spike di uno che la supera di molto.

Le due equazioni usano due simboli diversi per il potenziale, e conviene chiarire subito perché. All'interno di ogni passo avvengono due cose in sequenza. Prima il neurone riceve l'ingresso e si carica, e $H_k$ indica il potenziale in questo momento, prima che il neurone decida se emettere lo spike. Poi, se lo spike è stato emesso, il potenziale viene riportato verso il basso, ed è il valore che resta dopo questa operazione, indicato con $V_k$, a passare al passo successivo. Se lo spike non viene emesso, i due valori coincidono.

L'operazione che segue lo spike si chiama reset, e ne esistono due varianti. L'\emph{hard reset} riporta il potenziale a un valore fisso, di solito zero, mentre il \emph{soft reset}, adottato in questo lavoro, si limita a sottrarre la soglia,
\begin{equation}
\label{eq:soft_reset}
V_k=H_k-S_k V_{\mathrm{th}}.
\end{equation}
La differenza emerge quando il potenziale supera la soglia di un certo margine. Con soglia 1 e $H_k=1{,}3$, l'hard reset riparte da zero e perde l'eccedenza di $0{,}3$, mentre il soft reset riparte da $0{,}3$ e la conserva per il passo successivo~\cite{han2020rmp}. Qui e nel seguito si adotta la convenzione con potenziale di riposo nullo.

Il neurone IF ha una memoria perfetta, e al di fuori del reset un contributo ricevuto all'inizio della sequenza pesa esattamente quanto uno appena arrivato, perché il modello non prevede alcuna forma di oblio.\\
Il Leaky Integrate-and-Fire (LIF) introduce proprio questo elemento, una perdita o \emph{leak}. A ogni passo, prima di sommare il nuovo ingresso, il potenziale perde una frazione fissa del proprio valore,
\begin{equation}
\label{eq:lif_decadimento_discreto}
H_k=\lambda V_{k-1}+\eta X_k,
\qquad 0<\lambda<1.
\end{equation}
Il fattore di decadimento $\lambda$ indica quale frazione del potenziale sopravvive da un passo al successivo, e il coefficiente $\eta$ stabilisce quanto pesi l'ingresso. Con $\lambda=0{,}5$, in assenza di ingresso, il potenziale si dimezza a ogni passo. Il neurone IF si ritrova ponendo $\lambda=1$ ed $\eta=1$.

La perdita cambia in modo sostanziale il comportamento del neurone, come mostra la figura~\ref{fig:dinamica_lif}. Un neurone con soglia 1 riceve tre ingressi da $0{,}6$, una volta in tre passi consecutivi e una volta a tre passi di distanza l'uno dall'altro. Il neurone IF emette uno spike in entrambi i casi, non appena il secondo ingresso porta il potenziale a $1{,}2$. Il neurone LIF con $\lambda=0{,}5$ ed $\eta=1$, quando gli ingressi sono ravvicinati, sale a $0{,}6$, poi a $0{,}9$ e infine a $1{,}05$, e supera la soglia al terzo passo. Quando invece sono distanziati, tra un ingresso e l'altro il potenziale ha il tempo di scendere, non va mai oltre $0{,}68$ e lo spike non arriva. A parità di ingresso complessivo, il LIF risponde agli ingressi che arrivano vicini e lascia cadere quelli isolati, e si comporta in questo senso come un rilevatore di coincidenze.

\begin{figure}[t]
    \centering
    \includegraphics[width=\textwidth]{images/cap1/cap1_dinamica_lif.pdf}
    \caption[Integrazione, soglia e soft reset nei neuroni IF e LIF]{Integrazione, soglia e soft reset in un neurone IF (tratteggio arancione) e in un neurone LIF con $\lambda=0{,}5$ (linea blu), entrambi con soglia $V_{\mathrm{th}}=1$. I due neuroni ricevono tre ingressi da $0{,}6$ (in basso), consecutivi in (a) e distanziati in (b). I marcatori indicano il potenziale $H_k$ appena caricato, i tratti verticali il soft reset che segue uno spike, la fila in alto gli spike emessi. Il neurone IF, che non dimentica, emette uno spike in entrambi i casi. Il LIF lo emette soltanto quando gli ingressi sono abbastanza ravvicinati da sommarsi prima che il potenziale decada.}
    \label{fig:dinamica_lif}
\end{figure}

\subsection{Decadimento discreto e scala temporale fisica}
\label{subsec:decadimento_scala_fisica}

Il fattore $\lambda$ è definito per passo, mentre la memoria di un neurone, dal punto di vista della scena osservata, si misura in millisecondi. Per collegare le due cose bisogna risalire al modello da cui il LIF discende. Nella formulazione originale, a tempo continuo, la membrana si comporta come un condensatore che si scarica lentamente attraverso una resistenza,
\begin{equation}
\label{eq:lif_continuo}
\tau_m\frac{\mathrm d V(t)}{\mathrm dt}
=-\bigl(V(t)-V_{\mathrm{rest}}\bigr)+R I(t),
\end{equation}
dove $I(t)$ è la corrente d'ingresso, $R$ la resistenza di membrana, $V_{\mathrm{rest}}$ il potenziale di riposo e $\tau_m$ la costante di tempo di membrana~\cite{gerstner2014neuronal}. Il primo termine a destra riporta il potenziale verso il riposo con una spinta proporzionale alla distanza, il secondo lo sposta in base all'ingresso. Se l'ingresso si annulla, il potenziale torna verso il riposo lungo un'esponenziale. Trascorso un tempo pari a $\tau_m$ la distanza dal riposo si è ridotta a circa il $37\%$, poiché $e^{-1}\approx0{,}37$, e la costante $\tau_m$ misura dunque, in unità di tempo, per quanto il neurone ricorda un ingresso.

Per simulare questa dinamica a passi di durata $\Delta t$ occorre discretizzarla. Molte implementazioni, tra cui la libreria SpikingJelly~\cite{fang2023spikingjelly}, adottano la discretizzazione di Eulero, che sostituisce la derivata con un rapporto incrementale e ritrova proprio la forma dell'equazione~\eqref{eq:lif_decadimento_discreto},
\begin{equation}
\label{eq:lif_euler_discreto}
H_k=V_{k-1}+\frac{X_k-V_{k-1}}{\tau_d},
\qquad
\lambda=1-\frac1{\tau_d},\quad
\eta=\frac1{\tau_d}.
\end{equation}
Qui la costante $\tau_d$, che deve superare 1 perché $\lambda$ resti positivo, non è un tempo ma un numero puro, espresso in passi. Letta come discretizzazione dell'equazione continua vale $\tau_d=\tau_m/\Delta t$, ma Eulero approssima soltanto la soluzione esatta, per la quale varrebbe $\lambda=e^{-\Delta t/\tau_m}$. Con una costante pari a due passi si ottiene $\lambda=0{,}5$ nel primo caso e circa $0{,}61$ nel secondo. Quando un lavoro riporta una costante di tempo, quindi, bisogna sapere se si tratti di un tempo fisico o di un numero di passi, e con quale discretizzazione venga usata.

Srotolando la ricorrenza, cioè sostituendo a ritroso ogni potenziale con la sua espressione, si vede che cosa il neurone conservi della propria storia. In assenza di spike e reset, con coefficienti costanti, si ottiene
\begin{equation}
\label{eq:lif_espansione_temporale}
V_k=\lambda^k V_0+
\eta\sum_{j=1}^{k}\lambda^{k-j}X_j.
\end{equation}
Il potenziale è la somma di tutti gli ingressi ricevuti, ciascuno moltiplicato per $\lambda$ tante volte quanti sono i passi trascorsi dal suo arrivo. Con $\lambda=0{,}5$, rispetto all'ingresso appena arrivato, quello del passo precedente pesa la metà, quello di tre passi prima un ottavo e quello di dieci passi prima meno di un millesimo. Gli ingressi passati, però, non vi compaiono come osservazioni separate. Sono fusi in un unico numero, e su questa proprietà si tornerà alla fine della sezione.

La stessa espressione dice come tradurre $\lambda$ in un tempo fisico. Un ingresso arrivato $n$ passi prima, cioè $n\Delta t$ secondi prima, pesa $\lambda^n$. Cercando la costante $\tau_{\mathrm{eq}}$ per cui $\lambda^n=\exp(-n\Delta t/\tau_{\mathrm{eq}})$, come nel rilassamento continuo, si trova
\begin{equation}
\label{eq:tau_fisica_equivalente}
\tau_{\mathrm{eq}}=-\frac{\Delta t}{\ln\lambda}.
\end{equation}
Con $\lambda=0{,}5$ vale circa $1{,}44\,\Delta t$, quindi poco più di $14\,\mathrm{ms}$ con bin da $10\,\mathrm{ms}$ e poco più di $140\,\mathrm{ms}$ con bin da $100\,\mathrm{ms}$.

Il punto essenziale è che $\tau_{\mathrm{eq}}$ dipende dalla durata del bin, cioè dal clock dell'input. Lo stesso neurone, con lo stesso $\lambda$, dimentica due volte più in fretta in tempo reale se la durata dei bin si dimezza. Adeguare $\lambda$ al clock conserverebbe la memoria fisica, ma non renderebbe equivalenti i due modelli, perché cambiando i bin cambiano anche i conteggi che ciascuno contiene e il numero di occasioni in cui il neurone può emettere spike. Con bin di durata variabile, come nelle partizioni a eventi costanti, un coefficiente fisso per passo non corrisponde infine ad alcun decadimento uniforme nel tempo fisico, perché un passo lungo mezzo secondo e uno lungo cinque millisecondi fanno perdere al potenziale la stessa frazione.

Queste relazioni descrivono un singolo neurone isolato, al di sotto della soglia. La memoria di una rete, che combina molti neuroni, trasformazioni apprese e interazioni tra strati, non ha un orizzonte massimo così rigidamente definito.

\subsection{Addestramento con surrogate gradient}
\label{subsec:surrogate_gradient}

Le reti neurali si addestrano con la retropropagazione dell'errore, che calcola per ogni peso di quanto cambierebbe la loss modificandolo di poco, e lo corregge di conseguenza. La soglia che genera gli spike rende questo calcolo inutilizzabile. La funzione $\Theta$ è piatta ovunque tranne che nell'origine, dove salta da 0 a 1. Se il potenziale cambia di poco lo spike resta quello che era, e la derivata è quindi nulla quasi ovunque, mentre nell'unico punto in cui lo spike cambia la derivata non è definita. In entrambi i casi nessun segnale di errore riesce ad attraversare la soglia e a raggiungere i pesi degli strati precedenti.

I \emph{surrogate gradient} aggirano l'ostacolo con un compromesso~\cite{neftci2019surrogate}. Nel passaggio in avanti (forward pass), quello che produce le uscite della rete, lo spike resta esattamente quello dell'equazione~\eqref{eq:if_spike}. Nel passaggio all'indietro (backward pass), quello che calcola i gradienti, la sua derivata viene sostituita da una funzione regolare $g_\alpha(H_k-V_{\mathrm{th}})$. Una scelta diffusa è la derivata di una sigmoide. Con $\sigma(u)=(1+e^{-u})^{-1}$ si pone
\begin{equation}
\label{eq:surrogate_sigmoide}
g_\alpha(z)=\alpha\,\sigma(\alpha z)
                    \bigl[1-\sigma(\alpha z)\bigr],
\qquad \alpha>0,
\end{equation}
che è proprio la derivata di $\sigma(\alpha z)$. In pratica la rete calcola gli spike con il gradino ma, al momento di calcolare i gradienti, si comporta come se il gradino fosse stato una sigmoide (figura~\ref{fig:surrogate_gradient_schema}). Il parametro $\alpha$ regola la ripidità della sigmoide, e quindi l'altezza e la larghezza del gradiente surrogato, che raggiunge il massimo sulla soglia e si annulla rapidamente allontanandosene. Ricevono così un gradiente apprezzabile soltanto i neuroni il cui potenziale era vicino alla soglia, quelli cioè per cui una piccola modifica dei pesi avrebbe potuto cambiare la decisione di emettere o meno lo spike.

\begin{figure}[t]
    \centering
    \includegraphics[width=\textwidth]{images/cap1/cap1_surrogata.pdf}
    \caption[Principio del surrogate gradient]{Principio del surrogate gradient, in funzione della distanza $z=H_k-V_{\mathrm{th}}$ tra potenziale e soglia. (a)~Nel passaggio in avanti lo spike è prodotto dal gradino $\Theta$ (linea nera). La sigmoide $\sigma(4z)$ (tratteggio) non viene mai calcolata, ma è la funzione di cui il passaggio all'indietro usa la derivata. (b)~La derivata del gradino è nulla ovunque tranne nell'origine, dove non è definita (freccia), e non trasmette alcun gradiente. La surrogata $g_\alpha$ la sostituisce con una campana centrata sulla soglia, tanto più alta e stretta quanto maggiore è $\alpha$. Il valore evidenziato, $\alpha=4$, è quello adottato in questo lavoro.}
    \label{fig:surrogate_gradient_schema}
\end{figure}

La soglia non è l'unico ostacolo, perché il potenziale collega ogni passo al successivo e la loss dipende dai pesi anche attraverso il tempo. La Backpropagation Through Time (BPTT) affronta il problema srotolando la ricorrenza~\cite{wu2018spatiotemporal}. La rete viene vista come una catena di $T$ copie di sé stessa, una per ciascun passo, che condividono gli stessi pesi e sono collegate dagli stati che ciascuna passa alla successiva. Su questa catena, che non contiene più cicli, la retropropagazione ordinaria si applica senza difficoltà, e l'errore si propaga in due direzioni, attraverso gli strati e all'indietro nel tempo. Passando da un passo al precedente attraverso la membrana, il gradiente viene moltiplicato all'incirca per $\lambda$, lo stesso fattore con cui il neurone dimentica. Un neurone che dimentica in fretta impara quindi poco dalle conseguenze lontane dei propri stati, una forma della scomparsa del gradiente già nota nelle reti ricorrenti~\cite{pascanu2013difficulty}.

Il prezzo della BPTT è la memoria, dato che per percorrere la catena all'indietro servono i valori calcolati in avanti in ogni passo, e il solo stato neuronale da conservare cresce come $\mathcal O(T\sum_l N_l)$, dove $N_l$ è il numero di unità dello strato $l$~\cite{neftci2019surrogate}. In inferenza, invece, a una rete di soli neuroni LIF basta il potenziale attuale di ciascun neurone, cioè $\sum_l N_l$ valori qualunque sia $T$. Le due grandezze vanno tenute distinte, perché un modello può richiedere pochissimo stato durante l'esecuzione e, al tempo stesso, un grafo molto esteso durante l'ottimizzazione.

In questo lavoro la rete spiking viene addestrata direttamente, con surrogate gradient e BPTT. È un paradigma diverso dalla conversione ANN-SNN, in cui una rete convenzionale viene prima addestrata e poi tradotta in una rete spiking che ne approssima le attivazioni con la frequenza degli spike. La dinamica discreta, la funzione surrogata e il trattamento del reset effettivamente adottati sono specificati nella sezione~\ref{subsec:neurone_convenzioni}.

\subsection{Memoria neuronale e relazioni tra feature}
\label{subsec:memoria_relazioni}

Si torni all'equazione~\eqref{eq:lif_espansione_temporale}. Per un neurone LIF privo di ulteriori stati interni, tutto ciò che la storia degli ingressi può trasmettere al futuro passa per un unico numero, il potenziale $V_k$. L'intera sequenza viene compressa in un solo valore, e storie diverse possono lasciare lo stesso potenziale. Con $\lambda=0{,}5$ ed $\eta=1$, partendo da un potenziale nullo, un ingresso pari a $0{,}4$ arrivato al passo precedente e un ingresso pari a $0{,}2$ arrivato adesso lasciano entrambi $V_k=0{,}2$. Da quel momento, se gli ingressi successivi coincidono, il neurone si comporta in modo identico nei due casi, e non ha più modo di sapere se abbia visto qualcosa di intenso poco fa o qualcosa di debole in questo istante. La rete nel suo complesso può comunque conservare la differenza in altre unità o in altri moduli, sicché la compressione operata da un singolo potenziale non implica di per sé una perdita equivalente a livello dell'intera rete.

Si delineano così due capacità diverse, la \emph{persistenza}, cioè la capacità di conservare a lungo una traccia del passato, e l'\emph{accesso}, cioè la possibilità di recuperare separatamente ciò che è accaduto in un istante preciso per confrontarlo con il presente. Si potrebbe pensare di migliorare la memoria di un neurone avvicinando $\lambda$ a 1, ma la modifica agisce soltanto sulla prima. I contributi passati sopravvivono più a lungo, e proprio per questo un numero maggiore di ingressi si mescola nella stessa somma, che diventa ancora più difficile da scomporre. Si pensi alla barra della sezione~\ref{subsec:contrasto_eventi}, la cui direzione è scritta nell'ordine con cui si attivano pixel vicini. Per riconoscerla non basta sapere che entrambi i pixel sono stati attivi di recente, informazione che un potenziale conserva bene. Bisogna sapere quale dei due si è attivato prima, cioè mettere a confronto due istanti distinti, e un'unica somma pesata può farlo soltanto in modo indiretto.

La distinzione tra persistenza e accesso motiva l'interesse per operatori capaci di combinare o confrontare esplicitamente feature diverse. La sezione successiva descrive come l'attenzione sia stata integrata nelle architetture spiking, e sarà allora necessario distinguere le relazioni costruite all'interno di uno stesso passo da quelle che collegano istanti diversi.


% ==============================================================================
%  Dalle SNN agli Spiking Transformer
% ==============================================================================


\section{Dalle SNN agli Spiking Transformer}
\label{sec:snn_spiking_transformer}

La dinamica neuronale offre alle SNN un modo naturale di trasportare informazione da un passo temporale, o \emph{timestep}, al successivo. Resta però aperto un problema di altra natura, quello di mettere in relazione le numerose feature che compongono una rappresentazione visiva. Un neurone LIF integra nel tempo il proprio ingresso, ma nulla nella sua dinamica gli permette di stabilire quali regioni della scena siano rilevanti l'una per l'altra, né di collegare feature spazialmente lontane. Questa sezione ripercorre il modo in cui l'attenzione dei Transformer, nata per affrontare questo problema nelle reti convenzionali, è stata adattata al dominio spiking, fino all'architettura da cui muove la tesi.

\subsection{Self-attention e relazioni tra token}
\label{subsec:self_attention}

Le architetture convoluzionali affrontano il problema componendo operatori locali. Ogni convoluzione combina l'informazione di un intorno limitato, e soltanto impilando molti strati, che allargano progressivamente il campo recettivo, la rete arriva a costruire relazioni spaziali estese. Su questo principio si fonda l'efficacia delle CNN nella visione artificiale, ed è su di esso che sono state costruite molte delle prime SNN profonde.

Il Transformer adotta una prospettiva diversa e, anziché affidare la costruzione delle dipendenze a una lunga catena di operatori locali, permette agli elementi di una rappresentazione di interagire direttamente attraverso il meccanismo di self-attention~\cite{vaswani2017attention}. Trasferita alla visione, l'idea consiste nel suddividere l'immagine in porzioni, le patch, ciascuna delle quali viene proiettata in un vettore di feature detto token~\cite{dosovitskiy2021image}. I token possono poi influenzarsi a vicenda in base al contenuto che rappresentano, qualunque sia la distanza che li separa nell'immagine.

Il meccanismo si capisce bene a partire da un dizionario, la struttura dati che associa a ogni chiave un valore. Per recuperare un valore si presenta una richiesta (\emph{query}), la si confronta con le chiavi (\emph{key}) e si ottiene il valore (\emph{value}) associato alla chiave identica alla richiesta. Un dizionario di traduzioni, interrogato con la parola «gatto», trova la chiave «gatto» e restituisce «cat». L'attenzione rende graduale questa ricerca (figura~\ref{fig:attenzione_dizionario}). Richiesta e chiavi diventano vettori di feature, e il confronto non dà più un esito sì o no ma un punteggio di somiglianza con ciascuna chiave. Il risultato, di conseguenza, non è un singolo valore ma una media di tutti i valori, pesata con quei punteggi. Con pesi pari a $0{,}7$, $0{,}2$ e $0{,}1$ si ottiene $0{,}7\,v_1+0{,}2\,v_2+0{,}1\,v_3$, dove conta soprattutto il valore associato alla chiave più affine alla richiesta.

\begin{figure}[t]
    \centering
    \includegraphics[width=\textwidth]{images/cap1/cap1_attenzione_dizionario.pdf}
    \caption[L'attenzione come ricerca graduale in un dizionario]{L'attenzione come ricerca graduale in un dizionario. (a)~Un dizionario confronta la richiesta con le chiavi e restituisce soltanto il valore associato alla chiave identica. (b)~L'attenzione confronta la Query con tutte le Key, ricava da ciascun confronto un peso e restituisce la media dei Value pesata in questo modo. Lo spessore dei collegamenti è proporzionale al peso, e i valori numerici sono illustrativi.}
    \label{fig:attenzione_dizionario}
\end{figure}

Nella self-attention richiesta, chiavi e valori provengono tutti dagli stessi token, ed è da qui che viene il prefisso \emph{self}. Data una matrice di token $X\in\mathbb{R}^{N\times D}$, con $N$ token di dimensione $D$, tre trasformazioni lineari apprese ricavano da ciascun token tre vettori,
\begin{equation}
\label{eq:qkv_projections}
Q = X W_Q, \qquad K = X W_K, \qquad V = X W_V,
\end{equation}
la Query, la Key e il Value. Ogni token ricopre quindi tre ruoli, e con la Query esprime che cosa cerca negli altri token, con la Key descrive che cosa offre a chi lo interroga, con il Value fornisce il contenuto che gli altri riceveranno se lo trovano rilevante. Tenere separati i tre vettori permette alla rete di imparare che ciò che un token cerca può essere diverso da ciò che offre. Il token che copre un tratto di contorno, per esempio, può imparare a cercare i token che ne proseguono il profilo, anche se si trovano dall'altra parte dell'immagine.

Le formule traducono direttamente questa ricerca, e l'elemento $(i,j)$ del prodotto $QK^{\top}$ è il prodotto scalare tra la Query del token $i$ e la Key del token $j$, cioè il punteggio di somiglianza tra ciò che $i$ cerca e ciò che $j$ offre. La softmax trasforma ogni riga di punteggi in pesi positivi che sommano a 1, e questi pesi combinano i Value,
\begin{equation}
\label{eq:vanilla_softmax_attention}
Y = \operatorname{softmax}\left( \frac{Q K^{\top}}{\sqrt{d}} \right) V,
\end{equation}
dove la divisione per $\sqrt{d}$, con $d$ dimensione di Query e Key, evita che i prodotti scalari crescano troppo con la dimensione. La riga $i$ di $Y$ è il nuovo token $i$, cioè la media dei Value di tutti i token pesata secondo la loro compatibilità con la sua Query. Nel seguito questa forma viene indicata come VSA (\emph{Vanilla Self-Attention}). Ne deriva una capacità che la convoluzione non possiede nella stessa forma. Due regioni lontane possono interagire direttamente all'interno di un singolo blocco, senza che l'informazione debba attraversare una lunga successione di operatori locali.

Questa libertà ha un prezzo, dato che $QK^{\top}$ confronta ogni token con tutti gli altri, e la matrice di attenzione ha quindi dimensione $N\times N$, il costo del prodotto cresce come $\mathcal{O}(N^2D)$ e la memoria necessaria a conservare la matrice come $\mathcal{O}(N^2)$. Con 256 token la matrice contiene già più di $65\,000$ elementi. In una SNN se ne aggiunge un secondo vincolo, perché le stesse operazioni vanno ripetute a ogni timestep. Per un ingresso $X\in\mathbb{R}^{T\times N\times D}$ l'attenzione viene tipicamente applicata, istante per istante, alle feature di ciascun passo, e la dimensione temporale, che finora era apparsa come la sorgente della memoria neuronale, diventa così anche un moltiplicatore del costo.

\subsection{La self-attention in forma spiking}
\label{subsec:ssa}

L'incontro tra SNN e Transformer nasce da questa tensione. La VSA si basa su attivazioni continue, prodotti tra matrici dense e normalizzazioni, operazioni che non traggono alcun vantaggio dalla natura binaria e sparsa degli spike. Sostituire le attivazioni di un Transformer con neuroni LIF non basta dunque a ottenere uno Spiking Transformer. Occorre ripensare almeno in parte il meccanismo di attenzione, altrimenti la rappresentazione spiking verrebbe subito riconvertita in una computazione del tutto analoga a quella di una rete convenzionale (ANN).

Per capire che cosa si voglia preservare conviene anticipare un'idea che la sezione~\ref{sec:efficienza_neuromorfica_vincoli} riprenderà. Per applicare un peso a un'attivazione continua serve una moltiplicazione seguita da un accumulo (MAC). Se l'attivazione è uno spike binario la moltiplicazione scompare, perché basta sommare il peso quando lo spike c'è e non fare nulla quando manca, con un semplice accumulo (AC). Una rete si dice \emph{spike-driven} quando funziona così, cioè quando ogni operazione è innescata da uno spike e l'assenza di uno spike coincide con l'assenza di calcolo~\cite{yao2023spikedriven}.

Con Spikformer, Zhou et al. propongono nel 2023 la Spiking Self-Attention (SSA), mostrando che Query, Key e Value possono essere espressi in forma spiking e che la softmax può essere eliminata~\cite{zhou2023spikformer}. Dalle feature spiking si ricavano tre rappresentazioni binarie,
\begin{equation}
\label{eq:spiking_qkv}
Q_s, \; K_s, \; V_s \in \{0, 1\}^{N \times D},
\end{equation}
la cui interazione continua a svolgere il ruolo dell'attenzione. In forma semplificata,
\begin{equation}
\label{eq:ssa_spikformer}
\operatorname{SSA}(Q_s, K_s, V_s) = \operatorname{SN}\left( Q_s K_s^{\top} V_s \, \beta \right),
\end{equation}
dove $\operatorname{SN}(\cdot)$ indica lo strato di neuroni spiking che riporta il risultato in forma binaria e $\beta$ è un fattore di scala. La rinuncia alla softmax ha un peso particolare, dato che la softmax richiede esponenziali, divisioni e una normalizzazione che coinvolge tutti i coefficienti di una riga, operazioni estranee a una computazione guidata dagli eventi. Nel dominio spiking, inoltre, viene meno il suo compito principale, quello di rendere positivi i pesi, perché il prodotto di due matrici binarie non può produrre valori negativi~\cite{zhou2023spikformer}.

Spikformer dimostra così che self-attention e SNN sono compatibili, ma porta alla luce due questioni. Anzitutto, che Query, Key e Value siano spike non garantisce che l'intero blocco esegua soltanto operazioni economiche. L'elemento $(i,j)$ di $Q_sK_s^{\top}$ conta in quanti canali la Query del token $i$ e la Key del token $j$ sono attive insieme. Se la prima vale $[1,0,1,1]$ e la seconda $[1,1,1,0]$ il risultato è 2, perché coincidono nel primo e nel terzo canale. Calcolare questo numero è economico, perché basta incrementare un contatore quando entrambi i bit valgono 1. Il risultato, però, non è più uno spike ma un intero, che può arrivare al numero dei canali, e lo strato SN lo riporta in forma binaria soltanto dopo che è stato combinato con i Value e scalato. Un problema analogo si ripresenta nelle connessioni residue, dove la somma di due spike può già valere 2. Quando un valore intero raggiunge uno strato dotato di pesi, ciascun peso va moltiplicato per quel valore, e ricompaiono le moltiplicazioni che il paradigma spike-driven vorrebbe evitare.

La seconda questione riguarda la scalabilità, perché il prodotto $Q_sK_s^{\top}$ conserva una matrice tra tutte le coppie di token, e anche con dati sparsi la dimensione logica dell'operazione resta legata a $N^2$, un costo che diventa critico al crescere del numero dei token. Da questi due problemi sono nate linee di ricerca diverse e complementari, che le prossime sottosezioni esaminano nell'ordine.

\subsection{Attenzione e residui spike-driven}
\label{subsec:attenzione_residui}

Spike-driven Transformer affronta la prima questione, chiedendosi quanta parte della rete sia davvero spike-driven~\cite{yao2023spikedriven}. Gli autori osservano che diverse architetture precedenti mescolano accumuli e moltiplicazioni, dando luogo a una computazione ibrida, e si pongono l'obiettivo di mantenere binaria la comunicazione tra i blocchi. Il loro meccanismo, la Spike-Driven Self-Attention (SDSA), abbandona il prodotto completo $Q_sK_s^{\top}$. In una delle sue formulazioni, Query e Key vengono moltiplicate elemento per elemento e sommate lungo la dimensione dei token, e il risultato attraversa uno strato spiking,
\begin{equation}
\label{eq:sdsa_vector}
A_s = \operatorname{SN}\left( \operatorname{SUM}_N \left( Q_s \odot K_s \right) \right),
\end{equation}
dopodiché il vettore binario così ottenuto viene applicato ai Value,
\begin{equation}
\label{eq:sdsa_value}
Y_s = A_s \odot V_s,
\end{equation}
dove $\odot$ indica il prodotto elemento per elemento. Con operandi binari questo prodotto equivale a una maschera, perché ogni elemento viene semplicemente conservato oppure azzerato. Nel procedimento il prodotto $Q_s\odot K_s$ segna dove Query e Key sono attive insieme, la somma conta per ciascun canale quante volte ciò accade tra tutti i token, e il neurone trasforma questi conteggi in una decisione binaria, canale per canale. I canali promossi attraversano il blocco, gli altri vengono azzerati.

\begin{figure}[t]
    \centering
    \includegraphics[width=\textwidth]{images/cap1/cap1_operatori_attenzione.pdf}
    \caption[Riformulazione della self-attention nel dominio spiking]{Riformulazione della self-attention nel dominio spiking, calcolata su un piccolo esempio casuale con 6 token (righe) e 5 canali (colonne). Il colore indica il tipo di valore, e la cornice spessa la struttura che svolge il ruolo di attenzione. Dalla VSA alla SSA scompare la softmax e i dati in ingresso diventano binari, ma la matrice tra coppie di token resta e contiene conteggi. SDSA e QKTA sostituiscono la matrice con un vettore binario, che in uscita azzera interi canali nel primo caso e interi token nel secondo. La soglia degli strati SN è scelta per rendere leggibile l'esempio.}
    \label{fig:evoluzione_self_attention_spiking}
\end{figure}

Rispetto all'equazione~\eqref{eq:vanilla_softmax_attention} il cambiamento è sostanziale, come mostra la figura~\ref{fig:evoluzione_self_attention_spiking}. Non viene costruita alcuna matrice $N\times N$ di relazioni tra coppie di token, e l'attenzione si riduce a una scelta su quali canali lasciar passare. Per questo SDSA può essere interpretata come una forma particolare di attenzione lineare, la cui complessità cresce linearmente con il numero di token e di canali anziché quadraticamente con il numero di token~\cite{yao2023spikedriven}. Si rinuncia così a parte della ricchezza delle correlazioni tra coppie, in cambio di un'operazione molto più adatta alla computazione sparsa.

Spike-driven Transformer interviene anche su un elemento meno appariscente, ma importante per comprendere gli sviluppi successivi, cioè la posizione delle connessioni residue. Una connessione residua somma l'ingresso di un blocco alla trasformazione che il blocco stesso produce, e nelle reti profonde facilita la propagazione dell'informazione e dei gradienti. In una SNN, però, conta dove la somma viene eseguita, se tra spike oppure tra potenziali di membrana (figura~\ref{fig:residui_sew_ms}).

Nella connessione residua Spike-Element-Wise (SEW) adottata da Spikformer si sommano direttamente due tensori di spike. Il risultato può valere $0$, $1$ o $2$, e l'intervallo si allarga a ogni blocco attraversato, perché la somma di un blocco diventa un addendo di quella del successivo. Il segnale smette così di essere binario, e quando raggiunge uno strato pesato non basta più selezionare i pesi corrispondenti agli spike attivi, perché ciascun peso va moltiplicato per un valore intero, e ricompaiono moltiplicazioni tra interi e numeri in virgola mobile.

\begin{figure}[t]
    \centering
    \includegraphics[width=\textwidth]{images/cap1/cap1_residui.pdf}
    \caption[Connessioni residue SEW e MS]{Due posizioni della connessione residua, indicate con i valori che ogni tratto può assumere~\cite{zhou2026spikingformer}. (a)~Nel residuo SEW la somma avviene tra spike. I valori crescono a ogni blocco, e lo strato pesato successivo deve moltiplicare ciascun peso per un intero. (b)~Nel residuo MS la somma avviene tra potenziali di membrana, cioè tra valori reali, e un neurone spiking precede ogni strato pesato, che riceve così soltanto ingressi binari.}
    \label{fig:residui_sew_ms}
\end{figure}

Spike-driven Transformer sposta invece la somma residua sui potenziali di membrana, e colloca un neurone spiking subito prima di ogni strato pesato. La somma avviene così tra valori reali, come quella che ogni neurone esegue già al proprio interno, e ciò che entra negli strati pesati resta binario~\cite{yao2023spikedriven}. Il dettaglio contiene una lezione progettuale di portata più generale. Contare i neuroni spiking presenti in un'architettura dice poco sulla sua natura event-driven, che dipende piuttosto dalla posizione di quei neuroni rispetto alle operazioni lineari e ai percorsi residui.

Spikingformer approfondisce lo stesso punto in una linea di ricerca parallela~\cite{zhou2026spikingformer}. Gli autori analizzano le connessioni residue di Spikformer e mostrano che la variante SEW genera valori a più bit lungo i collegamenti diretti, mentre una Membrane Shortcut (MS) preserva ingressi binari per gli operatori successivi. Combinando questa soluzione con la SSA, il modello mira a conservare insieme comunicazione spike-driven e capacità di modellazione globale.

Il confronto mette in luce un compromesso che accompagna l'intera evoluzione degli Spiking Transformer. Semplificare l'attenzione in modo aggressivo ne riduce il costo, ma impoverisce le relazioni che il modello può rappresentare. Conservare una self-attention più globale offre maggiore espressività, ma obbliga a intervenire altrove per preservare l'efficienza spike-driven. La storia di queste architetture non va letta, allora, come una successione di modelli via via migliori. È piuttosto l'esplorazione di uno spazio definito da almeno tre assi, la capacità relazionale, la scalabilità e l'aderenza al paradigma spike-driven, lungo i quali ciascun lavoro sceglie la propria posizione.

\subsection{La struttura gerarchica}
\label{subsec:struttura_gerarchica}

Resta aperta una questione di altro genere, che riguarda il modo in cui scalare l'architettura stessa. Un Vision Transformer tradizionale mantiene spesso la stessa risoluzione lungo gran parte dell'encoder, e molte delle prime architetture Spiking Transformer seguono la stessa strategia. Spikformer, per esempio, usa uno Spiking Patch Splitting (SPS) per convertire l'ingresso in $N$ token di dimensione $D$, elaborati poi da una sequenza di blocchi che ne conservano il numero e la dimensione (figura~\ref{fig:scala_unica_gerarchia}a).

La soluzione è semplice, ma si discosta dalla struttura gerarchica che caratterizza molte CNN e numerosi Vision Transformer moderni. In una rete gerarchica le prime fasi operano su molte posizioni con pochi canali, mentre quelle profonde lavorano su pochi token descritti da feature più ricche. Questa progressione preserva il dettaglio spaziale nelle prime fasi mentre costruisce rappresentazioni via via più astratte, e riduce il numero dei token proprio dove la dimensione delle feature, e con essa il costo dei blocchi Transformer, tende a crescere.

\begin{figure}[t]
    \centering
    \includegraphics[width=\textwidth]{images/cap1/cap1_gerarchia.pdf}
    \caption[Struttura a scala unica e struttura gerarchica]{Struttura a scala unica e struttura gerarchica, secondo le descrizioni originali. I rettangoli rappresentano le feature, con l'altezza proporzionale al lato della griglia dei token e la larghezza al numero dei canali, mentre il colore distingue i blocchi con SSA da quelli con Q-K Token Attention (QKTA). (a)~In Spikformer tutti i blocchi operano sugli stessi $N$ token. (b)~In QKFormer ogni modulo SPEDS riduce la griglia e aumenta i canali. La dimensione temporale $T$ resta invariata lungo entrambe le reti.}
    \label{fig:scala_unica_gerarchia}
\end{figure}

Da questa esigenza nasce QKFormer~\cite{zhou2024qkformer}. Gli autori individuano due limiti congiunti delle architetture precedenti, la complessità quadratica della SSA rispetto al numero dei token e l'elevato costo di memoria che deriva dal combinare la dimensione spaziale con la simulazione temporale. Per superarli introducono una struttura gerarchica a tre stadi, collegati da moduli di Spiking Patch Embedding with Deformed Shortcut (SPEDS), e un nuovo meccanismo di attenzione, la Q-K Attention (figura~\ref{fig:scala_unica_gerarchia}b). Oltre a chiedersi come rendere più economico il prodotto tra Q, K e V, QKFormer si domanda se tutte e tre le componenti siano davvero necessarie per selezionare l'informazione. La Q-K Attention ne usa soltanto due, Query e Key, e riduce l'attenzione a un vettore binario che decide quali token, o quali canali, lasciar passare, come mostra l'ultima riga della figura~\ref{fig:evoluzione_self_attention_spiking}. La sezione~\ref{sec:qkformer_riferimento} ne descrive il funzionamento in dettaglio.

La tabella~\ref{tab:evoluzione_attenzione_spiking} riassume il percorso. QKFormer ne è il punto di partenza naturale per questa tesi, perché riunisce tre proprietà preziose nel regime compatto studiato più avanti, un'attenzione di costo lineare, una rappresentazione gerarchica e una trasmissione residua che attraversa i cambi di risoluzione.

\begin{table}[t]
\centering
\small
\renewcommand{\arraystretch}{1.2}
\caption[Evoluzione dei meccanismi di attenzione spiking]{Evoluzione dei meccanismi di attenzione spiking rilevanti per il lavoro.}
\label{tab:evoluzione_attenzione_spiking}
\begin{adjustbox}{max width=\textwidth}
\begin{tabular}{@{}>{\raggedright\arraybackslash}p{2.6cm} >{\raggedright\arraybackslash}p{3.3cm} >{\raggedright\arraybackslash}p{2.6cm} >{\raggedright\arraybackslash}p{2.1cm} >{\raggedright\arraybackslash}p{2.1cm} c@{}}
\hline
\textbf{Architettura} & \textbf{Problema affrontato} & \textbf{Operatore} & \textbf{Ambito} & \textbf{Costo} & \textbf{Spike-driven} \\
\hline
Transformer / ViT & Relazioni globali tra token & VSA con softmax & Coppie di token & Quadratica in $N$ & No \\
Spikformer & Rendere spiking la self-attention & SSA con Q, K, V spiking & Coppie di token & Quadratica in $N$ & Parziale \\
Spike-driven Transformer & Eliminare le operazioni non spike-driven e ridurre il costo & SDSA, maschere e addizioni sparse & Selezione di canali & Lineare & Sì \\
Spikingformer & Conservare l'attenzione globale con residui spike-driven & SSA con residui MS & Coppie di token & Quadratica in $N$ & Sì \\
QKFormer & Scalabilità e struttura gerarchica & QKTA o QKCA con SPEDS & Selezione di token o canali & Lineare nella dimensione & Nell'attenzione \\
\hline
\end{tabular}
\end{adjustbox}
\vspace{0.4em}
\begin{flushleft}
\footnotesize
\emph{Nota.} «Spike-driven» indica qui la proprietà architetturale dichiarata dai rispettivi lavori. Il costo effettivo dipende comunque da firing rate, operatori periferici, rappresentazione degli input e piattaforma hardware, e verrà discusso nella sezione~\ref{sec:efficienza_neuromorfica_vincoli}.
\end{flushleft}
\end{table}

Altrettanto importante è osservare quale problema QKFormer non si proponga di risolvere. Le operazioni di attenzione descritte finora costruiscono relazioni tra token o canali all'interno delle feature di un singolo timestep. Data una sequenza $X_1,X_2,\dots,X_T$, l'attenzione applicata all'istante $t$ riceve tipicamente la sola $X_t$, e il collegamento con $X_{t-1}$ resta affidato principalmente agli stati dei neuroni che attraversano il tempo. L'espressione \emph{self-attention} suggerisce, in modo intuitivo ma impreciso, un accesso globale all'intera sequenza. La globalità va invece sempre riferita alla dimensione lungo cui l'operatore agisce. Uno Spiking Transformer può disporre di un'attenzione globale sui token e, insieme, di una comunicazione temporale prevalentemente ricorrente e locale.

Resta da capire se la dinamica neuronale descritta nella sezione~\ref{sec:dinamica_neuronale_memoria} basti a governare anche l'interazione tra $X_t$ e $X_{t-k}$, ed è da questa domanda che muove la sezione~\ref{sec:strategie_modellazione_temporale}. Prima conviene però esaminare in dettaglio il punto di arrivo dell'evoluzione descritta, perché QKFormer determina direttamente la baseline su cui si fonda la parte sperimentale della tesi.


% ==============================================================================
%  QKFormer come riferimento architetturale
% ==============================================================================


\section{QKFormer come riferimento architetturale}
\label{sec:qkformer_riferimento}

La sezione~\ref{subsec:struttura_gerarchica} ha presentato QKFormer come risposta ai limiti di scalabilità delle architetture precedenti~\cite{zhou2024qkformer}. Qui se ne riprende ciò che serve al resto del lavoro, perché da questa architettura discende la baseline compatta del capitolo~\ref{chap:baseline_metodologia}. Contano in particolare tre aspetti, la forma esatta del blocco di attenzione, il peso del patch embedding nella gerarchia e il modo in cui il modello tratta il tempo.

\paragraph{Il blocco con Q-K Token Attention}
\label{subsec:qk_attention}

La Q-K Attention ricava dal proprio ingresso $X$ due sole proiezioni spiking,
\begin{equation}
\label{eq:qk_projections_qkformer}
Q = \operatorname{SN}_{Q} \left( \operatorname{BN}(XW_Q) \right), \qquad K = \operatorname{SN}_{K} \left( \operatorname{BN}(XW_K) \right),
\end{equation}
dove BN indica la batch normalization e $Q,K\in\{0,1\}^{T\times N\times D}$, con $T$ timestep, $N$ token e $D$ canali. Nella Q-K Token Attention (QKTA), la variante adottata dalla baseline, le due proiezioni svolgono compiti distinti. La Query decide quali token lasciar passare, la Key ne trasporta il contenuto. Considerando un solo timestep, si sommano per ciascun token gli spike della Query lungo i canali, e il conteggio attraversa un neurone che lo traduce in una decisione binaria,
\begin{equation}
\label{eq:qkta_vettore_attenzione}
A^{\mathrm{tok}} = \operatorname{SN} \left( \sum_{d=1}^{D} Q_{:,d} \right), \qquad A^{\mathrm{tok}} \in \{0, 1\}^{N \times 1}, \qquad X' = A^{\mathrm{tok}} \odot K.
\end{equation}
Un token la cui Query è abbastanza attiva conserva in uscita la propria Key, gli altri vengono azzerati su tutti i canali, come mostra l'ultima riga della figura~\ref{fig:evoluzione_self_attention_spiking}. La variante Q-K Channel Attention applica lo stesso principio all'altro asse, sommando lungo i token per decidere quali canali conservare, ma non viene impiegata in questo lavoro.

Conviene chiarire anche che cosa questa operazione non faccia. Ogni componente di $A^{\mathrm{tok}}$ dipende soltanto dalla Query del token corrispondente, perché le proiezioni $XW_Q$ e $XW_K$ agiscono sui canali di ciascun token separatamente. Diversamente da VSA e SSA, la QKTA non confronta quindi token diversi, ma ne seleziona alcuni in base alla loro attività. Le relazioni tra posizioni diverse restano affidate alle convoluzioni dei patch embedding, su scala locale, e alla SSA dell'ultimo stadio, su scala globale. Il capitolo~\ref{chap:riallocazione_budget} riesamina questa scelta nel regime compatto.

Il guadagno è nel costo, perché l'attenzione si riduce a un vettore di $N$ elementi, e gli autori riportano una complessità $\mathcal{O}(D)$ in tempo e $\mathcal{O}(N)$ in spazio, contro $\mathcal{O}(N^2D)$ e $\mathcal{O}(N^2+ND)$ della SSA~\cite{zhou2024qkformer}. La somma degli spike ha inoltre una varianza molto più contenuta del prodotto $Q_sK_s^{\top}V_s$, e il fattore di scala $\beta$ può essere eliminato. Questi valori, però, riguardano la sola operazione di attenzione. Le matrici $Q$ e $K$ vanno comunque calcolate e occupano $\mathcal{O}(ND)$, e attorno alla maschera restano proiezioni, normalizzazioni e neuroni. Un'attenzione di costo trascurabile non implica quindi un blocco di costo trascurabile.

Al mascheramento segue infatti una proiezione lineare, che ricombina le feature dei token selezionati,
\begin{equation}
\label{eq:qk_proiezione_lineare_uscita}
X'' = \operatorname{SN} \left( \operatorname{BN} \left( \operatorname{Linear}(X') \right) \right),
\end{equation}
e il blocco completo conserva la struttura di un encoder Transformer,
\begin{align}
\label{eq:blocco_qkformer_qka}
X'_l &= \operatorname{QKA}(X_{l-1}) + X_{l-1}, \\
\label{eq:blocco_qkformer_smlp}
X_l &= \operatorname{SMLP}(X'_l) + X'_l,
\end{align}
dove QKA indica l'intero modulo di attenzione e SMLP uno spiking MLP~\cite{zhou2024qkformer}. Le somme residue avvengono a valle dei neuroni, secondo lo schema SEW della figura~\ref{fig:residui_sew_ms}a. L'attenzione opera quindi con maschere e addizioni, ma ciò che circola tra i blocchi contiene piccoli interi, e gli strati pesati che li ricevono eseguono moltiplicazioni. Per questo i conteggi di questo lavoro distinguono MAC e AC in base all'operando e non al tipo di strato (sezione~\ref{sec:caratterizzazione_computazionale}).

\paragraph{Gerarchia e patch embedding}
\label{subsec:gerarchia_patch}

La struttura a tre stadi è quella della figura~\ref{fig:scala_unica_gerarchia}b. Poiché a ogni passaggio i token diminuiscono e i canali aumentano, un unico tipo di attenzione non è ugualmente adatto a tutta la rete. Gli autori adottano perciò un'integrazione mista, con la QKTA nei primi stadi, dove i token sono molti, e la SSA nell'ultimo, dove sono pochi e il costo quadratico pesa meno. Le prove di ablazione sostengono la scelta, e su CIFAR10-DVS la combinazione di QKTA e SSA raggiunge l'$84{,}0\%$, mentre la sola QKTA si ferma all'$80{,}7\%$~\cite{zhou2024qkformer}.\footnote{Gli autori segnalano nel repository ufficiale un errore nell'ordine di reshape e LIF della proiezione SSA, con possibile sovrastima dei risultati ImageNet. I punteggi dell'articolo vanno quindi distinti da una riproduzione verificata dell'implementazione corretta~\cite{qkformerImplementationNotice}.} La baseline eredita questa sequenza, e il capitolo~\ref{chap:riallocazione_budget} verifica se la SSA resti necessaria anche con una capacità molto inferiore.

Un collegamento residuo ordinario presuppone che ingresso e uscita abbiano le stesse dimensioni, ma il patch embedding che separa due stadi cambia sia il numero dei token sia quello dei canali. Il modulo SPEDS risolve il problema inserendo una proiezione leggera $W_d$ sul ramo residuo,
\begin{equation}
\label{eq:speds_formulazione}
Y = \mathcal{F}(X, \{W_i\}) + \operatorname{SN}(W_d X),
\end{equation}
dove $\mathcal{F}$ è il ramo principale, una sequenza di convoluzioni, normalizzazioni, neuroni e max-pooling, e $W_d$ una convoluzione $1\times 1$ con passo maggiore di uno, che adegua in un solo colpo risoluzione e canali~\cite{zhou2024qkformer}. Da qui il nome di \emph{deformed shortcut}, perché il collegamento diretto non può più essere un'identità e va «deformato» per portare l'ingresso nello spazio dell'uscita. Inserito in Spikformer al posto del modulo originale, SPEDS porta l'accuratezza su CIFAR10-DVS dall'$80{,}90\%$ all'$82{,}20\%$~\cite{zhou2024qkformer}.

La gerarchia ha però un rovescio che conta per questa tesi. Rendendo economica l'attenzione e concentrando la riduzione di scala in moduli dedicati, sposta una parte del budget verso il patch embedding, che lavora sulle mappe a risoluzione più alta. Nella baseline compatta il primo patch embedding ospita oltre l'$80\%$ dello stato persistente, e i due patch embedding insieme circa due terzi dei parametri (sezione~\ref{subsec:bilancio_baseline}). Da questa osservazione parte l'analisi del capitolo~\ref{chap:riallocazione_budget}.

Per i dati neuromorfici gli autori impiegano una configurazione ridotta, con $1{,}50$ milioni di parametri, una larghezza massima di 256 canali e blocchi distribuiti sui tre stadi secondo lo schema $\{0,1,1\}$, nel quale il primo stadio si riduce al solo patch embedding. Con 16 timestep raggiunge il $98{,}6\%$ su DVS128 Gesture e l'$84{,}0\%$ su CIFAR10-DVS~\cite{zhou2024qkformer}. La baseline compatta scende a circa un terzo di questi parametri (tabella~\ref{tab:confronto_qkformer}), e in un regime simile la riduzione non può essere meccanica. Componenti che in QKFormer occupano una frazione modesta del budget possono diventare dominanti, e il peso di ciascuno va rivalutato quando larghezza e profondità diminuiscono.

\paragraph{Il ruolo del tempo}
\label{subsec:qkformer_tempo}

Nelle equazioni precedenti $Q$ e $K$ hanno una dimensione $T$, e si potrebbe pensare che l'attenzione operi congiuntamente sui timestep, ma gli autori precisano che il tempo resta separato soltanto negli strati di neuroni, mentre negli altri operatori viene fuso con la dimensione del batch~\cite{zhou2024qkformer}. Per questi operatori ogni timestep è un campione a sé, e a ciascun istante $t$ l'attenzione lavora sulla sola rappresentazione di quel passo,
\begin{equation}
\label{eq:attenzione_istante_t}
A^{\mathrm{tok}}_t = \mathcal{A}(Q_t).
\end{equation}
Il passato può influenzare questa decisione soltanto attraverso lo stato dei neuroni, compreso quello che produce il vettore, e non compare alcun termine in cui $Q_t$ venga messa a confronto con $Q_{t-1}$ o con istanti precedenti. QKFormer possiede dunque una struttura fortemente gerarchica nello spazio, ma nessuna gerarchia corrispondente lungo il tempo. Ciò non significa che il modello ignori l'informazione temporale, ma che l'interazione temporale non è l'oggetto del suo meccanismo di attenzione.

Per questa tesi la distinzione è essenziale, perché QKFormer fornisce un backbone in cui l'elaborazione spaziale e di canale è già stata progettata con cura per essere gerarchica ed efficiente, e consente così di porre una domanda più specifica di quella affrontata dal lavoro originale.
\begin{center}
\emph{Una volta portata l'architettura in un regime fortemente compatto, la sola memoria LIF è il modo più efficace di spendere il budget destinato alla dimensione temporale?}
\end{center}
Il lavoro originale non offre elementi per rispondere, e la risposta va cercata sperimentalmente. La baseline compatta nasce da una scelta conservativa. Mantiene i principi di QKFormer con fedeltà sufficiente a costituire un riferimento riconoscibile, ma ne riduce drasticamente la scala. Soltanto dopo averne caratterizzato il comportamento diventa possibile decidere quali componenti preservare, quali semplificare e quali sostituire. La sezione successiva esamina intanto le strategie proposte in letteratura per rendere esplicito il legame tra istanti diversi.


% ==============================================================================
%  Strategie per la modellazione temporale esplicita
% ==============================================================================


\section{Strategie per la modellazione temporale esplicita}
\label{sec:strategie_modellazione_temporale}

La sezione~\ref{subsec:memoria_relazioni} ha distinto la persistenza di una traccia del passato dall'accesso a ciò che è accaduto in un istante preciso, e la sezione~\ref{subsec:qkformer_tempo} ha mostrato che in QKFormer il passato raggiunge il presente soltanto attraverso lo stato dei neuroni. Il limite non è esclusivo di questa architettura, e Shi et al. hanno cercato di quantificarlo con il Temporal Interaction Coefficient (TIC)~\cite{shi2026multidelay}. Il coefficiente guarda a quanto l'uscita dell'attenzione al passo $q$ sia sensibile agli ingressi dei passi precedenti, e misura con un'entropia quanto questa sensibilità sia distribuita su più istanti. Su CIFAR10-DVS, per SSA, SDSA e QKTA, la dipendenza risulta fortemente concentrata sul passo corrente, mentre i passi lontani contribuiscono pochissimo. Il TIC descrive la forma della dipendenza e non la sua utilità, ma rende misurabile un'osservazione altrimenti qualitativa.

Le due strategie esaminate di seguito affiancano alla memoria dei neuroni un accesso esplicito al passato, e la figura~\ref{fig:confronto_integrazione_passato} le mette a confronto. Entrambe contano il tempo in passi, quindi il significato fisico delle loro distanze dipende dal clock dell'input (sezione~\ref{subsec:discretizzazione_clock}).

\begin{figure}[t]
    \centering
    \includegraphics[width=\textwidth]{images/cap1/cap1_accesso_passato.pdf}
    \caption[Tre modalità di accesso al passato]{Tre modalità di accesso al passato, con il tempo in orizzontale e il passo corrente $t$ a destra. In tutti i casi lo stato dei neuroni collega ogni passo al successivo. (a)~Con la sola memoria neuronale il passato arriva al presente compresso in un unico stato, e il peso di ciascun istante dipende soltanto dalla distanza, come indica l'intensità del colore. (b)~L'attenzione a blocchi confronta direttamente le feature degli istanti che appartengono allo stesso blocco, qui con $B=4$. (c)~Il mixing a ritardi legge alcuni istanti scelti, ciascuno con il proprio peso, qui con i ritardi di 1, 2, 4 e 8 passi del mixer usato in questo lavoro. Lo spessore dei collegamenti è illustrativo.}
    \label{fig:confronto_integrazione_passato}
\end{figure}

\paragraph{Attenzione spazio-temporale a blocchi}
\label{subsec:statten}

STAtten estende la SSA a più timestep~\cite{lee2025spiking}. Anziché calcolare l'attenzione separatamente per ciascun istante, raggruppa $B$ timestep consecutivi e dispone token e istanti lungo un unico asse, come se ogni coppia formata da una posizione e da un istante fosse un token a sé. Con $T=4$ e $B=2$, per esempio, vengono elaborati insieme i timestep $(0,1)$ e $(2,3)$. L'informazione del primo elemento di un blocco non deve più sopravvivere nello stato neuronale fino al secondo, perché le feature di entrambi partecipano alla stessa operazione. Di conseguenza l'operazione non è causale all'interno del blocco, dato che ogni istante vede anche quelli che lo seguono.

Si procede a blocchi anzitutto perché una correlazione sull'intera sequenza richiederebbe di conservare tutte le feature temporali, con un costo di memoria elevato. Inoltre il numero di neuroni attivi dopo il prodotto tra Query, Key e Value diminuisce con la distanza temporale, sicché correlare istanti molto lontani diventa poco produttivo~\cite{lee2025spiking}. La dimensione $B$ definisce così un orizzonte di interazione esplicita, oltre il quale la relazione torna a essere affidata ai neuroni. Sui dati neuromorfici gli autori usano $B=4$ con 16 timestep, e integrata in una configurazione di QKFormer a quattro blocchi STAtten guadagna un punto di accuratezza su CIFAR10-DVS. Il prezzo è che le feature di più timestep devono essere disponibili insieme, un costo che non compare nel conteggio dei parametri e che in un modello compatto pesa in modo particolare.

\paragraph{Mixing mediante ritardi}
\label{subsec:mixing_ritardi}

Il Multi-Delay Mixer (MD-Mixer) segue una strada diversa. Non apre l'attenzione a più timestep, ma modifica le feature che la raggiungono~\cite{shi2026multidelay}. Per ciascun canale $i$ considera $K$ ritardi $d_i^{(1)},\dots,d_i^{(K)}$ e altrettanti coefficienti, tutti appresi durante l'addestramento, e calcola
\begin{equation}
\label{eq:md_mixer_combinazione}
\widetilde X_{t,i} = \sum_{k=1}^{K} \alpha_i^{(k)} X_{t-d_i^{(k)},i}.
\end{equation}
Ogni canale attinge così a più istanti della propria storia, e canali diversi possono imparare ritardi diversi. Il cambiamento rispetto al LIF è qualitativo, perché nell'equazione~\eqref{eq:lif_espansione_temporale} è la distanza temporale a fissare il peso di ogni contributo, attraverso il decadimento, mentre qui è la rete a stabilire quanta importanza attribuire a ciascun ritardo, e un contributo lontano raggiunge il presente senza dover sopravvivere a tutti gli aggiornamenti intermedi del potenziale.

Nell'architettura completa il mixer prende il posto della proiezione lineare di Key e Value, mentre la Query resta allineata all'istante corrente. L'attenzione continua a lavorare un passo alla volta, ma su feature che contengono già una parte del passato. Integrato in una configurazione di QKFormer a due blocchi, MD-Mixer guadagna poco più di un punto di accuratezza su CIFAR10-DVS, contro meno di mezzo punto di STAtten nello stesso confronto. Sembra contare soprattutto la possibilità di scegliere quale passato leggere. In un esperimento di ablazione su Spike-driven Transformer, sullo stesso dataset, ritardi casuali e fissi portano circa mezzo punto, mentre ritardi appresi quasi due punti e mezzo~\cite{shi2026multidelay}. Poiché uno spostamento discreto non è derivabile rispetto al ritardo, durante l'addestramento ogni ritardo è una combinazione morbida di più spostamenti, che una temperatura decrescente rende via via più netta fino a una scelta discreta.

È l'approccio più vicino al mixer causale a più ritardi del modello di riferimento di questa tesi (capitolo~\ref{chap:riallocazione_budget}), che ne condivide l'idea di fondo ma se ne discosta in tre punti. Agisce nei patch embedding anziché su Key e Value, combina i canali tra loro con una matrice per ciascun ritardo anziché trattarli separatamente, e usa ritardi fissi di 1, 2, 4 e 8 passi, uguali per tutti i canali.

% TODO: eventuale parte aggiuntiva della 1.5 sugli operatori che confrontano passato e presente
% (correlatori di Hassenstein-Reichardt, termini bilineari). Per ora sono trattati nel Capitolo 5,
% sezione 5.1. Decidere a campagna conclusa; il paragrafo seguente va adattato di conseguenza.

Attenzione a blocchi e mixing a ritardi affiancano dunque la memoria neuronale con esigenze diverse di conservazione delle feature. Per confrontarli occorre guardare, oltre che alle prestazioni, a parametri, operazioni e stato persistente, le grandezze introdotte nella sezione successiva.


% ==============================================================================
%  Efficienza neuromorfica e vincoli hardware
% ==============================================================================

\section{Efficienza neuromorfica e vincoli hardware}
\label{sec:efficienza_neuromorfica_vincoli}

Un modello spiking viene spesso definito efficiente perché comunica mediante spike binari e presenta attivazioni sparse. Queste proprietà sono il presupposto di una computazione a basso consumo, ma non la garantiscono. Il vantaggio si materializza soltanto quando l'assenza di attività si traduce in lavoro effettivamente evitato, e dipende quindi anche dall'hardware su cui il modello viene eseguito. Questo lavoro non realizza i modelli su un dispositivo, ma caratterizza ogni architettura con grandezze che ne anticipano il costo. La sezione le introduce e ne chiarisce i limiti.

\paragraph{Dalle operazioni alla stima energetica}
\label{subsec:operazioni_energia}

La sezione~\ref{subsec:ssa} ha già introdotto la differenza tra una moltiplicazione seguita da un accumulo (MAC), necessaria quando l'attivazione è un valore a più bit, e un semplice accumulo (AC), che basta quando è uno spike binario. Il risparmio dipende però da quanti spike vengono emessi. Il firing rate $\rho$ misura la frazione di neuroni attivi, mediata sui passi. Una rete con $\rho=0{,}05$ attiva in media il $5\%$ dei neuroni a ogni timestep, mentre una con $\rho=0{,}60$ resta formalmente spiking, ma offre molte meno occasioni di saltare le operazioni associate agli zeri.

Il legame tra attività e costo si esprime con il numero di operazioni sinaptiche (\emph{synaptic operations}, SOP). Se uno strato $l$ richiederebbe $F_l$ operazioni in forma densa, in forma spiking ne esegue circa
\begin{equation}
\label{eq:stima_sop}
\operatorname{SOP}_l \approx \rho_l \, T \, F_l,
\end{equation}
perché le operazioni vengono ripetute a ogni timestep, ma soltanto la frazione corrispondente agli spike emessi dà luogo a calcolo. Il profilo di questo lavoro riporta separatamente il conteggio potenziale, in cui ogni possibile spike è presente, e quello pesato per l'attività osservata (sezione~\ref{sec:caratterizzazione_computazionale}).

QKFormer adotta questa impostazione e traduce i conteggi in una stima dell'energia, assegnando a ciascuna operazione un costo di riferimento~\cite{zhou2024qkformer},
\begin{equation}
\label{eq:energia_snn_semplificata}
E_{\mathrm{SNN}} \approx E_{\mathrm{MAC}}N_{\mathrm{MAC}} + E_{\mathrm{AC}}N_{\mathrm{AC}}.
\end{equation}
Nei valori di riferimento una MAC costa circa cinque volte un AC, e il risultato dipende quindi molto da come si attribuiscono le operazioni alle due voci. Nel calcolo di QKFormer le MAC provengono dal solo primo strato, che codifica l'ingresso in spike, e tutto il resto conta come accumulo. Il conteggio di questo lavoro è più severo, perché considera MAC ogni operazione con un operando non binario, comprese quelle che ricevono le somme residue descritte nella sezione~\ref{subsec:qk_attention}. Le stime dei due lavori non vanno perciò confrontate direttamente.

La formula va letta per quello che è, un modello aritmetico e non la misura del consumo di un sistema fisico. Restano fuori, tra l'altro, gli accessi in memoria, il trasferimento dei dati e la logica di controllo, e due modelli con la stessa stima possono comportarsi in modo molto diverso una volta realizzati. In questa tesi la stima serve a confrontare architetture a parità di input, e non viene mai presentata come energia misurata.

\paragraph{Stato persistente e buffer temporali}
\label{subsec:stato_buffer}

I pesi non sono l'unica memoria che un modello spiking occupa. Ogni neurone LIF deve conservare il proprio potenziale di membrana da un passo al successivo, e un'esecuzione in streaming deve leggerlo, aggiornarlo e riscriverlo a ogni timestep. Questo stato persistente segue il numero dei neuroni, e non quello dei pesi. Una convoluzione con pochi parametri, applicata a una mappa ad alta risoluzione, richiede un potenziale per ogni posizione e per ogni canale. Per questa ragione, nella baseline compatta, il primo patch embedding concentra la maggior parte dello stato (sezione~\ref{subsec:gerarchia_patch}) pur contenendo meno di un quinto dei parametri.

Una capacità temporale esplicita può accrescere ulteriormente questo costo. Un mixer che legge fino a $d_{\max}$ passi indietro deve conservare in un buffer le feature di quei passi, e occupa circa $d_{\max}$ volte lo spazio di una singola feature, mentre un'attenzione a blocchi deve avere disponibili insieme le feature di $B$ timestep. Per questo gli elementi di stato e i buffer temporali vengono riportati sempre accanto ai parametri. Un modulo temporale con pochi pesi può rivelarsi poco conveniente, se costringe a conservare molte feature storiche.

Ne discende la regola progettuale che guida il resto del lavoro. La capacità di un modello va giudicata per il luogo in cui viene spesa, non soltanto per la sua quantità. Ridurre i canali delle prime feature map può liberare pochi parametri, ma ridurre in misura molto maggiore i potenziali di membrana e le operazioni ripetute a ogni timestep. Allo stesso modo, un piccolo mixer temporale su feature a bassa risoluzione aggiunge pochi parametri e richiede un buffer ridotto, e può risultare più conveniente di un'attenzione temporale sulle feature ad alta risoluzione. In questo consiste un'allocazione della capacità consapevole dell'hardware (\emph{hardware-aware capacity allocation}), che il capitolo~\ref{chap:riallocazione_budget} applica alla baseline.

\paragraph{Dalla stima all'hardware}
I contatori descrivono il modello, non la macchina che lo esegue. Una GPU non è ottimizzata per la sparsità, perché le ampie regioni inattive dei tensori spiking non vengono sfruttate, e il continuo trasferimento dei dati tra memoria e unità di calcolo aggiunge un costo considerevole~\cite{karamimanesh2025fpga}. Il tempo di inferenza su GPU dice quindi poco sull'efficienza neuromorfica di un modello. L'hardware dedicato cerca invece di rendere visibile la sparsità. L'acceleratore di Li et al. per Spike-driven Transformer opera sugli indirizzi degli spike attivi, con addizioni e confronti anziché moltiplicazioni~\cite{li2024sparseacc}, e quello di Du et al. su FPGA non impiega alcun blocco DSP~\cite{du2024spike}. Il vantaggio dura però finché i circuiti che gestiscono la sparsità costano meno delle operazioni che permettono di evitare~\cite{karamimanesh2025fpga}.

Su una FPGA il costo si esprime infine in risorse del dispositivo, cioè logica (LUT), registri (FF), memoria (BRAM) e blocchi aritmetici (DSP), oltre che in latenza e potenza. Queste grandezze richiedono una sintesi che esula da questo lavoro. I conteggi qui descritti ne sono al più un'anticipazione, e le prospettive di una futura realizzazione vengono riprese nelle conclusioni.

Si chiude così il percorso del capitolo, che ha seguito l'informazione dal segnale prodotto dalla event camera alla sua discretizzazione, dalla memoria neuronale all'attenzione spiking, fino alla modellazione temporale esplicita e al suo costo. L'obiettivo non è costruire il modello con il minor numero di operazioni, né quello più accurato a qualunque costo, ma trovare una soluzione in cui l'informazione temporale aggiuntiva valga le risorse necessarie a rappresentarla. Il capitolo successivo applica questi criteri ai modelli effettivamente implementati, tenendo distinte le scelte relative all'input, i meccanismi architetturali e le condizioni di valutazione.


% Capitolo 2: baseline compatta e metodologia sperimentale.
\chapter{Baseline compatta e metodologia sperimentale}
\label{chap:baseline_metodologia}

% Capitolo rivisto il 2026-11-11 (figure da figure_scripts/cap2; statistiche appaiate da
% src/etsr/evaluation/comparison_metrics.py, bootstrap sulla differenza di accuratezza).
% Fonti verificate il 2026-11-04 sul repository temporal-event-spiking-research,
% branch Predictive-Temporal-Coding, commit 5a03fec: src/etsr/models/{mini_qkformer,layers,spiking,sequence}.py,
% src/etsr/encoders/count.py, configs/{models,recipes,datasets,splits}, config_resolved.yaml e
% summary.json dei run B E0 (dvslip 20260825_211710, dvsgesture 20261009_174645,
% cifar10dvs 20261009_184314, thu_eact_50_chl 20261011_093547), hardware_profile_v4.json di B su DVS-Lip,
% audit di sviluppo in artifacts/datasets/*. Nessun dato del test ufficiale.

Ogni modifica studiata nei capitoli successivi viene misurata rispetto a un riferimento, e quella misura vale quanto la precisione con cui il riferimento è stato fissato. Questo capitolo lo descrive, presentando i quattro dataset, il modo in cui il flusso di eventi diventa l'ingresso della rete, l'architettura della baseline compatta, il protocollo con cui i modelli vengono addestrati e confrontati e le grandezze con cui se ne misurano prestazioni e costo.

La parte sperimentale procede per passi successivi, ciascuno dei quali cambia un solo fattore. Il capitolo~\ref{chap:riallocazione_budget} ridistribuisce la capacità all'interno della baseline, a parità di ingresso, e ne ricava il modello di riferimento. Il capitolo~\ref{chap:clock_input} tiene ferma questa architettura e cambia il clock dell'input. Il capitolo~\ref{chap:memoria_relazionale}, con il clock così scelto, arricchisce la memoria temporale. Ciò che resta costante lungo questo percorso è descritto qui una volta per tutte.


% ==============================================================================
%  Dataset
% ==============================================================================

\section{Dataset}
\label{sec:dataset}

La sperimentazione si svolge su quattro benchmark neuromorfici di classificazione, scelti perché affidano all'informazione temporale ruoli molto diversi. In DVS-Lip le classi sono parole pronunciate, distinte da movimenti labiali brevi e spesso simili tra loro, e l'ordine in cui si susseguono le configurazioni della bocca è parte essenziale del segnale. DVS128 Gesture raccoglie gesti ampi e in larga parte ripetitivi, alcuni dei quali si distinguono soltanto per il verso del movimento. In CIFAR10-DVS la dinamica nasce dal modo in cui immagini statiche vengono mostrate al sensore, più che dal loro contenuto. THU-EACT-50-CHL, infine, contiene azioni umane complete, riprese a una risoluzione più alta e con durate molto variabili. Un meccanismo temporale utile in uno di questi scenari non lo è necessariamente negli altri, ed è proprio questa eterogeneità a rendere il confronto informativo. La figura~\ref{fig:esempi_dataset} ne mostra un campione per ciascuno.

\paragraph{DVS-Lip}
\label{subsec:dvslip}
È un dataset di lettura labiale a eventi, introdotto insieme alla rete MSTP~\cite{tan2022dvslip}. Raccoglie 100 parole inglesi pronunciate da 40 parlanti e rilasciate come regioni di $128\times128$ pixel attorno alla bocca, con una partizione ufficiale che assegna 30 parlanti al training e 10 al test~\cite{tan2022dvslip}. Le registrazioni durano circa un secondo e contengono pochi eventi, perché il segnale utile si concentra in un'area ristretta e in variazioni di forma molto fini. Il vocabolario include di proposito 25 coppie di parole visivamente simili, come \emph{spend} e \emph{spent}, \emph{million} e \emph{billion} o \emph{around} e \emph{ground}, che differiscono per un solo suono e producono movimenti labiali quasi identici. L'elenco completo è riportato in appendice, nella tabella~\ref{tab:coppie_estese}. Le 50 parole che le compongono e le restanti 50 formano due gruppi, sui quali si calcolano due accuratezze parziali, Acc1 e Acc2.

\paragraph{DVS128 Gesture}
\label{subsec:dvsgesture}
È stato acquisito con un sensore DVS128 e comprende 11 classi, dieci gesti di mani e braccia e una classe residua che raccoglie movimenti diversi~\cite{amir2017dvsgesture}. Le registrazioni coinvolgono 29 soggetti in diverse condizioni di illuminazione, e la partizione ufficiale riserva al test gli ultimi sei. Ogni registrazione contiene più gesti consecutivi, delimitati da annotazioni temporali, e in questo lavoro un campione coincide con un segmento annotato. I segmenti durano in mediana più di sei secondi, ma alcuni superano i diciotto. Diverse classi si distinguono soltanto per il braccio impiegato o per il verso di rotazione, e per questo trasformazioni come il ribaltamento orizzontale, che le scambierebbero tra loro, sono escluse.

\paragraph{CIFAR10-DVS}
\label{subsec:cifar10dvs}
Converte in eventi $10\,000$ immagini di CIFAR-10, mille per classe~\cite{li2017cifar10dvs}. Le immagini, mostrate su un monitor, vengono fatte muovere con un movimento regolare e ripetuto davanti a un sensore DVS128, così che i loro contorni generino eventi~\cite{li2017cifar10dvs}. Ne risultano registrazioni di durata quasi uniforme, poco più di un secondo, e le più dense del gruppo. Poiché il contenuto è statico, la struttura temporale riflette soprattutto il movimento imposto durante l'acquisizione, e il dataset fa da termine di confronto in cui è lecito attendersi che la memoria temporale pesi meno. Non esiste una partizione ufficiale, e in questo lavoro i primi cento campioni di ciascuna classe formano un insieme riservato, escluso dallo sviluppo.

\paragraph{THU-EACT-50-CHL}
\label{subsec:thueact}
THU-EACT-50 è un benchmark per il riconoscimento di 50 azioni umane con event camera~\cite{gao2023action}. Questo lavoro ne adotta la versione THU-EACT-50-CHL, acquisita con un sensore DAVIS346 a $346\times260$ pixel e composta da $2\,330$ registrazioni di 18 soggetti, che gli autori ripartiscono tra training e test senza soggetti in comune. Le durate hanno una coda lunga, da una mediana inferiore ai tre secondi fino a quasi venti, e il numero di eventi varia da circa diecimila a oltre quattro milioni per registrazione. La risoluzione maggiore moltiplica inoltre il numero di posizioni spaziali, e quindi di token, che la rete elabora a ogni passo.

\begin{figure}[t]
    \centering
    \includegraphics[width=\textwidth]{images/cap2/cap2_dataset.pdf}
    \caption[Un campione per ciascun dataset]{Un campione di sviluppo per ciascun dataset, scelto con una durata vicina alla mediana. A sinistra, tre bin della rappresentazione a finestra fissa della sezione~\ref{sec:input_baseline}, in rosso dove prevalgono gli eventi ON e in blu dove prevalgono gli OFF. A destra, la frequenza degli eventi sull'intera finestra, con le bande scure in corrispondenza dei bin mostrati e in grigio chiaro la parte successiva all'ultimo evento. I campioni sono la parola \emph{warning} per DVS-Lip, un saluto con la mano destra per DVS128 Gesture e un aereo per CIFAR10-DVS.}
    \label{fig:esempi_dataset}
\end{figure}

\paragraph{Partizioni di sviluppo}
\label{subsec:partizioni_sviluppo}

I test ufficiali non partecipano allo sviluppo, e ogni decisione, dalla configurazione dell'input alla scelta del checkpoint, si basa su una validazione ricavata dal solo training ufficiale, chiamata nel seguito validazione di sviluppo (\emph{development validation}). Le partizioni sono fissate prima di qualsiasi esperimento e sono le stesse per tutti i modelli, e ogni confronto avviene quindi sugli stessi campioni. La tabella~\ref{tab:dataset_sintesi} ne riporta le dimensioni insieme alle durate delle registrazioni.

\begin{table}[t]
\centering
\small
\renewcommand{\arraystretch}{1.2}
\caption[Dataset e partizioni di sviluppo]{Dataset e partizioni di sviluppo. La partizione riservata è il test ufficiale, tranne che per CIFAR10-DVS, per il quale è l'insieme escluso dallo sviluppo. Le durate si riferiscono alle partizioni di sviluppo e riportano il valore mediano e, tra parentesi, il massimo.}
\label{tab:dataset_sintesi}
\setlength{\tabcolsep}{4pt}
\begin{tabular}{@{}l c c r r r c@{}}
\hline
\textbf{Dataset} & \textbf{Classi} & \textbf{Risoluz.} & \textbf{Addestr.} & \textbf{Valid.} & \textbf{Riservati} & \textbf{Durata [s]} \\
\hline
DVS-Lip & 100 & $128\times128$ & $11\,901$ & $2\,995$ & $4\,975$ & $1{,}08$ ($1{,}89$) \\
DVS128 Gesture & 11 & $128\times128$ & 912 & 264 & 288 & $6{,}36$ ($18{,}46$) \\
CIFAR10-DVS & 10 & $128\times128$ & $8\,100$ & 900 & $1\,000$ & $1{,}30$ ($1{,}46$) \\
THU-EACT-50-CHL & 50 & $346\times260$ & $1\,387$ & 399 & 544 & $2{,}71$ ($19{,}91$) \\
\hline
\end{tabular}
\end{table}


% ==============================================================================
%  Rappresentazione dell'input nella baseline
% ==============================================================================

\section{Rappresentazione dell'input}
\label{sec:input_baseline}

La baseline riceve il flusso di eventi come frame di conteggio a finestra fissa (\emph{fixed-window count frames}), la più diretta tra le rappresentazioni introdotte nella sezione~\ref{sec:rappresentazione_event_camera}. A partire dall'origine del campione si considera una finestra di durata $D_{\mathrm{win}}$, la si divide in $T$ bin della stessa durata secondo l'equazione~\eqref{eq:finestre_fisse}, e in ciascun bin si contano, pixel per pixel, gli eventi ON e quelli OFF, come nell'equazione~\eqref{eq:event_frame_accumulo}. Il campione diventa così un tensore
\begin{equation}
\label{eq:input_baseline}
X=\bigl(F_0,\ldots,F_{T-1}\bigr)\in\mathbb N_0^{T\times2\times H\times W},
\qquad
\Delta t=\frac{D_{\mathrm{win}}}{T},
\end{equation}
con un canale per ciascuna polarità. Gli istogrammi di eventi a due canali sono tra le rappresentazioni più diffuse nell'apprendimento profondo su dati neuromorfici~\cite{maqueda2018event,gallego2022event}, e la rete li riceve così come sono, senza alcuna normalizzazione. Nel seguito questa scelta verrà indicata come rappresentazione a finestra fissa.

In tutti i dataset la sequenza conta $T=40$ passi, un budget comune di passi e non una risoluzione temporale comune, perché la finestra, e con essa la durata dei bin, viene scelta per ciascun dataset in base alle durate di sviluppo (tabella~\ref{tab:finestre_input}). Su tre dataset la finestra contiene per intero anche le registrazioni più lunghe. Su THU-EACT-50-CHL, dove la coda delle durate è molto lunga, una finestra estesa fino al campione più lungo avrebbe richiesto bin di mezzo secondo. Quella di $8\,\mathrm{s}$ conserva invece bin da $200\,\mathrm{ms}$ e il $99\%$ degli eventi, al prezzo di troncare due registrazioni della validazione.

\begin{table}[t]
\centering
\small
\renewcommand{\arraystretch}{1.2}
\caption[Parametri della rappresentazione a finestra fissa]{Parametri della rappresentazione a finestra fissa, con $T=40$ passi in tutti i dataset. Le ultime due colonne si riferiscono alle partizioni di sviluppo. La frazione di finestra occupata è il rapporto tra la durata mediana delle registrazioni e quella della finestra.}
\label{tab:finestre_input}
\begin{tabularx}{\textwidth}{@{}>{\raggedright\arraybackslash}X c c >{\centering\arraybackslash}p{2.0cm} >{\centering\arraybackslash}p{2.0cm}@{}}
\hline
\textbf{Dataset} & \textbf{Finestra} & \textbf{Bin} & \textbf{Eventi conservati} & \textbf{Finestra occupata} \\
\hline
DVS-Lip & $2\,\mathrm{s}$ & $50\,\mathrm{ms}$ & $100\%$ & $54\%$ \\
DVS128 Gesture & $20\,\mathrm{s}$ & $500\,\mathrm{ms}$ & $100\%$ & $32\%$ \\
CIFAR10-DVS & $2\,\mathrm{s}$ & $50\,\mathrm{ms}$ & $100\%$ & $65\%$ \\
THU-EACT-50-CHL & $8\,\mathrm{s}$ & $200\,\mathrm{ms}$ & $99{,}06\%$ & $34\%$ \\
\hline
\end{tabularx}
\end{table}

Una finestra prefissata ha però un effetto collaterale, ben visibile nella figura~\ref{fig:esempi_dataset}. Le registrazioni più brevi della finestra ne lasciano vuoti gli ultimi bin, che non sono padding ma osservazioni valide di un intervallo senza attività (sezione~\ref{subsec:tempo_valido_padding}). La rete li elabora comunque, e il readout ne include i contributi nella propria media. Su DVS128 Gesture e THU-EACT-50-CHL un campione tipico trascorre così circa due terzi dei passi senza ricevere alcun evento. Il capitolo~\ref{chap:clock_input} analizza le conseguenze di questa scelta e la confronta con altri clock. Fino ad allora la rappresentazione a finestra fissa resta invariata, perché ogni differenza osservata nel capitolo~\ref{chap:riallocazione_budget} possa essere attribuita all'architettura.


% ==============================================================================
%  Architettura della baseline compatta
% ==============================================================================

\section{Architettura della baseline compatta}
\label{sec:architettura_baseline}

La baseline compatta è un'implementazione scritta per questo lavoro, che conserva i principi di QKFormer (sezione~\ref{sec:qkformer_riferimento}) riducendone drasticamente la scala. Il backbone ha due stadi gerarchici (figura~\ref{fig:architettura_baseline}a). Un primo patch embedding porta l'ingresso a un ottavo della risoluzione con 64 canali, un blocco con Q-K Token Attention elabora questa rappresentazione, un secondo patch embedding dimezza ancora la risoluzione portando i canali a 128, e un blocco con spiking self-attention precede il readout. Le differenze rispetto all'architettura originale sono riassunte alla fine della sezione.

\begin{figure}[t]
    \centering
    \includegraphics[width=\textwidth]{images/cap2/cap2_architettura.pdf}
    \caption[Architettura della baseline compatta]{Architettura della baseline compatta per un ingresso di $128\times128$ pixel. (a)~Sequenza dei blocchi, con le dimensioni di uscita di ciascuno per ogni passo, nella forma canali $\times$ altezza $\times$ larghezza. (b)~Il primo patch embedding, l'unico modulo che lavora a piena risoluzione, con il ramo principale di quattro convoluzioni e la shortcut che parte dalla prima mappa di feature. (c)~Ripartizione dei parametri e dello stato persistente tra i blocchi su DVS-Lip, secondo la tabella~\ref{tab:architettura_baseline}. I colori dei blocchi corrispondono a quelli del pannello (a).}
    \label{fig:architettura_baseline}
\end{figure}

\paragraph{Neurone e convenzioni di calcolo}
\label{subsec:neurone_convenzioni}

Tutti i neuroni della rete sono LIF, discretizzati con lo schema di Eulero dell'equazione~\eqref{eq:lif_euler_discreto} con $\tau_d=2$ e dotati del soft reset dell'equazione~\eqref{eq:soft_reset} con soglia $V_{\mathrm{th}}=1$. Il potenziale viene azzerato all'inizio di ogni campione. Fanno eccezione i neuroni che prendono le decisioni dell'attenzione, descritti più avanti, la cui soglia è dimezzata. In addestramento la derivata dello spike è sostituita dalla surrogata sigmoidale dell'equazione~\eqref{eq:surrogate_sigmoide} con $\alpha=4$, e il termine di reset è escluso dal calcolo del gradiente.

Con $\tau_d=2$ il potenziale si dimezza a ogni passo, e la costante di tempo equivalente dell'equazione~\eqref{eq:tau_fisica_equivalente} vale circa $1{,}44\,\Delta t$. Lo stesso neurone dimentica quindi in tempi fisici molto diversi da un dataset all'altro, in circa $72\,\mathrm{ms}$ su DVS-Lip e CIFAR10-DVS, $290\,\mathrm{ms}$ su THU-EACT-50-CHL e $720\,\mathrm{ms}$ su DVS128 Gesture. È un primo segno di quanto il clock dell'input fissi anche l'orizzonte della memoria neuronale, un tema che il capitolo~\ref{chap:clock_input} riprende.

Convoluzioni, normalizzazioni e proiezioni lineari vengono applicate a ogni passo in modo indipendente, come se i passi fossero campioni distinti. L'unico canale attraverso cui l'informazione passa da un istante al successivo è quindi lo stato dei neuroni, come in QKFormer (sezione~\ref{subsec:qkformer_tempo}). Anche i collegamenti residui seguono QKFormer (sezione~\ref{subsec:qk_attention}). Sommano spike a valle dei neuroni, e i tensori scambiati tra i blocchi contengono perciò piccoli interi anziché valori binari, con conseguenze dirette sul conteggio delle operazioni (sezione~\ref{sec:caratterizzazione_computazionale}).

\paragraph{Primo patch embedding}
\label{subsec:primo_embedding}
Il primo patch embedding traduce i conteggi in spike e riduce la risoluzione di un fattore otto (figura~\ref{fig:architettura_baseline}b). Il ramo principale è una cascata di quattro convoluzioni $3\times3$, ciascuna seguita da batch normalization e da un neurone LIF. La prima porta i canali da 2 a 32 a piena risoluzione, le altre tre dimezzano a turno la risoluzione e portano i canali a 64. In parallelo, una shortcut preleva le feature del primo strato e le porta direttamente alla risoluzione finale con una convoluzione $1\times1$ di passo 8. Le due uscite, entrambe spiking, vengono sommate,
\begin{equation}
\label{eq:primo_embedding}
E^{(1)}_k=M^{(1)}_k+R^{(1)}_k\in\{0,1,2\}^{64\times\frac H8\times\frac W8},
\end{equation}
dove $M^{(1)}_k$ e $R^{(1)}_k$ sono gli spike emessi al passo $k$ dal ramo principale e dalla shortcut. È il principio della shortcut deformata di SPEDS (equazione~\eqref{eq:speds_formulazione}), con due differenze. Il sottocampionamento avviene con convoluzioni a passo 2 anziché con max-pooling, e la shortcut attraversa l'intero cambio di scala in un solo salto. È anche l'unico modulo della rete che opera a piena risoluzione, e il suo peso sul bilancio complessivo emerge più avanti.

\paragraph{Stadio con Q-K Token Attention}
\label{subsec:stadio_qkta}
Il primo stadio contiene un unico blocco Transformer con Q-K Token Attention. Per un ingresso di $128\times128$ pixel lavora su $16\times16=256$ token con 64 canali, divisi in 8 teste da 8 canali. Query e Key si ottengono come nell'equazione~\eqref{eq:qk_projections_qkformer}, e per ogni testa $h$ e ogni token $n$ la somma degli spike della Query sui canali della testa attraversa un neurone con soglia $0{,}5$,
\begin{equation}
\label{eq:qkta_baseline}
a_{h,n}=\operatorname{SN}_{0{,}5}\Bigl(\sum_{d\in h}Q_{n,d}\Bigr),
\qquad
X'_{n,d}=a_{h,n}\,K_{n,d}\quad\text{per } d\in h.
\end{equation}
La decisione binaria lascia passare o azzera i canali della Key di quella testa per quel token. Rispetto all'equazione~\eqref{eq:qkta_vettore_attenzione} ogni token dispone così di otto decisioni indipendenti, una per testa. Il neurone che le produce possiede inoltre uno stato, e la selezione di un token dipende anche dalla storia recente della sua Query. Seguono la proiezione di uscita dell'equazione~\eqref{eq:qk_proiezione_lineare_uscita}, il collegamento residuo e uno spiking MLP che espande i canali da 64 a 128 e li riporta a 64, secondo lo schema delle equazioni~\eqref{eq:blocco_qkformer_qka} e~\eqref{eq:blocco_qkformer_smlp}.

\paragraph{Secondo patch embedding e stadio con SSA}
\label{subsec:stadio_ssa}
Il secondo patch embedding segue lo stesso schema a due rami. Il ramo principale applica una convoluzione $3\times3$ da 64 a 128 canali e una seconda con passo 2, la shortcut una convoluzione $1\times1$ con passo 2. Ne risultano 128 canali a un sedicesimo della risoluzione d'ingresso, cioè 64 token. Il secondo stadio impiega la SSA, con Query, Key e Value spiking~\eqref{eq:spiking_qkv} divise in 8 teste da 16 canali. Per ciascuna testa
\begin{equation}
\label{eq:ssa_baseline}
Y=\operatorname{SN}_{0{,}5}\bigl(\beta\,Q(K^{\top}V)\bigr),
\qquad
\beta=\frac{1}{\sqrt{16}},
\end{equation}
cui seguono la proiezione di uscita, il collegamento residuo e uno spiking MLP che porta i canali da 128 a 256 e di nuovo a 128. In assenza di softmax il prodotto è associativo, e l'implementazione calcola prima $K^{\top}V$, una matrice $16\times16$ per testa. Il risultato coincide con quello dell'equazione~\eqref{eq:ssa_spikformer}, ma la matrice $N\times N$ tra token non viene mai costruita e il costo cresce linearmente con il numero dei token.

\paragraph{Readout}
\label{subsec:readout_baseline}
Il readout media la rappresentazione finale prima sulle posizioni e poi sui passi,
\begin{equation}
\label{eq:readout_media}
z=\frac{1}{T\,N_2}\sum_{k=0}^{T-1}\sum_{n=1}^{N_2}Z^{(2)}_{k,n},
\qquad
\hat y=W_c\,z+b_c,
\end{equation}
dove $Z^{(2)}_{k,n}\in\mathbb R^{128}$ è l'uscita del secondo stadio al passo $k$ e nella posizione $n$, $N_2$ il numero dei suoi token, $W_c$ e $b_c$ i parametri della testa lineare. Le feature mediate appartengono al flusso residuo e non sono binarie. La media comprende tutti i $T$ passi, anche quelli successivi all'ultimo evento, e la decisione riassume così l'intera sequenza, con lo stesso peso per ogni passo.

\paragraph{Bilancio di parametri e stato}
\label{subsec:bilancio_baseline}

La tabella~\ref{tab:architettura_baseline} riporta, blocco per blocco, le dimensioni, i parametri e lo stato persistente, cioè il numero di potenziali di membrana che devono sopravvivere da un passo al successivo. Il backbone conta $487\,808$ parametri, ai quali si aggiunge la testa, che dipende dal numero delle classi, per un totale che va da $489\,098$ su CIFAR10-DVS a $500\,708$ su DVS-Lip. THU-EACT-50-CHL ha invece un ingresso di $346\times260$ pixel, e su questo dataset i due stadi lavorano quindi su $1\,452$ e 374 token, con uno stato che cresce in proporzione al numero delle posizioni.

\begin{table}[t]
\centering
\small
\renewcommand{\arraystretch}{1.2}
\caption[Struttura della baseline compatta]{Struttura della baseline compatta per un ingresso di $128\times128$ pixel. Le uscite sono indicate come canali $\times$ altezza $\times$ larghezza per ciascun passo. Lo stato conta i potenziali di membrana di tutti i neuroni del blocco. I parametri della testa dipendono dal numero $C$ delle classi, e il totale si riferisce a DVS-Lip.}
\label{tab:architettura_baseline}
\begin{tabularx}{\textwidth}{@{}>{\raggedright\arraybackslash}p{2.6cm} >{\raggedright\arraybackslash}X c r r@{}}
\hline
\textbf{Blocco} & \textbf{Operazioni principali} & \textbf{Uscita} & \textbf{Parametri} & \textbf{Stato} \\
\hline
Primo patch embedding & 4 conv.\ $3\times3$ ($2\to32\to64\to64\to64$), shortcut $1\times1$ a passo 8 & $64\times16\times16$ & $95\,360$ & $884\,736$ \\
Stadio 1 & QKTA a 8 teste, MLP $64\to128\to64$ & $64\times16\times16$ & $29\,440$ & $100\,352$ \\
Secondo patch embedding & conv.\ $3\times3$ ($64\to128$), conv.\ $3\times3$ a passo 2, shortcut $1\times1$ a passo 2 & $128\times8\times8$ & $230\,144$ & $49\,152$ \\
Stadio 2 & SSA a 8 teste, MLP $128\to256\to128$ & $128\times8\times8$ & $132\,864$ & $65\,536$ \\
Readout e testa & media su spazio e tempo, lineare $128\to C$ & $C$ & $129\,C$ & 0 \\
\hline
\textbf{Totale} & DVS-Lip, $C=100$ & & $500\,708$ & $1\,099\,776$ \\
\hline
\end{tabularx}
\end{table}

La figura~\ref{fig:architettura_baseline}c mostra che parametri e stato si distribuiscono in modo quasi speculare. I due patch embedding contengono circa due terzi dei parametri, e il secondo da solo quasi la metà. Lo stato è invece dominato dal primo, che lavora sulle mappe più risolute e ne ospita oltre l'$80\%$, e nel complesso supera di più del doppio il numero dei parametri. Nel regime compatto capacità e costo si concentrano dunque lontano dall'attenzione, ed è da questa osservazione che parte il capitolo~\ref{chap:riallocazione_budget}.

\paragraph{Differenze rispetto a QKFormer}
\label{subsec:differenze_qkformer}

Rispetto alla configurazione neuromorfica di QKFormer (sezione~\ref{subsec:gerarchia_patch}), la baseline conserva la sequenza di un blocco QKTA e di un blocco SSA, il patch embedding con shortcut e la separazione tra memoria neuronale e interazioni all'interno del passo. Cambiano la scala e alcuni dettagli costruttivi, riassunti nella tabella~\ref{tab:confronto_qkformer}. La differenza nel numero degli stadi è meno marcata di quanto sembri, perché nella configurazione originale il primo stadio non contiene blocchi e si riduce al patch embedding. Più sostanziali sono la larghezza massima, dimezzata da 256 a 128 canali, il numero dei parametri, sceso a circa un terzo, e la lunghezza delle sequenze, di 40 passi anziché 16. Nell'articolo originale, tra i dataset neuromorfici, QKFormer è valutato soltanto su DVS128 Gesture e CIFAR10-DVS~\cite{zhou2024qkformer}, che si lasciano classificare bene anche con pochi passi. I 40 passi servono soprattutto a DVS-Lip e THU-EACT-50-CHL, dove con 16 passi ogni bin coprirebbe $125\,\mathrm{ms}$ di una parola che dura circa un secondo o mezzo secondo di un'azione, e vengono mantenuti su tutti i dataset perché il budget di passi resti comune.

Un'ultima differenza riguarda l'implementazione della SSA, che come gli altri operatori lavora con il tempo fuso nella dimensione del batch (sezione~\ref{subsec:qkformer_tempo}), ma il neurone che ne produce l'uscita deve far evolvere il proprio potenziale lungo i passi di uno stesso campione, e prima di raggiungerlo i due assi vanno separati di nuovo. Nel codice ufficiale di QKFormer la separazione avveniva dopo il neurone anziché prima, e il neurone finiva così per mescolare i potenziali di campioni diversi. Gli autori hanno segnalato l'errore nel repository, indicando un effetto rilevante soprattutto sui risultati su ImageNet~\cite{qkformerImplementationNotice}. Nella baseline tempo e batch vengono separati prima del neurone, e il problema non si presenta.

\begin{table}[t]
\centering
\small
\renewcommand{\arraystretch}{1.2}
\caption[QKFormer neuromorfico e baseline compatta a confronto]{Confronto tra la configurazione neuromorfica di QKFormer, secondo la descrizione della sezione~\ref{sec:qkformer_riferimento}, e la baseline compatta.}
\label{tab:confronto_qkformer}
\begin{tabularx}{\textwidth}{@{}>{\raggedright\arraybackslash}p{3.8cm} >{\raggedright\arraybackslash}X >{\raggedright\arraybackslash}X@{}}
\hline
\textbf{Aspetto} & \textbf{QKFormer, configurazione neuromorfica} & \textbf{Baseline compatta} \\
\hline
Stadi e blocchi & tre stadi, blocchi $\{0,1,1\}$ & due stadi, un blocco ciascuno \\
Attenzione & QKTA, poi SSA & QKTA, poi SSA \\
Larghezza massima & 256 canali & 128 canali \\
Parametri & circa $1{,}50$ milioni & da $489\,098$ a $500\,708$ \\
Sottocampionamento & convoluzione, BN, neurone e max-pooling & convoluzioni a passo 2 \\
Passi temporali & 16 & 40 \\
\hline
\end{tabularx}
\end{table}


% ==============================================================================
%  Protocollo sperimentale
% ==============================================================================

\section{Protocollo sperimentale}
\label{sec:protocollo_addestramento}

Tutti i modelli vengono addestrati da zero con la stessa ricetta, e due modelli messi a confronto differiscono per un solo elemento. Questa sezione descrive la ricetta, i vincoli che ne derivano per i confronti e i modi in cui un componente della rete può essere interrogato.

\paragraph{Ricetta di addestramento}
\label{subsec:ricetta}
La ricetta, riassunta nella tabella~\ref{tab:ricetta_addestramento}, è comune a tutti i modelli e a tutti i dataset. L'ottimizzatore è AdamW~\cite{loshchilov2019adamw}, con un tasso di apprendimento che cresce nelle prime epoche e decresce poi con un profilo cosinusoidale~\cite{loshchilov2017sgdr} lungo 128 epoche. La funzione di perdita è l'entropia incrociata con label smoothing~\cite{szegedy2016rethinking}, la norma del gradiente viene limitata~\cite{pascanu2013difficulty} e i calcoli avvengono in precisione mista~\cite{micikevicius2018mixed}. La ricetta non viene adattata ai singoli modelli né ai singoli dataset, e quando una campagna la modifica, la modifica vale per tutti i modelli del confronto.

\paragraph{Microbatch e batch normalization}
\label{subsec:microbatch_bn}
Ogni aggiornamento dei pesi usa un batch di 32 campioni. Quando la memoria non basta a elaborarli tutti insieme, il batch viene diviso in microbatch più piccoli, elaborati uno dopo l'altro, e i loro gradienti vengono sommati prima dell'aggiornamento. Per gli strati convolutivi e lineari il risultato è lo stesso che si otterrebbe con un unico batch da 32 campioni. Per la batch normalization~\cite{ioffe2015batch} non lo è, perché in addestramento questo strato normalizza ogni canale con la media e la varianza dei dati che sta elaborando in quel momento, cioè del solo microbatch corrente, e statistiche stimate su 8 campioni non coincidono con quelle stimate su 32. La dimensione dei microbatch, che cambia da un dataset all'altro con la memoria richiesta, fa quindi parte del protocollo, e ogni eccezione viene dichiarata dove si presenta.

Una seconda conseguenza, meno evidente, nasce dal fatto che il tempo è fuso con la dimensione del batch (sezione~\ref{sec:architettura_baseline}), sicché media e varianza vengono calcolate insieme su tutti i passi di tutti i campioni del microbatch, come avviene comunemente nelle SNN addestrate direttamente~\cite{zheng2021going}. In addestramento, quindi, la normalizzazione di un passo dipende anche dai passi successivi della stessa sequenza, e la rete nel suo complesso non è strettamente causale, anche se lo è ciascuno dei suoi operatori. In inferenza la batch normalization usa invece statistiche fisse, accumulate durante l'addestramento, e ogni passo torna a dipendere soltanto da quelli precedenti.

\begin{table}[t]
\centering
\small
\renewcommand{\arraystretch}{1.2}
\caption[Ricetta di addestramento comune]{Ricetta di addestramento comune a tutti i modelli e a tutti i dataset.}
\label{tab:ricetta_addestramento}
\begin{tabularx}{\textwidth}{@{}>{\raggedright\arraybackslash}p{3.6cm} >{\raggedright\arraybackslash}X@{}}
\hline
\textbf{Componente} & \textbf{Scelta} \\
\hline
Ottimizzatore & AdamW, tasso di apprendimento $3\times10^{-4}$, weight decay $5\times10^{-4}$ \\
Schedule & warmup lineare di 4 epoche da $0{,}01$ volte il valore nominale, poi decadimento cosinusoidale fino a $10^{-6}$ \\
Durata & 128 epoche \\
Funzione di perdita & entropia incrociata con label smoothing $0{,}1$ \\
Stabilizzazione & limite di $1$ sulla norma globale del gradiente \\
Precisione & precisione mista automatica su GPU \\
Batch effettivo & 32 campioni, ottenuti per accumulo su più microbatch \\
Selezione del checkpoint & epoca con la massima Macro-F1 sulla validazione di sviluppo \\
\hline
\end{tabularx}
\end{table}

\paragraph{Test di overfitting}
\label{subsec:verifiche_preliminari}
Prima di un addestramento completo, ogni nuova configurazione deve superare un test di overfitting controllato. Il modello viene addestrato su pochi campioni, quattro per ciascuna di al più 16 classi, e deve imparare a classificarli quasi perfettamente, mantenendo per cinque epoche consecutive un'accuratezza superiore al $95\%$. È una verifica diffusa nella pratica dell'apprendimento profondo~\cite{goodfellow2016deep}, e il suo principio è semplice. Una rete con centinaia di migliaia di parametri dovrebbe riuscire a memorizzare qualche decina di esempi senza difficoltà, e se non ci riesce il problema non sta nella generalizzazione ma nel modello o nella sua implementazione, per esempio in un errore nel flusso dei dati, in un gradiente che non raggiunge alcuni strati o in una dinamica che ostacola l'ottimizzazione. Il test non dice quindi nulla sulla capacità di generalizzare, ma mostra che il modello è in grado di apprendere, e intercetta a un costo trascurabile errori che altrimenti emergerebbero soltanto alla fine di un addestramento completo. Una configurazione che non lo supera non viene addestrata.

\paragraph{Confronti tra modelli}
\label{sec:disegno_controlli}
Due modelli messi a confronto differiscono per un solo elemento, mentre dati, partizioni, ingresso, ricetta e regola di selezione restano identici. Quando l'oggetto dell'esperimento è l'ingresso, è l'architettura a restare ferma. Lo stesso componente, però, può essere interrogato in tre modi, che rispondono a domande diverse.

Un'ablazione riaddestrata rimuove o sostituisce il componente e addestra il modello da zero. Misura che cosa riesca a ottenere l'architettura senza di esso, ed è la sola prova che il componente serva davvero.

Un intervento sul checkpoint modifica invece un modello già addestrato, soltanto in inferenza. Spegnere un ramo, eliminare un ritardo o rimescolare l'ordine dei passi misura quanto la soluzione appresa dipenda da quel componente, non se il componente sia necessario. Una rete addestrata con un modulo può essersi adattata alla sua presenza e degradare quando lo si rimuove, anche se un'architettura che ne è priva avrebbe raggiunto da zero un risultato analogo. Il ruolo dell'attenzione nel secondo stadio viene esaminato combinando i due punti di vista, con un test di necessità che la disattiva sul checkpoint addestrato e un test di sostituibilità che la rimpiazza con un mixer spaziale locale e riaddestra il modello.

Un controllo di capacità, infine, separa l'effetto di un operatore da quello delle risorse che aggiunge. Ogni nuovo modulo porta con sé parametri e operazioni, e un miglioramento potrebbe dipendere soltanto da questi. Il confronto con un'alternativa di costo comparabile ma priva della struttura in esame, come un allargamento uniforme della rete o un percettrone multistrato con lo stesso numero di parametri, permette di attribuire il guadagno alla struttura e non alla quantità.


% ==============================================================================
%  Valutazione e incertezza
% ==============================================================================

\section{Valutazione e incertezza}
\label{sec:metriche_valutazione}

\paragraph{Metriche}
\label{subsec:macro_f1}
La metrica principale è la Macro-F1, cioè la media non pesata, sulle $C$ classi, degli F1 delle singole classi~\cite{sokolova2009systematic},
\begin{equation}
\label{eq:macro_f1}
\text{Macro-F1}=\frac{1}{C}\sum_{c=1}^{C}\frac{2P_cR_c}{P_c+R_c},
\end{equation}
dove $P_c$ e $R_c$ sono la precisione e il richiamo della classe $c$. A differenza dell'accuratezza, la Macro-F1 attribuisce a tutte le classi lo stesso peso, indipendentemente dalla loro numerosità, e penalizza un modello che ottenga un buon risultato complessivo trascurando alcune classi. L'accuratezza conta infatti i campioni e non le classi, e quando poche classi raccolgono la maggior parte dei campioni basta riconoscere bene quelle per ottenere un valore alto. Se una classe comprendesse il $90\%$ dei campioni, un modello che la predicesse sempre raggiungerebbe un'accuratezza del $90\%$ senza aver imparato nulla delle altre, mentre la sua Macro-F1 resterebbe molto bassa, perché tutte le altre classi avrebbero un F1 nullo.\footnote{Su THU-EACT-50-CHL la media comprende tutte le 50 classi della testa, anche se quattro di esse compaiono soltanto in soggetti del training e mancano sia nella validazione di sviluppo sia nel test ufficiale. Il loro F1 è nullo per costruzione, e il valore massimo raggiungibile della Macro-F1 è del $92\%$. Il vincolo è identico per tutti i modelli e non altera i confronti, ma va tenuto presente nella lettura dei valori assoluti.} Accanto a essa si riportano l'accuratezza e, su DVS-Lip, le accuratezze Acc1 e Acc2 dei due gruppi di parole descritti nella sezione~\ref{sec:dataset}.

\paragraph{Selezione del checkpoint}
\label{subsec:selezione_incertezza}
Ogni modello viene valutato sulla validazione di sviluppo al termine di ciascuna epoca, e il checkpoint conservato è quello con la Macro-F1 più alta. Selezione e misura avvengono quindi sugli stessi dati, e il valore riportato è una stima ottimistica della prestazione su dati nuovi, tanto più quanto più la curva di validazione oscilla. Per questo si registrano anche la media e la deviazione standard della Macro-F1 nelle ultime 16 epoche, un indicatore di stabilità meno esposto alla scelta del massimo.

\paragraph{Incertezza dei confronti}
Quando due modelli vengono valutati sulla stessa validazione, le loro predizioni sono appaiate campione per campione, e il confronto può sfruttare questa corrispondenza con due strumenti che rispondono a domande complementari.

Il test esatto di McNemar chiede se la differenza tra i due modelli possa essere nata per caso~\cite{mcnemar1947note,dietterich1998approximate}. Trascura i campioni su cui i modelli sono entrambi corretti o entrambi in errore, che non dicono nulla sulla loro differenza, e guarda soltanto quelli su cui discordano. Se i due modelli fossero equivalenti, i campioni classificati correttamente soltanto dal primo e quelli classificati correttamente soltanto dal secondo dovrebbero dividersi all'incirca a metà, come gli esiti di una moneta lanciata più volte, e il test misura quanto sarebbe improbabile, in quell'ipotesi, lo sbilanciamento osservato. Il bootstrap appaiato stima invece quanto sia grande la differenza e con quale margine~\cite{efron1993bootstrap}. Costruisce molte validazioni alternative estraendo con reinserimento i campioni di quella disponibile, ciascuno con le predizioni di entrambi i modelli, ricalcola su ognuna la differenza di accuratezza e ne ricava un intervallo di confidenza al $95\%$. Se l'intervallo non comprende lo zero, il vantaggio di un modello difficilmente dipende dalla particolare scelta dei campioni di validazione. I valori di $p$ non sono corretti per i confronti multipli, e vanno letti come indicazioni più che come verdetti.

Entrambi gli strumenti misurano l'incertezza dovuta alla dimensione finita della validazione, a modelli fissati. Nessuno dei due misura la variabilità dell'addestramento, cioè quanto cambierebbe il risultato ripetendolo con un'inizializzazione diversa, che può essere dello stesso ordine delle differenze tra i modelli confrontati~\cite{bouthillier2021accounting}. Le esplorazioni usano perciò un solo addestramento per configurazione, e i loro risultati sono trattati come indicazioni. Il modello di riferimento del capitolo~\ref{chap:riallocazione_budget} e i modelli con due milioni di parametri del capitolo~\ref{chap:scalabilita} vengono addestrati tre volte, e solo per questi si riportano media e deviazione standard.

\paragraph{Valutazione sui prefissi}
\label{subsec:prefissi}
Si misura anche come la decisione si formi nel tempo. Lo stesso checkpoint viene valutato su porzioni iniziali della sequenza, con il readout calcolato sui soli passi osservati, e un indice sintetico, indicato come PrefixAUC, riassume l'andamento dell'accuratezza al crescere della porzione osservata. È una misura diagnostica, che non interviene nella selezione dei modelli e non autorizza a ridurre la finestra di ingresso, perché un modello addestrato su sequenze complete non è stato ottimizzato per decidere su sequenze troncate.

\paragraph{Partizioni riservate}
\label{subsec:partizioni_riservate}
Tutte le metriche riportate nei capitoli successivi sono misurate sulla validazione di sviluppo, salvo indicazione esplicita. Le partizioni riservate restano escluse dallo sviluppo, dalla selezione dei checkpoint e dalla caratterizzazione computazionale. L'unica eccezione riguarda il test ufficiale di DVS-Lip e THU-EACT-50-CHL, valutato in sola inferenza e soltanto a posteriori, su checkpoint già selezionati. Serve a descrivere il comportamento dei modelli scelti (sezioni~\ref{sec:uso_memoria} e~\ref{sec:errori}) e a confrontare i modelli finali con la letteratura (sezione~\ref{sec:configurazioni_finali}), e non viene usato per sceglierne altri.


% ==============================================================================
%  Caratterizzazione computazionale
% ==============================================================================

\section{Caratterizzazione computazionale}
\label{sec:caratterizzazione_computazionale}

Il costo di ciascun modello viene misurato in inferenza sul checkpoint selezionato, usando un sottoinsieme fisso di 64 campioni della validazione di sviluppo, lo stesso per tutti i modelli di un dataset. Le grandezze raccolte sono quelle introdotte nella sezione~\ref{sec:efficienza_neuromorfica_vincoli}, e la tabella~\ref{tab:contatori_profilo} ne precisa il significato.

\begin{table}[t]
\centering
\small
\renewcommand{\arraystretch}{1.2}
\caption[Grandezze della caratterizzazione computazionale]{Grandezze raccolte nella caratterizzazione computazionale di ciascun checkpoint.}
\label{tab:contatori_profilo}
\begin{tabularx}{\textwidth}{@{}>{\raggedright\arraybackslash}p{3.2cm} >{\raggedright\arraybackslash}X@{}}
\hline
\textbf{Grandezza} & \textbf{Che cosa conta} \\
\hline
Parametri & pesi addestrabili, testa inclusa \\
Stato persistente & potenziali di membrana e buffer temporali che devono sopravvivere da un passo al successivo \\
MAC & moltiplicazioni e accumuli con un operando non binario, come i conteggi in ingresso, le somme residue e le correnti dei moduli temporali \\
AC potenziali & accumuli con ingresso binario, contati come se ogni possibile spike fosse presente \\
SOP & accumuli effettivamente innescati dagli spike osservati, comprese le operazioni dell'attenzione, secondo l'equazione~\eqref{eq:stima_sop} \\
Firing rate & frazione di spike emessi da ciascuno strato \\
\hline
\end{tabularx}
\end{table}

La distinzione tra MAC e accumuli segue il tipo di operando, non il tipo di strato. Uno strato che legge spike esegue accumuli, mentre lo stesso strato applicato a conteggi o a somme residue esegue moltiplicazioni vere e proprie. Per questo i collegamenti residui della baseline, che producono piccoli interi (sezione~\ref{sec:architettura_baseline}), pesano sul bilancio, perché una parte delle proiezioni della rete riceve valori non binari.

Per riassumere i conteggi in un'unica grandezza si adotta la stima dell'equazione~\eqref{eq:energia_snn_semplificata}, con i costi di riferimento per operazioni a 32 bit in tecnologia a $45\,\mathrm{nm}$ riportati da Horowitz~\cite{horowitz2014energy}, cioè $4{,}6\,\mathrm{pJ}$ per una MAC, $0{,}9\,\mathrm{pJ}$ per un'addizione e $3{,}7\,\mathrm{pJ}$ per una moltiplicazione. È la convenzione adottata dalla maggior parte dei lavori sugli spiking Transformer, da Spikformer e Spike-driven Transformer fino a QKFormer~\cite{zhou2023spikformer,yao2023spikedriven,zhou2024qkformer}, perché si basa su costi pubblicati, non dipende da un particolare acceleratore e si applica a qualsiasi rete a partire dal solo conteggio delle operazioni. Adottarla rende le stime di questo lavoro confrontabili con quelle della letteratura. La stima viene calcolata in due varianti, una densa, basata sugli accumuli potenziali, e una pesata per l'attività osservata. Come discusso nella sezione~\ref{sec:efficienza_neuromorfica_vincoli}, è una proxy parziale del solo costo aritmetico, utile per confrontare architetture a parità di input e in nessun caso una misura dell'energia consumata. Lo stato persistente, infine, viene riportato sempre accanto ai parametri, perché nella baseline le due grandezze differiscono di più di un fattore due e un modulo con pochi pesi può richiedere un buffer molto più grande.


% ==============================================================================
%  Il riferimento nei capitoli successivi
% ==============================================================================

\section{Il riferimento nei capitoli successivi}
\label{sec:sintesi_protocollo}

Dataset e partizioni, ricetta di addestramento, regola di selezione del checkpoint, metriche e grandezze di costo restano gli stessi in tutti i capitoli sperimentali. Cambiano invece, da un capitolo all'altro, la rete di partenza e l'ingresso, secondo il percorso descritto all'inizio del capitolo e riassunto nella tabella~\ref{tab:mappa_riferimenti}. In ogni capitolo varia un solo fattore, e i confronti avvengono sempre con il riferimento di quel capitolo. I dettagli necessari a ripetere gli esperimenti, dal frazionamento del batch agli iperparametri dei singoli moduli, sono raccolti nell'appendice~\ref{app:riproducibilita}.

\par\medskip\noindent
\begin{minipage}{\textwidth}
\small
\renewcommand{\arraystretch}{1.2}
\captionof{table}[Il riferimento nei capitoli sperimentali]{Rete di partenza, ingresso e fattore studiato nei capitoli sperimentali. Tutti gli altri elementi del protocollo restano quelli descritti in questo capitolo.}
\label{tab:mappa_riferimenti}
\begin{tabularx}{\textwidth}{@{}>{\raggedright\arraybackslash}p{3.3cm} >{\raggedright\arraybackslash}p{3.0cm} >{\raggedright\arraybackslash}X >{\raggedright\arraybackslash}X@{}}
\hline
\textbf{Capitolo} & \textbf{Rete di partenza} & \textbf{Ingresso} & \textbf{Fattore studiato} \\
\hline
\ref{chap:riallocazione_budget}, riallocazione del budget & baseline compatta & finestra fissa & distribuzione della capacità nella rete \\
\ref{chap:clock_input}, clock dell'input & modello di riferimento & finestra fissa, eventi costanti, lunghezza variabile & rappresentazione dell'ingresso \\
\ref{chap:memoria_relazionale}, memoria relazionale & modello di riferimento & lunghezza variabile & memoria temporale \\
\hline
\end{tabularx}
\end{minipage}
\par\medskip

Il capitolo successivo parte dal bilancio della sezione~\ref{sec:architettura_baseline}, che concentra stato e operazioni lontano dall'attenzione, e si chiede come spendere meglio un budget di capacità così ridotto.


% Capitolo 3: in attesa delle riesecuzioni che isolano i singoli interventi.
\chapter{Riallocazione del budget nella baseline compatta}
\label{chap:riallocazione_budget}

% Fonti verificate il 2026-11-11 sul repository temporal-event-spiking-research:
% notebooks/evoluzione_architetturale.ipynb e docs/EVOLUZIONE_DELLA_SPERIMENTAZIONE_v2.md;
% summary.json dei run DVS-Lip E0 seed 42: E0_B_M__20260825_211710, E0_B_M--pyramidal__20260908_094118_261095,
% E0_C0_M--f-temporal-capacity-d8__20260913_161320_867938 (4 ritardi, Token-QK), E0_C0_M__20260914_093427_434930,
% varianti E0_B_M--{no-cross-time,readout-time-last-event,capacity-1m,capacity-2m},
% profili comuni in artifacts/reports/architecture_evolution/<run>/hardware_profile_v4.json (64 campioni, FP32, CPU);
% Sezioni 3.3-3.6: ripetizioni E0_{B,C0}_M seed 42-44 (DVS-Lip), E0_{B,C0}_M sui quattro dataset con le ripetizioni
% disponibili del riferimento, varianti --capacity-{1m,2m}; necessità in artifacts/reports/necessity/<dataset>/,
% sostituibilità E0_C0_M--ssa-local (Gesture, CIFAR10-DVS); E0_C0_M--f-tcap-stage1-dwc3-learnable-delays,
% E0_D_M (gate di contenuto, solo patch_embed2) ed E0_D_M--static-control; profili hardware_profile_v4.json dei run.
% codice: src/etsr/models/{layers,temporal,mini_qkformer,local_mixer,diagnostics}.py. Figure da figure_scripts/cap3.
% Aggiunte 2026-11-18 (revisione sulle note del PDF): righe a 192 canali della tabella tab:capacita_lip dai
% hardware_profile_v4.json dei run E0_{B,C0}_M--capacity-1m (DVS-Lip, 64 campioni, stima = covered activity proxy);
% sostituibilità su DVS-Lip e THU: artifacts/overfit/<dataset>/E0_C0_M--ssa-local__* (overfit gate fallito dopo 500
% epoche, Lip max 4 epoche consecutive >= 95%, THU max 64%; il riferimento E0_C0_M lo supera in 310 e 429 epoche);
% figura cap3_ritardi_appresi da learned_delay_trajectory.csv del run learnable-delays (best_epoch 125);
% Max-Former (fang2025maxformer): configurazione neuromorfica Embed-Max+ / DWC-3 nello stadio 1, SSA nello stadio 2.

La baseline compatta dispone di circa mezzo milione di parametri, e con un budget così ridotto conta meno quanta capacità si aggiunge e molto di più dove la si colloca. Questo capitolo parte da una diagnosi della baseline, che mostra dove si concentrano le operazioni e che cosa limita le prestazioni, e procede con tre modifiche successive, ciascuna valutata cambiando un solo elemento. Il risultato è il modello di riferimento usato nei capitoli successivi, che su DVS-Lip guadagna undici punti di Macro-F1 con appena 320 parametri in più. Le sezioni finali ne verificano la solidità, esaminano il ruolo dell'attenzione nel secondo stadio e discutono due varianti in cui i ritardi diventano adattivi.

In tutto il capitolo l'ingresso è la rappresentazione a finestra fissa della sezione~\ref{sec:input_baseline}, la ricetta è quella della sezione~\ref{sec:protocollo_addestramento} e le prestazioni sono misurate sulla validazione di sviluppo.


% ==============================================================================
%  Diagnosi della baseline
% ==============================================================================

\section{Diagnosi della baseline}
\label{sec:diagnosi_baseline}

Per capire dove si concentri il costo, la figura~\ref{fig:diagnosi_baseline}a ripartisce tra i blocchi le operazioni potenziali della baseline su DVS-Lip, cioè quelle che la rete eseguirebbe se ogni possibile spike fosse presente. I due patch embedding ne concentrano oltre il $90\%$. Il primo, che lavora a piena risoluzione, raccoglie quasi nove decimi degli accumuli, e insieme al secondo totalizza tre quarti delle moltiplicazioni. Queste ultime nascono dove gli operandi non sono binari, perché il primo strato riceve i conteggi degli eventi, il secondo patch embedding riceve il flusso residuo del primo stadio, che contiene piccoli interi (sezione~\ref{sec:architettura_baseline}), e ne eredita da solo metà delle MAC dell'intera rete. I due stadi con l'attenzione, che danno il nome all'architettura, contano per un quarto delle moltiplicazioni e per meno del $5\%$ degli accumuli. Lo stesso squilibrio, come si è visto nella figura~\ref{fig:architettura_baseline}c, riguarda lo stato persistente, dominato dal primo patch embedding.

\begin{figure}[t]
    \centering
    \includegraphics[width=\textwidth]{images/cap3/cap3_diagnosi.pdf}
    \caption[Diagnosi della baseline su DVS-Lip]{Diagnosi della baseline su DVS-Lip. (a)~Ripartizione tra i blocchi delle moltiplicazioni con operando non binario (MAC) e degli accumuli potenziali per campione. (b)~Macro-F1 della baseline e di quattro varianti, ciascuna addestrata da zero con un solo cambiamento. Senza memoria tra i passi i neuroni non conservano lo stato da un passo al successivo. La variante con readout fino all'ultimo evento esclude dalla media i passi della coda vuota. Le due varianti più larghe portano a 192 e 256 i canali della rete, che conta così circa uno e due milioni di parametri. La linea tratteggiata indica la baseline.}
    \label{fig:diagnosi_baseline}
\end{figure}

Quanto alla memoria, se i neuroni non conservano il proprio stato da un passo al successivo, la rete può soltanto mediare le feature dei singoli passi, e la Macro-F1 crolla da $44{,}2$ a $13{,}0$ (figura~\ref{fig:diagnosi_baseline}b). Nella baseline, dunque, quasi tutta l'informazione temporale utile passa attraverso i potenziali di membrana, l'unico canale che collega un istante al successivo (sezione~\ref{sec:architettura_baseline}). Conta molto meno il momento in cui si legge la decisione, perché escludere dalla media del readout i passi successivi all'ultimo evento costa circa due punti, un primo segnale del fatto che la rete fa uso anche della coda vuota, su cui torna il capitolo~\ref{chap:clock_input}.

Quanto alla capacità, allargando la rete a 192 e a 256 canali, con circa uno e due milioni di parametri, la Macro-F1 sale a $49{,}4$ e a $51{,}8$. Quadruplicare il budget frutta meno di otto punti, e lo fa amplificando proprio i blocchi che già concentrano il costo. Insieme, le tre osservazioni suggeriscono di spostare la capacità invece di aumentarla, togliendo risorse al frontend, che costa molto, e investendole in una memoria temporale più esplicita di quella offerta dai soli neuroni. La sezione successiva realizza questo spostamento un passo alla volta.


% ==============================================================================
%  Dalla baseline al modello di riferimento
% ==============================================================================

\section{Dalla baseline al modello di riferimento}
\label{sec:percorso_riferimento}

Il percorso si compone di tre modifiche, mostrate nella figura~\ref{fig:modello_riferimento}, che seguono l'ordine suggerito dalla diagnosi. La prima libera risorse dove il costo si concentra, sostituendo il primo patch embedding con un frontend piramidale che riduce la risoluzione fin dai primi strati. La seconda investe il margine così ottenuto nella memoria temporale, inserendo in entrambi i patch embedding un mixer causale che riporta nel presente le correnti di alcuni passi prima. La terza alleggerisce il primo stadio, dove una convoluzione depthwise prende il posto della Q-K Token Attention. Ogni passo modifica soltanto l'elemento indicato, a parità di ingresso e ricetta, ed è accompagnato da una tabella con prestazioni e costi misurati sugli stessi campioni.

\begin{figure}[t]
    \centering
    \includegraphics[width=\textwidth]{images/cap3/cap3_modello_riferimento.pdf}
    \caption[Il modello di riferimento]{Il modello di riferimento. (a)~Sequenza dei blocchi per un ingresso di $128\times128$ pixel, con il bordo scuro sui blocchi modificati rispetto alla baseline e l'indicazione dei patch embedding che ospitano il mixer a ritardi. (b)~Il frontend piramidale, con il sottocampionamento anticipato e i canali che raddoppiano a ogni dimezzamento della risoluzione. (c)~Il mixer causale a più ritardi, che somma alla corrente presente di uno strato le trasformazioni lineari delle correnti di uno, due, quattro e otto passi prima.}
    \label{fig:modello_riferimento}
\end{figure}

\subsection{Un frontend piramidale}
\label{subsec:frontend_piramidale}

La diagnosi indica con chiarezza da dove cominciare, perché il primo patch embedding ospita l'$80\%$ dello stato e quasi nove decimi degli accumuli potenziali, e la ragione sta quasi interamente nel suo primo strato. Una convoluzione con 32 canali applicata alla risoluzione piena dell'ingresso produce, a ogni passo, più di mezzo milione di potenziali di membrana, quasi la metà di quelli dell'intera rete. Gli strati successivi riducono la risoluzione, ma partono comunque da una mappa così ampia, e tutto il modulo ne eredita il costo.

Il frontend piramidale (figura~\ref{fig:modello_riferimento}b) ridisegna il modulo intorno a un principio semplice, dimezzare a ogni livello il numero dei neuroni. Il primo strato lavora ancora a $128\times128$ pixel, ma con 8 canali anziché 32. Gli strati successivi raddoppiano i canali, da 16 a 32 e infine a 64, e dimezzano a turno la risoluzione con un max-pooling $2\times2$ collocato dopo convoluzione e normalizzazione ma prima dei neuroni, che vivono così già sulla mappa ridotta. Ne risulta una piramide in cui ogni livello conta la metà dei potenziali del precedente, da circa 131 mila nel primo a 16 mila nell'ultimo. La shortcut resta, ma parte dalla seconda mappa e la porta alla risoluzione finale con una convoluzione $1\times1$ di passo 4. L'uscita ha le stesse dimensioni di quella originale, 64 canali su una griglia di $16\times16$ posizioni, e il resto della rete non cambia.

\begin{table}[t]
\centering
\small
\renewcommand{\arraystretch}{1.2}
\caption[Effetto del frontend piramidale]{Effetto del frontend piramidale su DVS-Lip. I costi sono per campione e provengono da profili misurati sugli stessi 64 campioni. Gli accumuli sono potenziali, lo stato è espresso in migliaia di valori persistenti e la stima è quella aritmetica della sezione~\ref{sec:caratterizzazione_computazionale}.}
\label{tab:frontend_piramidale}
\begin{adjustbox}{max width=\textwidth}
\begin{tabular}{@{}l r r r r r r@{}}
\hline
\textbf{Configurazione} & \textbf{Macro-F1} & \textbf{Parametri} & \textbf{MAC} & \textbf{Accumuli} & \textbf{Stato} & \textbf{Stima} \\
 & [\%] & & [G] & [G] & [migliaia] & [mJ] \\
\hline
Baseline & $44{,}2$ & $500\,708$ & $1{,}53$ & $5{,}56$ & $1\,100$ & $7{,}32$ \\
Frontend piramidale & $46{,}1$ & $431\,076$ & $1{,}25$ & $2{,}90$ & $477$ & $6{,}19$ \\
\hline
Variazione & $+2{,}0$ & $-13{,}9\%$ & $-18{,}5\%$ & $-47{,}7\%$ & $-56{,}6\%$ & $-15{,}4\%$ \\
\hline
\end{tabular}
\end{adjustbox}
\end{table}

Con il frontend piramidale lo stato si riduce di oltre la metà e gli accumuli potenziali quasi altrettanto, mentre moltiplicazioni e parametri scendono in misura più contenuta (tabella~\ref{tab:frontend_piramidale}). Le moltiplicazioni nascono infatti soprattutto nel secondo patch embedding, che questo passo non tocca. Le prestazioni, che ci si sarebbe potuti attendere in calo, salgono anzi leggermente. Il guadagno è piccolo, ma basta a dire che la capacità rimossa non svolgeva un lavoro indispensabile.

C'è però un rovescio, visibile nell'attività della rete, il cui firing rate medio passa dal $6\%$ al $14\%$, come se la stessa informazione dovesse attraversare un numero minore di neuroni, ciascuno dei quali si attiva più spesso. Sul costo l'effetto è trascurabile, perché gli accumuli effettivamente eseguiti restano una frazione minima della stima energetica, ma per un'esecuzione che sfrutti la sparsità è un aspetto da tenere presente (sezione~\ref{sec:efficienza_neuromorfica_vincoli}). Il frutto principale di questo passo, in ogni caso, è il margine che lascia, circa settantamila parametri e oltre seicentomila potenziali di membrana, ed è da qui che parte il passo successivo.

\subsection{Un mixer causale a più ritardi}
\label{subsec:mixer_ritardi}

Il margine liberato viene investito nella memoria, dato che la diagnosi ha mostrato che quasi tutta l'informazione temporale della baseline passa per i potenziali di membrana, e la sezione~\ref{sec:dinamica_neuronale_memoria} ha descritto i limiti di questo canale. Con i parametri adottati un neurone dimezza il proprio potenziale a ogni passo, e su DVS-Lip un ingresso arrivato otto passi prima, cioè $400\,\mathrm{ms}$ prima, pesa meno di un duecentocinquantesimo di quello appena arrivato. Una parola dura in media circa un secondo. La memoria dei neuroni copre dunque una piccola parte dell'evento da riconoscere, e per di più lo riassume in un'unica somma.

Il mixer causale a più ritardi (\emph{causal multi-lag channel mixer}, CMLCM) affianca a questo canale una via diretta verso il passato. Agisce sulla corrente $z_t$ che uno strato produce al passo $t$, prima che raggiunga i neuroni, e le somma le correnti dello stesso strato a quattro distanze precise, ciascuna trasformata da una propria matrice,
\begin{equation}
\label{eq:cmlcm}
u_t=z_t+\sum_{k=1}^{4}W_k\,z_{t-d_k},
\qquad
(d_1,d_2,d_3,d_4)=(1,2,4,8).
\end{equation}
Ogni $W_k$ combina i canali tra loro, come una convoluzione $1\times1$ priva di bias, così che un'informazione registrata su un canale possa tornare nel presente anche su altri. Nei primi passi della sequenza, quando un ritardo punta a un istante non ancora osservato, il termine corrispondente è semplicemente assente. Le matrici $W_k$ sono inizializzate a zero, così che all'avvio dell'addestramento la rete coincide con quella senza mixer, e ogni contributo del passato viene appreso, non imposto.

I ritardi crescono in progressione geometrica, e quattro soli termini coprono così otto passi ($400\,\mathrm{ms}$ per DVS-Lip) con una risoluzione fine per il passato recente e via via più grossolana per quello lontano. I ritardi sono poi fissi e uguali per tutti i canali, a differenza di quanto avviene in MD-Mixer (sezione~\ref{sec:strategie_modellazione_temporale}), una scelta che la sezione~\ref{sec:ritardi_adattivi} rimette in discussione. Il mixer, infine, è inserito in due soli punti, entrambi a risoluzione ridotta (figura~\ref{fig:modello_riferimento}a). Il primo è l'ultimo strato del frontend piramidale, che produce la mappa di $16\times16$ posizioni in uscita dal primo patch embedding. Il secondo si trova nel ramo principale del secondo patch embedding, che applica in sequenza due convoluzioni $3\times3$ (sezione~\ref{sec:architettura_baseline}), ed è la seconda di queste, quella con passo 2 che porta la mappa a $8\times8$ posizioni. In entrambi i casi il mixer agisce sulla corrente diretta ai neuroni, subito prima che li raggiunga, dopo convoluzione e normalizzazione e, nel frontend, dopo il max-pooling. Il buffer che conserva le correnti passate cresce con il numero delle posizioni, e collocarlo a risoluzione ridotta ne contiene il costo.

\begin{table}[t]
\centering
\small
\renewcommand{\arraystretch}{1.2}
\caption[Effetto del mixer a ritardi]{Effetto del mixer a ritardi su DVS-Lip, aggiunto alla rete con frontend piramidale. Il mixer opera sulle correnti e lascia invariati gli accumuli, e lo stato comprende i buffer delle correnti passate. Grandezze e unità sono quelle della tabella~\ref{tab:frontend_piramidale}.}
\label{tab:mixer_ritardi}
\begin{adjustbox}{max width=\textwidth}
\begin{tabular}{@{}l r r r r r r@{}}
\hline
\textbf{Configurazione} & \textbf{Macro-F1} & \textbf{Parametri} & \textbf{MAC} & \textbf{Accumuli} & \textbf{Stato} & \textbf{Stima} \\
 & [\%] & & [G] & [G] & [migliaia] & [mJ] \\
\hline
Frontend piramidale & $46{,}1$ & $431\,076$ & $1{,}25$ & $2{,}90$ & $477$ & $6{,}19$ \\
+ mixer a ritardi & $53{,}9$ & $512\,996$ & $1{,}58$ & $2{,}90$ & $674$ & $7{,}53$ \\
\hline
Variazione & $+7{,}7$ & $+19{,}0\%$ & $+26{,}9\%$ & $0{,}0\%$ & $+41{,}2\%$ & $+21{,}5\%$ \\
\hline
\end{tabular}
\end{adjustbox}
\end{table}

È il passo che pesa di più, con quasi otto punti di Macro-F1 (tabella~\ref{tab:mixer_ritardi}). Il prezzo è di circa ottantamila parametri per le quattro matrici, quattro quinti dei quali nel secondo patch embedding, dove i canali sono 128, e di circa duecentomila valori di stato per i buffer degli ultimi otto passi. Gli accumuli restano invariati, mentre le moltiplicazioni crescono di un quarto, perché le correnti non sono binarie e il mixer le combina con moltiplicazioni vere e proprie. Grazie al margine del passo precedente, però, i parametri superano quelli della baseline di meno del $3\%$, e lo stato ne resta poco più della metà.

L'interpretazione del guadagno è coerente con quanto discusso nel capitolo~\ref{chap:fondamenti_stato_arte}. Il mixer non confronta passato e presente, si limita a sommarli, ma rende così disponibile al neurone, nel momento in cui decide, una versione del passato che la sola membrana avrebbe in gran parte già dimenticato, e lascia alla soglia e agli strati successivi il compito di metterla in relazione con il presente. È la differenza tra persistenza e accesso della sezione~\ref{subsec:memoria_relazioni}, e il capitolo~\ref{chap:memoria_relazionale} riparte proprio da qui, chiedendosi se il confronto possa diventare esplicito.

Il guadagno, infine, non si spiega con i soli parametri aggiunti, e i controlli di capacità della sezione~\ref{sec:solidita_riferimento} mostrano che una baseline con un numero di parametri quattro volte maggiore resta al di sotto del modello di riferimento.

\subsection{Dall'attenzione Q-K alla convoluzione depthwise}
\label{subsec:depthwise}

Con il frontend alleggerito e la memoria rafforzata, resta da riconsiderare il primo stadio. Max-Former~\cite{fang2025maxformer}, uno spiking Transformer che condivide con QKFormer l'architettura gerarchica, osserva che i neuroni spiking si comportano come filtri passa-basso e tendono ad attenuare le componenti ad alta frequenza delle feature, cioè i dettagli fini e i contorni netti. Gli stadi iniziali di una rete visiva sono però proprio quelli che ne hanno più bisogno, perché costruiscono le feature locali da cui gli stadi successivi ricavano relazioni via via più globali. Per questo Max-Former sostituisce l'attenzione dei primi stadi con una convoluzione depthwise, un operatore locale che preserva quei dettagli con pochi pesi e un costo lineare nel numero dei token, e conserva la self-attention soltanto più in profondità. Nella configurazione per i dataset neuromorfici adotta esattamente una convoluzione depthwise $3\times3$ nel primo stadio e la SSA nel secondo.

Il ragionamento si adatta bene alla baseline compatta, visto che la sezione~\ref{sec:qkformer_riferimento} ha mostrato che la Q-K Token Attention seleziona i token in base all'attività della loro Query, senza mettere in relazione token diversi. Nel primo stadio i token sono 256, su una griglia di $16\times16$, e ciascuno descrive una porzione di circa $8\times8$ pixel dell'ingresso. In un dataset come DVS-Lip il segnale sta nella forma della bocca, cioè nei rapporti tra porzioni vicine dell'immagine, e un'operazione che decide token per token, senza guardare i vicini, ne coglie soltanto una parte.

L'ultimo passo sostituisce quindi l'attenzione con una convoluzione depthwise $3\times3$, seguita da normalizzazione e neurone. A differenza di una convoluzione ordinaria, che combina tutti i canali, una convoluzione depthwise filtra ciascun canale separatamente, con un nucleo di nove pesi che mette in relazione ogni posizione con le otto adiacenti. È un mixing spaziale puramente locale, complementare allo spiking MLP del blocco, che mescola i canali in ogni posizione senza guardare le vicine. Le tre proiezioni dell'attenzione costavano circa dodicimila parametri, la convoluzione depthwise ne richiede meno di settecento. Collegamento residuo, MLP, ritardi del mixer e self-attention del secondo stadio restano invariati, e la rete conserva così un operatore globale nello stadio in cui i token sono pochi.

\begin{table}[t]
\centering
\small
\renewcommand{\arraystretch}{1.2}
\caption[Effetto della convoluzione depthwise]{Effetto della convoluzione depthwise nel primo stadio su DVS-Lip, con frontend piramidale e mixer a ritardi. La seconda riga è il modello di riferimento. Grandezze e unità sono quelle della tabella~\ref{tab:frontend_piramidale}.}
\label{tab:depthwise}
\begin{adjustbox}{max width=\textwidth}
\begin{tabular}{@{}l r r r r r r@{}}
\hline
\textbf{Primo stadio} & \textbf{Macro-F1} & \textbf{Parametri} & \textbf{MAC} & \textbf{Accumuli} & \textbf{Stato} & \textbf{Stima} \\
 & [\%] & & [G] & [G] & [migliaia] & [mJ] \\
\hline
Q-K Token Attention & $53{,}9$ & $512\,996$ & $1{,}58$ & $2{,}90$ & $674$ & $7{,}53$ \\
convoluzione depthwise & $55{,}2$ & $501\,028$ & $1{,}51$ & $2{,}86$ & $639$ & $7{,}16$ \\
\hline
Variazione & $+1{,}3$ & $-2{,}3\%$ & $-4{,}9\%$ & $-1{,}4\%$ & $-5{,}2\%$ & $-4{,}8\%$ \\
\hline
\end{tabular}
\end{adjustbox}
\end{table}

La sostituzione conviene sotto ogni aspetto, perché riduce parametri, moltiplicazioni e stato e migliora la Macro-F1 di oltre un punto (tabella~\ref{tab:depthwise}). È un guadagno piccolo in assoluto, ma ogni grandezza si muove nella direzione di una rete al tempo stesso più leggera e più accurata. Nel regime compatto un'operazione locale e quasi gratuita rende più di un'attenzione che seleziona i token senza metterli in relazione. Resta aperta una domanda più ampia, se la stessa logica valga anche per la self-attention del secondo stadio, che è globale per costruzione. La sezione~\ref{sec:ssa_secondo_stadio} la affronta con due test complementari.

\begin{figure}[t]
    \centering
    \includegraphics[width=\textwidth]{images/cap3/cap3_percorso.pdf}
    \caption[Dalla baseline al modello di riferimento su DVS-Lip]{Dalla baseline al modello di riferimento su DVS-Lip. Ogni passo aggiunge una modifica a quelli precedenti. (a)~Macro-F1 sulla validazione di sviluppo, con il guadagno di ciascun passo alla base delle barre. (b)~Costi per campione dopo ciascun passo, in percentuale della baseline, con le riduzioni in blu e gli aumenti in arancio. Le MAC comprendono quelle del mixer a ritardi, gli accumuli sono potenziali e la stima energetica è quella aritmetica della sezione~\ref{sec:caratterizzazione_computazionale}. Parametri, MAC e stima energetica tornano alla fine sui valori della baseline, mentre stato e accumuli si riducono di circa la metà.}
    \label{fig:percorso_riferimento}
\end{figure}

\paragraph{Bilancio del percorso}
Nel percorso complessivo, riassunto dalla figura~\ref{fig:percorso_riferimento}, il modello di riferimento guadagna undici punti di Macro-F1 rispetto alla baseline, con 320 parametri in più, quasi la metà degli accumuli potenziali e il $42\%$ di stato in meno. Le moltiplicazioni e la stima energetica restano invece quasi invariate, e il motivo è istruttivo. Nella baseline le MAC rappresentano il $96\%$ della stima, e nascono dagli operandi interi, i conteggi in ingresso e le somme residue, che il percorso non elimina. Il risparmio si concentra così sugli accumuli e sullo stato, due voci che la stima aritmetica pesa poco ma che, come discusso nella sezione~\ref{sec:efficienza_neuromorfica_vincoli}, contano molto in una realizzazione hardware. In sintesi, il budget non è stato aumentato ma spostato, dalla risoluzione del frontend alla memoria temporale e dall'attenzione del primo stadio a un mixing locale, e lo spostamento è valso più di quanto avrebbe reso quadruplicarlo.


% ==============================================================================
%  Solidità del risultato
% ==============================================================================

\section{Solidità del risultato}
\label{sec:solidita_riferimento}

Il percorso della sezione precedente poggia su un solo dataset e su un solo addestramento per passo. Prima di adottare il modello di riferimento conviene quindi sottoporne il vantaggio a tre verifiche, ripetendo l'addestramento, cambiando dataset e confrontandolo con versioni più grandi della baseline.

Per misurare la variabilità dell'addestramento, su DVS-Lip baseline e modello di riferimento sono stati addestrati tre volte, con inizializzazioni diverse. La baseline raggiunge una Macro-F1 di $43{,}9\pm0{,}4$, il modello di riferimento di $55{,}5\pm0{,}9$, in media e deviazione standard. La variabilità dovuta all'addestramento resta sotto il punto, mentre la distanza tra i due modelli supera gli undici, e la migliore ripetizione della baseline resta oltre dieci punti sotto la peggiore del modello di riferimento. Lo stesso confronto aiuta però a leggere i singoli passi del percorso. Il frontend piramidale e soprattutto il mixer a ritardi producono guadagni superiori a questa variabilità. Il punto guadagnato con la convoluzione depthwise ne è invece dello stesso ordine, e va considerato un'indicazione, sostenuta dalla riduzione concorde di tutti i costi, più che un miglioramento accertato.

\medskip
Resta da capire se il vantaggio dipenda dal dataset su cui è stato costruito. Esteso agli altri tre dataset (figura~\ref{fig:solidita_riferimento}a), il confronto mostra che il modello di riferimento migliora la baseline su tutti e quattro, con guadagni compresi tra sei e quindici punti, e le ripetizioni disponibili non cambiano il quadro. Il guadagno è più contenuto dove la baseline è già alta, come su DVS128 Gesture, il cui limite principale, come mostrerà il capitolo~\ref{chap:clock_input}, sta nella rappresentazione dell'ingresso più che nella rete. È invece massimo su THU-EACT-50-CHL, dove la baseline parte più in basso. Anche su CIFAR10-DVS, dove ci si poteva attendere che la memoria temporale contasse meno (sezione~\ref{sec:dataset}), il vantaggio supera i nove punti. Le modifiche, scelte guardando soltanto DVS-Lip, non sono quindi legate a quel dataset.

\begin{figure}[t]
    \centering
    \includegraphics[width=\textwidth]{images/cap3/cap3_solidita.pdf}
    \caption[Solidità del modello di riferimento]{Solidità del modello di riferimento, sulla validazione di sviluppo. (a)~Baseline compatta e modello di riferimento sui quattro dataset, con il guadagno indicato sopra ogni coppia. Le barre corrispondono a un addestramento per modello, i punti alle ripetizioni con inizializzazioni diverse, disponibili su DVS-Lip per entrambi i modelli e sugli altri dataset per il solo modello di riferimento. (b)~Macro-F1 su DVS-Lip in funzione dei parametri, per reti con 128, 192 e 256 canali. La linea tratteggiata indica il modello di riferimento nella versione compatta.}
    \label{fig:solidita_riferimento}
\end{figure}

\medskip
Quanto alla capacità, la diagnosi aveva mostrato che allargare la baseline rende poco, ma non escludeva che il modello di riferimento dovesse il suo vantaggio al regime compatto in cui è stato trovato. La figura~\ref{fig:solidita_riferimento}b mette a confronto le due architetture alle stesse tre larghezze, 128, 192 e 256 canali, cioè circa mezzo, uno e due milioni di parametri. Le due curve salgono con pendenze simili e restano separate da circa dieci punti a ogni larghezza, e la baseline più grande resta oltre tre punti sotto il modello di riferimento più piccolo. Il vantaggio, dunque, viene dalla struttura della rete e si conserva quando la capacità cresce.

\begin{table}[t]
\centering
\small
\renewcommand{\arraystretch}{1.2}
\caption[Struttura e capacità a confronto su DVS-Lip]{Struttura e capacità a confronto su DVS-Lip, per la baseline compatta e il modello di riferimento con 128, 192 e 256 canali, cioè circa mezzo, uno e due milioni di parametri. Grandezze e unità sono quelle della tabella~\ref{tab:frontend_piramidale}.}
\label{tab:capacita_lip}
\begin{adjustbox}{max width=\textwidth}
\begin{tabular}{@{}l r r r r r r@{}}
\hline
\textbf{Modello} & \textbf{Macro-F1} & \textbf{Parametri} & \textbf{MAC} & \textbf{Accumuli} & \textbf{Stato} & \textbf{Stima} \\
 & [\%] & & [G] & [G] & [migliaia] & [mJ] \\
\hline
Baseline, 128 canali & $44{,}2$ & $500\,708$ & $1{,}53$ & $5{,}56$ & $1\,100$ & $7{,}32$ \\
Baseline, 192 canali & $49{,}4$ & $1\,113\,508$ & $3{,}16$ & $12{,}50$ & $1\,649$ & $15{,}13$ \\
Baseline, 256 canali & $51{,}8$ & $1\,967\,972$ & $5{,}37$ & $22{,}23$ & $2\,198$ & $25{,}68$ \\
\hline
Riferimento, 128 canali & $55{,}2$ & $501\,028$ & $1{,}51$ & $2{,}86$ & $639$ & $7{,}16$ \\
Riferimento, 192 canali & $59{,}8$ & $1\,114\,468$ & $3{,}31$ & $6{,}44$ & $958$ & $15{,}76$ \\
Riferimento, 256 canali & $61{,}6$ & $1\,969\,892$ & $5{,}82$ & $11{,}45$ & $1\,278$ & $27{,}60$ \\
\hline
\end{tabular}
\end{adjustbox}
\end{table}

La tabella~\ref{tab:capacita_lip} mostra quanto costerebbe inseguire il modello di riferimento con la sola capacità, senza raggiungerlo. Con un quarto dei parametri, il modello compatto esegue poco più di un quarto delle moltiplicazioni della baseline più grande, circa un ottavo dei suoi accumuli potenziali, e conserva meno di un terzo del suo stato. Le righe intermedie rendono il confronto ancora più netto, e con 192 canali il modello di riferimento supera di otto punti la baseline più grande, con poco più della metà dei suoi parametri, circa tre quinti delle sue moltiplicazioni e meno della metà del suo stato. A parità di larghezza, invece, appena si esce dal regime compatto il modello di riferimento esegue qualche moltiplicazione in più della baseline, il $5\%$ con 192 canali e l'$8\%$ con 256, perché le matrici del mixer crescono con il quadrato del numero di canali, mentre il risparmio del frontend piramidale cresce soltanto in proporzione. Nel regime compatto i due effetti si compensano.

\begin{table}[t]
\centering
\small
\renewcommand{\arraystretch}{1.2}
\caption[Capacità del modello di riferimento sui quattro dataset]{Macro-F1 del modello di riferimento al crescere della larghezza, sui quattro dataset. Il numero dei parametri varia leggermente da un dataset all'altro con le dimensioni della testa.}
\label{tab:capacita_dataset}
\begin{tabular}{@{}l r r r@{}}
\hline
\textbf{Dataset} & \textbf{128 canali} & \textbf{192 canali} & \textbf{256 canali} \\
 & {\footnotesize($0{,}5$ milioni)} & {\footnotesize($1{,}1$ milioni)} & {\footnotesize($2{,}0$ milioni)} \\
\hline
DVS-Lip & $55{,}2$ & $59{,}8$ & $61{,}6$ \\
DVS128 Gesture & $86{,}5$ & $83{,}6$ & $85{,}5$ \\
CIFAR10-DVS & $67{,}8$ & $68{,}4$ & $70{,}3$ \\
THU-EACT-50-CHL & $45{,}6$ & $57{,}8$ & $63{,}0$ \\
\hline
\end{tabular}
\end{table}

Ripetuta sugli altri dataset per il solo modello di riferimento (tabella~\ref{tab:capacita_dataset}), la stessa prova mostra che il bisogno di capacità cambia molto da un dataset all'altro. THU-EACT-50-CHL guadagna oltre diciassette punti passando da mezzo a due milioni di parametri, DVS-Lip sei, CIFAR10-DVS due e mezzo, mentre su DVS128 Gesture le versioni più larghe non migliorano quella compatta. Il mezzo milione di parametri è quindi un vincolo stringente per THU-EACT-50-CHL e quasi irrilevante per DVS128 Gesture. I capitoli successivi mantengono comunque la dimensione compatta, che è l'oggetto di questo lavoro, tenendo presente che su THU-EACT-50-CHL i valori assoluti risentono in misura rilevante del budget.


% ==============================================================================
%  La self-attention del secondo stadio
% ==============================================================================

\section{La self-attention del secondo stadio}
\label{sec:ssa_secondo_stadio}

Dopo la sostituzione della Q-K Token Attention, la SSA del secondo stadio è l'unico operatore della rete che mette in relazione posizioni lontane. Il secondo stadio lavora su una griglia di $8\times8$ token, ciascuno dei quali descrive una porzione di $16\times16$ pixel dell'ingresso, e viene naturale chiedersi se anche qui, come nel primo stadio, un'operazione locale possa bastare. La domanda ha due versioni, che la sezione~\ref{sec:disegno_controlli} ha distinto. La prima riguarda il modello addestrato, e chiede quanto la sua decisione dipenda dalla SSA. La seconda riguarda l'architettura, e chiede se possa farne a meno.

\paragraph{Test di necessità}
Il test di necessità disattiva la SSA sul checkpoint del modello di riferimento, soltanto in inferenza. Il blocco del secondo stadio salta il ramo dell'attenzione, così che il flusso residuo raggiunge direttamente l'MLP, e tutti i pesi restano invariati. La tabella~\ref{tab:ssa_secondo_stadio} riporta il risultato, un collasso completo su tutti e quattro i dataset. L'accuratezza scende all'$1{,}9\%$ su DVS-Lip, al $9{,}8\%$ su DVS128 Gesture, al $13{,}2\%$ su CIFAR10-DVS e all'$8{,}3\%$ su THU-EACT-50-CHL, valori prossimi a quelli di una risposta casuale, e su DVS128 Gesture e CIFAR10-DVS circa nove campioni su dieci vengono assegnati alla stessa classe. La decisione del modello addestrato passa dunque interamente per la SSA.

\paragraph{Test di sostituibilità}
Il test di sostituibilità pone la domanda opposta, sostituendo la SSA con un mixer spaziale locale e addestrando il modello da zero. Il mixer applica in sequenza una convoluzione depthwise $3\times3$, un'espansione $1\times1$ che raddoppia i canali e una proiezione $1\times1$ che li riporta al valore originale, ciascuna seguita da normalizzazione e neuroni. Ha quasi gli stessi parametri della SSA, appena $1\,152$ in più, ma nessuna relazione tra posizioni lontane. Su DVS128 Gesture il modello così ottenuto supera di un punto quello con la SSA, su CIFAR10-DVS lo eguaglia, e in entrambi i casi esegue circa l'$8\%$ di moltiplicazioni in meno. Nella SSA, infatti, le proiezioni di Query, Key e Value leggono il flusso residuo, che non è binario, mentre nel mixer locale lo legge soltanto la convoluzione depthwise, che ha pochi pesi.

Su DVS-Lip e THU-EACT-50-CHL l'esito è opposto, perché il modello con il mixer locale non ha superato il test di overfitting della sezione~\ref{subsec:verifiche_preliminari}, e secondo il protocollo non è stato addestrato in forma completa. Dopo 500 epoche sui pochi campioni del test, su DVS-Lip oscilla appena sotto la soglia richiesta senza mantenerla per cinque epoche consecutive, mentre su THU-EACT-50-CHL non riesce a classificarne correttamente più di due su tre. Nelle stesse condizioni il modello con la SSA supera il test su entrambi i dataset. L'esito non dice quale risultato avrebbe raggiunto un addestramento completo, ma indica che proprio sui due dataset in cui il segnale dipende di più da movimenti fini e articolati la rete priva di SSA fatica già ad adattarsi a pochi esempi.

\begin{table}[t]
\centering
\small
\renewcommand{\arraystretch}{1.2}
\caption[Necessità e sostituibilità della SSA nel secondo stadio]{Necessità e sostituibilità della SSA nel secondo stadio. Macro-F1 [\%] del modello di riferimento, dello stesso checkpoint con la SSA disattivata in inferenza, e del modello addestrato da zero con un mixer spaziale locale al posto della SSA. Su DVS-Lip e THU-EACT-50-CHL il modello con il mixer locale non ha superato il test di overfitting (n.d.).}
\label{tab:ssa_secondo_stadio}
\begin{tabular}{@{}l r r r@{}}
\hline
\textbf{Dataset} & \textbf{Riferimento} & \textbf{SSA disattivata} & \textbf{Mixer locale} \\
 & & {\footnotesize(stesso checkpoint)} & {\footnotesize(addestrato da zero)} \\
\hline
DVS-Lip & $55{,}2$ & $0{,}3$ & n.d. \\
DVS128 Gesture & $86{,}5$ & $3{,}6$ & $87{,}6$ \\
CIFAR10-DVS & $67{,}8$ & $7{,}1$ & $67{,}8$ \\
% Test di necessita' su THU-EACT-50-CHL eseguito il 20/11/2026 su mustafar (checkpoint E0_C0_M__20261011_092915_581860,
% artifacts/reports/necessity/thu_eact_50_chl): intatto 45,55 F1, SSA disattivata 2,73 F1 e 8,27% di accuratezza.
THU-EACT-50-CHL & $45{,}6$ & $2{,}7$ & n.d. \\
\hline
\end{tabular}
\end{table}

Letti insieme, i due test raccontano una storia coerente ma non ovvia. Il collasso sul checkpoint non dimostra che la SSA sia necessaria, ma che la rete, addestrata con la SSA, ha costruito intorno a essa la propria soluzione, al punto che l'MLP che la segue non sa più interpretare un flusso residuo privo del suo contributo. Su DVS128 Gesture e CIFAR10-DVS, addestrata senza SSA, la stessa architettura trova un'altra soluzione, altrettanto buona, con un operatore che guarda soltanto le posizioni vicine. Un intervento sul checkpoint misura quindi la dipendenza della soluzione appresa da un componente, e soltanto un'ablazione riaddestrata può dire se il componente serva. È la regola con cui la tesi legge tutti gli interventi di questo tipo, come quelli della sezione~\ref{sec:uso_memoria}.

Sul piano architetturale il quadro si divide in due, perché su DVS128 Gesture e CIFAR10-DVS, nel secondo stadio di una rete così piccola, la relazione globale tra i token non porta un vantaggio misurabile. È plausibile che vi contribuiscano la dimensione ridotta della griglia, che pochi strati locali coprono quasi per intero, e il readout, che media comunque su tutte le posizioni. Su DVS-Lip e THU-EACT-50-CHL, invece, la rete senza SSA non supera nemmeno la verifica preliminare, e la relazione globale sembra svolgere un ruolo che il mixer locale non riesce a sostituire. Il modello di riferimento conserva quindi la SSA, l'unica delle due scelte che funziona su tutti e quattro i dataset, che mantiene la rete nella famiglia degli spiking Transformer da cui la baseline deriva, rendendo diretto il confronto con i modelli discussi nel capitolo~\ref{chap:fondamenti_stato_arte}. Il mixer locale resta un'alternativa documentata, più economica e altrettanto accurata dove la relazione globale conta poco, da considerare quando il costo diventa il criterio principale.


% ==============================================================================
%  Ritardi adattivi
% ==============================================================================

\section{Ritardi adattivi}
\label{sec:ritardi_adattivi}

Nel mixer del modello di riferimento i ritardi sono fissati prima dell'addestramento, uguali per tutti i canali e per tutti i campioni. Per renderli adattivi esistono due modi naturali, lasciare all'addestramento la scelta del ritardo di ciascun canale, che resta poi fisso in inferenza, come in MD-Mixer (sezione~\ref{sec:strategie_modellazione_temporale}), oppure conservare i ritardi rendendo il peso di ciascuno variabile da un campione all'altro e da un istante all'altro.

\paragraph{Ritardi appresi}
Nella prima variante ogni ramo del mixer ha, per ciascun canale in ingresso, un proprio ritardo continuo, inizializzato al valore fisso e vincolato tra uno e otto passi. Poiché uno spostamento discreto non è derivabile, durante l'addestramento la corrente ritardata è una media dei passi interi vicini, pesata da un nucleo triangolare la cui ampiezza si riduce da quattro passi a mezzo passo, fino a selezionarne uno solo. È la stessa strategia adottata da MD-Mixer. In inferenza ogni ritardo viene arrotondato all'intero più vicino. La variante aggiunge 768 parametri e non cambia lo stato, perché il buffer conserva comunque gli ultimi otto passi.

Su DVS-Lip la variante non migliora il mixer a ritardi fissi. Raggiunge una Macro-F1 di $54{,}6$, contro $55{,}2$, e la piccola differenza rientra nella variabilità tra ripetizioni. Nel primo patch embedding, come mostrano i ritardi appresi (figura~\ref{fig:ritardi_appresi}), quasi tutti i canali conservano il ritardo iniziale. Nel secondo si spostano in misura apprezzabile soltanto i due rami intermedi, e sempre di un passo verso il passato, da due a tre passi per poco meno della metà dei canali e da quattro a cinque per circa tre quarti. Gli spostamenti si accumulano gradualmente mentre il nucleo si stringe, e si arrestano attorno alla centesima epoca, quando la sua ampiezza è ormai vicina a mezzo passo. La griglia geometrica è dunque già vicina a quella che la rete sceglierebbe da sé.

\begin{figure}[t]
    \centering
    \includegraphics[width=\textwidth]{images/cap3/cap3_ritardi_appresi.pdf}
    \caption[Ritardi appresi nel mixer su DVS-Lip]{Ritardi appresi nel mixer su DVS-Lip, nella variante in cui ogni canale sceglie il proprio ritardo. (a,\,b)~Spostamento medio del ritardo rispetto al valore iniziale, sui canali di ciascun ramo dei due patch embedding, durante l'addestramento. La linea tratteggiata corrisponde al valore iniziale, quella punteggiata allo spostamento di mezzo passo oltre il quale un canale cambia ritardo quando lo si arrotonda, e la linea verticale all'epoca del checkpoint selezionato. Il ramo da otto passi coincide con il limite superiore e non può spostarsi verso il passato. (c)~Frazione dei canali il cui ritardo, arrotondato all'intero come in inferenza, differisce da quello iniziale nel checkpoint selezionato.}
    \label{fig:ritardi_appresi}
\end{figure} Il risultato differisce da quello riportato per MD-Mixer, in cui i ritardi appresi valgono circa due punti in più di quelli fissi~\cite{shi2026multidelay}, e una spiegazione plausibile sta nella diversa struttura dei due mixer. In MD-Mixer ogni canale combina pochi istanti della propria storia con coefficienti scalari, e scegliere i ritardi è l'unico modo di raggiungere il passato utile. Nel mixer del modello di riferimento ogni ritardo porta con sé una matrice completa tra i canali, e quattro ritardi in progressione geometrica offrono già abbastanza punti di accesso perché le matrici ne ricavino ciò che serve.

\paragraph{Gate di contenuto}
La seconda variante moltiplica il contributo di ciascun ritardo per un gate che dipende dalla corrente presente,
\begin{equation}
\label{eq:cmlcm_gate}
u_t=z_t+\sum_{k=1}^{4}g_{k,t}\,W_k\,z_{t-d_k},
\qquad
g_{k,t}=4\,a_t\,\pi_{k,t}.
\end{equation}
In ogni posizione una piccola rete, formata da due convoluzioni $1\times1$, legge $z_t$ e ne ricava due grandezze. L'ampiezza $a_t$, compresa tra zero e due, stabilisce quanto passato far entrare nel complesso, e la distribuzione $\pi_{k,t}$ sui quattro ritardi stabilisce come ripartirlo. L'ultimo strato della piccola rete parte da zero, così che all'inizio ogni gate vale uno e il mixer coincide con quello a ritardi fissi.

Al di là della formula, il meccanismo si comporta come un instradatore dinamico, un \emph{router}. Il mixer a ritardi fissi apre sempre le stesse quattro vie verso il passato, con lo stesso peso per ogni campione e ogni istante. Con il gate, invece, ogni posizione guarda che cosa sta ricevendo in quel momento e decide da quali vie attingere e in che misura. Un tratto in cui l'immagine cambia poco può chiudere quasi del tutto l'accesso al passato, mentre un movimento appena iniziato può privilegiare il ritardo più breve, e uno che si sviluppa lentamente quello più lungo. La variante è applicata al solo mixer del secondo patch embedding e aggiunge circa $8\,600$ parametri, oltre a una sigmoide e a una softmax da calcolare in ogni posizione e a ogni passo. Per distinguere l'effetto del contenuto da quello di una semplice ripesatura dei ritardi, un controllo usa la stessa forma con gate costanti, appresi ma indipendenti dall'ingresso, al prezzo di cinque soli parametri.

\begin{table}[t]
\centering
\small
\renewcommand{\arraystretch}{1.2}
\caption[Gate di contenuto sui quattro dataset]{Mixer a ritardi fissi del modello di riferimento e variante con gate di contenuto. Macro-F1 [\%] sulla validazione di sviluppo e parametri aggiunti rispetto al mixer a ritardi fissi.}
\label{tab:ritardi_adattivi}
\begin{adjustbox}{max width=\textwidth}
\begin{tabular}{@{}l r r r r r@{}}
\hline
\textbf{Mixer} & \textbf{Parametri} & \textbf{DVS-Lip} & \textbf{DVS128} & \textbf{CIFAR10-} & \textbf{THU-EACT-} \\
 & \textbf{aggiunti} & & \textbf{Gesture} & \textbf{DVS} & \textbf{50-CHL} \\
\hline
Ritardi fissi (riferimento) & -- & $55{,}2$ & $86{,}5$ & $67{,}8$ & $45{,}6$ \\
Gate di contenuto & $8\,581$ & $58{,}3$ & $84{,}1$ & $67{,}8$ & $50{,}2$ \\
\hline
\end{tabular}
\end{adjustbox}
\end{table}

Su DVS-Lip i gate costanti lasciano la Macro-F1 al livello del modello di riferimento, a $55{,}5$, mentre i gate di contenuto la alzano di circa tre punti (tabella~\ref{tab:ritardi_adattivi}). Il vantaggio viene quindi dalla possibilità di decidere, posizione per posizione e istante per istante, quanto passato usare, non da una ripartizione fissa migliore tra i ritardi. Sugli altri dataset, però, il quadro cambia, con un guadagno che si ripete su THU-EACT-50-CHL, dove tuttavia le due ripetizioni del modello di riferimento differiscono già di tre punti, è nullo su CIFAR10-DVS e diventa una perdita su DVS128 Gesture. L'assenza di effetto su CIFAR10-DVS è coerente con la natura del dataset, in cui la struttura temporale riflette soprattutto il movimento imposto all'immagine e lascia poco da modulare. Resta da capire perché il gate penalizzi DVS128 Gesture, e un solo addestramento per dataset non basta a stabilirlo.

Per queste ragioni la variante condizionale non entra nel modello di riferimento. Il suo guadagno non si trasferisce in modo uniforme, è stato misurato una sola volta per dataset, e aggiunge al percorso del passato funzioni non lineari da calcolare a ogni passo, che il mixer a ritardi fissi non richiede. Il capitolo~\ref{chap:memoria_relazionale} la riprende come termine di confronto, quando la rappresentazione dell'ingresso sarà cambiata.


% ==============================================================================
%  Il modello di riferimento
% ==============================================================================

\section{Il modello di riferimento}
\label{sec:modello_riferimento}

Il modello di riferimento è dunque la baseline compatta con tre modifiche, un frontend piramidale, un mixer causale a più ritardi nei due patch embedding e una convoluzione depthwise al posto della Q-K Token Attention del primo stadio. La SSA del secondo stadio, il readout e la ricetta di addestramento restano quelli della baseline. La figura~\ref{fig:modello_riferimento} ne mostra lo schema, e la tabella~\ref{tab:riferimento_sintesi} ne raccoglie le caratteristiche su DVS-Lip accanto a quelle della baseline.

\begin{table}[t]
\centering
\small
\renewcommand{\arraystretch}{1.2}
\caption[Baseline compatta e modello di riferimento]{Baseline compatta e modello di riferimento su DVS-Lip. La Macro-F1 è la media di tre addestramenti, con la deviazione standard. I costi sono per campione, misurati sugli stessi 64 campioni della validazione di sviluppo, e le operazioni dell'attenzione sono contate nelle SOP come se ogni spike fosse presente. Secondo stadio con SSA, readout e ricetta sono comuni ai due modelli.}
\label{tab:riferimento_sintesi}
\begin{tabularx}{\textwidth}{@{}l >{\raggedright\arraybackslash}X >{\raggedright\arraybackslash}X@{}}
\hline
 & \textbf{Baseline compatta} & \textbf{Modello di riferimento} \\
\hline
Primo patch embedding & 32 canali a piena risoluzione & piramidale, da 8 a 64 canali \\
Mixer a ritardi & assente & ritardi 1, 2, 4, 8 nei due patch embedding \\
Primo stadio & Q-K Token Attention & convoluzione depthwise $3\times3$ \\
\hline
Macro-F1 [\%] & $43{,}9\pm0{,}4$ & $55{,}5\pm0{,}9$ \\
Parametri & $500\,708$ & $501\,028$ \\
Stato persistente [migliaia] & $1\,100$ & $639$ \\
MAC [G] & $1{,}53$ & $1{,}51$ \\
Accumuli potenziali [G] & $5{,}56$ & $2{,}86$ \\
SOP [G] & $0{,}30$ & $0{,}26$ \\
Firing rate [\%] & $5{,}9$ & $9{,}3$ \\
Stima energetica [mJ] & $7{,}32$ & $7{,}16$ \\
\hline
\end{tabularx}
\end{table}

Il bilancio dei costi è quello già discusso alla fine della sezione~\ref{sec:percorso_riferimento}, e la tabella aggiunge un solo elemento. Il firing rate sale dal $6\%$ al $9\%$, ma le operazioni sinaptiche effettivamente eseguite diminuiscono comunque, perché i neuroni che le innescano sono molto meno numerosi. Le verifiche di questo capitolo mostrano che il vantaggio resiste a nuovi addestramenti, si ripete su tutti e quattro i dataset e non si raggiunge quadruplicando la capacità della baseline.

Il modello lascia però aperte alcune questioni, che orientano il seguito della tesi. I ritardi del mixer, come la costante di tempo dei neuroni, sono espressi in passi, e il loro significato fisico dipende interamente da come l'ingresso viene suddiviso nel tempo. La diagnosi, inoltre, ha mostrato che la rete fa uso della coda vuota che la finestra fissa lascia dopo l'ultimo evento. Il capitolo~\ref{chap:clock_input} tiene ferma l'architettura e cambia proprio il clock dell'input. Il mixer, infine, somma passato e presente senza confrontarli, e il capitolo~\ref{chap:memoria_relazionale} si chiede se un confronto esplicito possa rendere di più.


% Capitolo 4: il clock dell'input come variabile sperimentale.
\chapter{Il clock dell'input come variabile sperimentale}
\label{chap:clock_input}

% Fonti verificate il 2026-11-04 sul repository temporal-event-spiking-research (commit 5a03fec), aggiornate il
% 2026-11-11 con il controllo a durata normalizzata (commit fe0688e, encoder src/etsr/encoders/fixed_bin_count.py).
% Risultati di addestramento: summary.json dei run C0 con encoding E0, FEC, VC e FBC
%   (dvslip: E0_C0_M__20260914_093427_434930__seed42 [+seed43/44], FEC_C0_M__20261011_121941_300598__seed42,
%    VC_C0_M__20261023_195514_528865__seed42 [32x1] e VC_C0_M__20261023_103414_923006__seed42 [8x4],
%    FBC_C0_M--vc-horizon__20261106_061637_836906__seed42 [32x1, ritardi 2,4,7,15];
%    dvsgesture: E0_C0_M__20261009_235016_816383__seed42, E0 seed43, FEC_C0_M__20261011_134033_534484__seed42,
%    VC_C0_M__20261025_135526_186863__seed42, VC seed44;
%    cifar10dvs: E0 seed42/43/44, FEC_C0_M__20261011_142919_379011__seed42, VC seed42 [32x1 e 16x2], VC seed43/44;
%    thu_eact_50_chl: E0 seed42/45, VC_C0_M__20261024_025125_082871__seed42 [4x8],
%    FBC_C0_M--vc-horizon__20261106_083202_342388__seed42 [4x8, ritardi 3,6,11,23]; FEC: overfit gate fallito,
%    FBC: overfit gate superato su entrambi i dataset).
% Diagnostica appaiata E0/VC in FP32: artifacts/reports/c0_e0_vc/<dataset>/{E0,VC}/clean*; predizioni FBC e VC
% in runs/<dataset>/<run>/validation_predictions.csv (McNemar esatto FBC-VC: DVS-Lip 420/431, p=0,73; THU 30/59, p=0,003).
% Rotazioni FEC: artifacts/reports/dvsgesture/mechanistic/dvsgesture_c0_fixed_event_count__*/record.npz.
% Figure generate da figure_scripts/cap4 (fig_*.py), immagini in images/cap4/.

Nei capitoli precedenti l'ingresso della rete è rimasto fermo, e la rappresentazione a finestra fissa, introdotta insieme alla baseline, ha fatto da sfondo a ogni modifica architetturale. Proprio per questo nessuna di quelle modifiche poteva dire qualcosa sull'ingresso stesso. Questo capitolo inverte i ruoli, congelando l'architettura e cambiando soltanto il modo in cui il flusso di eventi viene suddiviso in passi.

La questione può sembrare un dettaglio di preprocessing, ma tocca il cuore del problema poiché la rete conta il tempo proprio in passi. I neuroni perdono una frazione fissa del proprio potenziale a ogni passo, il mixer temporale guarda indietro di un numero fisso di passi, il readout media sui passi. Che cosa sia un passo, però, lo decide il clock dell'input, e con esso si decide che cosa la rete possa davvero ricordare e confrontare. Le pagine che seguono mostrano che questa scelta pesa sulle prestazioni in misura paragonabile a una modifica dell'architettura, e che i suoi effetti cambiano da un dataset all'altro in modi che la sola accuratezza finale non lascerebbe intuire.

Per tutto il capitolo la rete è il modello di riferimento della sezione~\ref{sec:modello_riferimento}, mostrato nella figura~\ref{fig:modello_riferimento}. Di questa architettura interessa qui soprattutto il fatto che i ritardi $d_k$ del mixer dell'equazione~\eqref{eq:cmlcm}, come la costante $\tau_d$ dei neuroni, sono espressi in passi, e che il loro significato fisico dipende interamente dal clock.


% ==============================================================================
%  I limiti della finestra fissa
% ==============================================================================

\section{I limiti della finestra fissa}
\label{sec:limiti_finestra_fissa}

La rappresentazione a finestra fissa ha due pregi evidenti, la semplicità e la confrontabilità. Ogni campione produce esattamente 40 passi, ogni passo corrisponde alla stessa durata fisica, e il modello può essere addestrato a batch senza alcun mascheramento. La stessa rigidità che rende comodo questo schema produce però tre effetti collaterali di natura molto diversa, che riguardano ciò che accade dopo la fine della registrazione, le registrazioni più lunghe della finestra e ciò che si perde dentro un singolo bin quando il movimento è più rapido del clock.

\subsection{Code vuote e registrazioni troncate}
\label{subsec:coda_vuota}

Ogni registrazione più breve della finestra lascia dopo l'ultimo evento una sequenza di bin vuoti, già incontrata nella sezione~\ref{sec:input_baseline}. Misurata sul modello di riferimento, questa coda occupa in mediana il $45\%$ dei passi su DVS-Lip, circa un terzo su CIFAR10-DVS e circa due terzi su DVS128 Gesture e THU-EACT-50-CHL (figura~\ref{fig:coda}a). In un campione tipico di DVS128 Gesture la rete trascorre più tempo ad aspettare che a osservare.

Aspettare, però, non significa restare inerti, perché i potenziali di membrana non si azzerano quando gli eventi finiscono, il mixer temporale continua a restituire le correnti di alcuni passi prima e i collegamenti residui propagano tutto ciò che resta. Nella figura~\ref{fig:coda}b le operazioni dense, che si ripetono a ogni passo indipendentemente dall'attività, seguono esattamente la quota dei passi vuoti, ma anche le operazioni sinaptiche innescate dagli spike restano consistenti, tra il $10\%$ e il $28\%$ del totale a seconda del dataset. Una parte non trascurabile del calcolo viene quindi spesa elaborando il silenzio.

\begin{figure}[t]
    \centering
    \includegraphics[width=\textwidth]{images/cap4/cap4_coda.pdf}
    \caption[La coda vuota della rappresentazione a finestra fissa]{La coda vuota della rappresentazione a finestra fissa per il modello di riferimento, sulla validazione di sviluppo. (a)~Frazione dei 40 passi successivi al bin che contiene l'ultimo evento, per ciascun campione, con mediana e intervallo interquartile. (b)~Quota dei passi, e quindi delle MAC dense, e quota delle operazioni sinaptiche spese dopo l'ultimo evento, sommate sull'intera validazione. (c)~Accuratezza dello stesso checkpoint quando elabora l'intera finestra e quando l'elaborazione si interrompe al bin che contiene l'ultimo evento, con la differenza in punti percentuali a destra.}
    \label{fig:coda}
\end{figure}

Resta da capire se quel calcolo serva a qualcosa, e un intervento sul checkpoint permette di rispondere almeno in parte. Lo stesso modello, addestrato con la finestra fissa, viene valutato interrompendo l'elaborazione al bin che contiene l'ultimo evento e calcolando il readout sui soli passi osservati. Se la coda fosse un riempitivo innocuo, l'accuratezza non dovrebbe cambiare, e invece su DVS-Lip e THU-EACT-50-CHL crolla di circa quindici punti, mentre su DVS128 Gesture e CIFAR10-DVS ne perde soltanto due o tre (figura~\ref{fig:coda}c).

Il modello ha dunque imparato a usare la coda, e le sue decisioni continuano a maturare dopo la fine dell'ingresso, mentre la dinamica interna si assesta e il readout accumula evidenza. Va però ricordato che si tratta di un intervento sul checkpoint, che misura una dipendenza appresa e non la necessità della coda (sezione~\ref{sec:disegno_controlli}). La diagnosi della baseline lo suggeriva già, dato che la baseline compatta, riaddestrata con il readout limitato ai passi che precedono la fine della registrazione, perdeva su DVS-Lip circa due punti (sezione~\ref{sec:diagnosi_baseline}), molto meno dei diciassette persi qui dal checkpoint. Il confronto riguarda un modello diverso, ma indica che gran parte della dipendenza dalla coda è un'abitudine appresa più che un bisogno. Una rete addestrata fin dall'inizio senza coda potrebbe quindi imparare a decidere in modo diverso, ed è ciò che verifica la rappresentazione a lunghezza variabile introdotta nella sezione~\ref{sec:tre_clock}.

Il problema opposto, quello delle registrazioni che eccedono la finestra, riguarda soltanto THU-EACT-50-CHL, dove la finestra di $8\,\mathrm{s}$ tronca due registrazioni della validazione (sezione~\ref{sec:input_baseline}). Anche nell'ipotesi più favorevole, in cui la registrazione completa bastasse a correggere entrambe le predizioni, il guadagno non supererebbe mezzo punto di accuratezza.

\subsection{Quando il movimento è più rapido del clock}
\label{subsec:aliasing}

Il terzo effetto è il più sottile, e lo si osserva con chiarezza sulle rotazioni di DVS128 Gesture. Quattro classi del dataset descrivono un braccio che ruota, il destro o il sinistro, in senso orario o antiorario, e si distinguono soltanto per il verso. Con bin da $500\,\mathrm{ms}$ e una rotazione che compie un giro in poco più di un secondo, tra il centro di un bin e quello del successivo il braccio percorre circa mezzo giro. A quel punto due posizioni consecutive sono compatibili tanto con mezzo giro in un verso quanto con mezzo giro nell'altro, e il verso diventa ambiguo. È la ragione per cui una ruota ripresa con troppi pochi fotogrammi al secondo sembra girare al contrario, un caso di aliasing temporale che compare quando il campionamento non è almeno due volte più rapido del moto osservato~\cite{shannon1949communication}.

La figura~\ref{fig:aliasing} rende il fenomeno concreto su una registrazione della validazione, una rotazione antioraria del braccio sinistro. La posizione del braccio è approssimata dal baricentro degli eventi di ciascun bin, e l'angolo attorno al centro della scena ne ricostruisce la rotazione. Con bin da $20\,\mathrm{ms}$ la traccia di riferimento compie quasi sette giri sempre nello stesso verso. Con i bin da $500\,\mathrm{ms}$ della finestra fissa, invece, la metà degli intervalli tra bin consecutivi raggiunge o supera il mezzo giro, e la rotazione ricostruita procede nel verso opposto. Le due rappresentazioni alternative della sezione~\ref{sec:tre_clock}, che usano passi più brevi, restano quasi sempre sotto la soglia. Il modello addestrato con la finestra fissa classifica proprio questo campione come una rotazione oraria dello stesso braccio.

\begin{figure}[t]
    \centering
    \includegraphics[width=\textwidth]{images/cap4/cap4_aliasing.pdf}
    \caption[Aliasing temporale su una rotazione del braccio]{Aliasing temporale su una rotazione antioraria del braccio sinistro di DVS128 Gesture, dalla validazione di sviluppo. (a)~Giri ricostruiti dall'angolo del baricentro degli eventi di ciascun passo, per una traccia di riferimento a $20\,\mathrm{ms}$ e per i tre clock. (b)~Variazione angolare della traccia di riferimento tra i centri di passi consecutivi. In alto, il numero di intervalli che raggiungono o superano il mezzo giro.}
    \label{fig:aliasing}
\end{figure}

Il caso è stato scelto a posteriori, e il baricentro è un'approssimazione grossolana della posizione. Inoltre ogni bin accumula conteggi lungo tutta la sua durata, anziché campionare una singola posizione, e l'esempio illustra il meccanismo senza dimostrare che i tensori delle due classi siano indistinguibili né stimare quanto spesso accada. Il comportamento dei modelli sull'intera validazione va però nella stessa direzione, e sugli 88 campioni di rotazione il modello addestrato con la finestra fissa commette 13 errori, e 11 di questi scambiano il verso a parità di braccio (figura~\ref{fig:rotazioni}). Con le due rappresentazioni a passo più breve le inversioni scompaiono del tutto. Su questo dataset, circa un terzo di tutti gli errori della finestra fissa è un errore di verso.

\begin{figure}[t]
    \centering
    \includegraphics[width=\textwidth]{images/cap4/cap4_rotazioni.pdf}
    \caption[Matrici di confusione sulle classi di rotazione]{Matrici di confusione ristrette alle quattro classi di rotazione di DVS128 Gesture, per il modello di riferimento addestrato con i tre clock, sulla validazione di sviluppo. Le righe indicano la classe vera, le colonne la classe predetta, con D e S per il braccio destro e sinistro e le frecce per il verso. Le celle tratteggiate corrispondono alle inversioni di verso a parità di braccio, la colonna \emph{altro} alle predizioni di classi diverse dalle rotazioni.}
    \label{fig:rotazioni}
\end{figure}


% ==============================================================================
%  Tre clock a confronto
% ==============================================================================

\section{Tre clock a confronto}
\label{sec:tre_clock}

Le tre criticità hanno origini diverse ma una radice comune, perché la finestra fissa stabilisce la durata e il numero dei passi senza guardare il campione. Le due alternative esaminate in questa sezione allentano questo vincolo in modi opposti, la prima rinunciando al tempo fisico e tagliando la registrazione secondo il numero di eventi, la seconda conservando il tempo fisico e lasciando che il numero dei passi segua la durata della registrazione. La figura~\ref{fig:tre_clock} mostra come i tre schemi suddividono una stessa registrazione, costruita apposta con fasi di attività molto diverse, e la tabella~\ref{tab:contratti_clock} ne confronta le proprietà, insieme a quelle della rappresentazione di controllo della sezione~\ref{sec:cosa_conta_clock}.

\begin{figure}[t]
    \centering
    \includegraphics[width=\textwidth]{images/cap4/cap4_tre_clock.pdf}
    \caption[Una registrazione sintetica suddivisa secondo i tre clock]{Una registrazione sintetica di $15\,\mathrm{s}$, con due fasi intense, tratti quieti e una pausa senza eventi, suddivisa secondo i tre clock con i parametri di DVS128 Gesture. In alto, la frequenza degli eventi su intervalli di $10\,\mathrm{ms}$. Sotto, i passi delle tre rappresentazioni, con i passi privi di eventi in tinta chiara. La finestra fissa usa 40 bin da $500\,\mathrm{ms}$ e lascia vuoti gli ultimi dieci, successivi alla fine della registrazione. I gruppi a eventi costanti coprono esattamente la registrazione, ma la loro durata segue l'attività, da circa $100\,\mathrm{ms}$ nelle fasi intense a oltre $3\,\mathrm{s}$ nel tratto che comprende la pausa. I bin a lunghezza variabile hanno tutti la stessa durata, coprono anch'essi l'intera registrazione e lasciano visibile la pausa come una serie di passi validi senza eventi.}
    \label{fig:tre_clock}
\end{figure}

\begin{table}[tb]
\centering
\small
\renewcommand{\arraystretch}{1.2}
\caption[Proprietà dei clock dell'input]{Proprietà dei tre clock dell'input e della rappresentazione a durata normalizzata, usata come controllo nella sezione~\ref{sec:cosa_conta_clock}.}
\label{tab:contratti_clock}
\begin{tabularx}{\textwidth}{@{}>{\raggedright\arraybackslash}p{2.75cm} >{\raggedright\arraybackslash}X >{\raggedright\arraybackslash}X >{\raggedright\arraybackslash}X >{\raggedright\arraybackslash}X@{}}
\hline
 & \textbf{Finestra fissa} & \textbf{Eventi costanti} & \textbf{Lunghezza variabile} & \textbf{Durata normalizzata} {\footnotesize(controllo)} \\
\hline
Durata di un passo & fissa per dataset & variabile, segue l'attività & fissa per dataset, calibrata sull'addestramento & fissa nel campione, segue la sua durata \\
Numero di passi & 40 & 40 & variabile, 40 in media & 40 \\
Parte di registrazione elaborata & quella entro la finestra & intera & intera & intera \\
Passi dopo l'ultimo evento & sì, la coda vuota & no & no & no \\
Significato fisico di un ritardo & costante & variabile & costante & variabile tra campioni \\
Informazione necessaria in anticipo & nessuna & numero totale di eventi & nessuna & istante dell'ultimo evento \\
\hline
\end{tabularx}
\end{table}

\subsection{Due alternative alla finestra fissa}
\label{subsec:due_alternative}

\paragraph{Frame a eventi costanti}
\label{subsec:eventi_costanti}
La rappresentazione a eventi costanti (\emph{fixed event-count frames}) divide l'intera registrazione in 40 gruppi consecutivi che contengono lo stesso numero di eventi, a meno di uno, secondo l'equazione~\eqref{eq:gruppi_eventi_bilanciati}, e in ciascun gruppo conta gli eventi ON e OFF di ogni pixel come la finestra fissa. Suddivisioni di questo tipo sono diffuse tanto nella visione a eventi, dove le fette a numero costante di eventi si affiancano a quelle a durata costante~\cite{liu2018adaptive,gallego2022event}, quanto nelle librerie per le SNN~\cite{fang2023spikingjelly}.

In questo modo ogni passo riceve la stessa quantità di attività, la coda vuota scompare e una fase molto intensa viene distribuita su più passi. In cambio, la durata fisica di un passo dipende dal campione e, all'interno del campione, dall'intensità dell'attività, sicché un ritardo di 8 passi nel mixer temporale può corrispondere, a seconda del dataset e del tratto di registrazione, a poche decine di millisecondi oppure a interi secondi. Inoltre, per dividere la registrazione in 40 parti uguali occorre conoscere in anticipo il numero totale degli eventi, e la rappresentazione è perciò applicabile soltanto a registrazioni già concluse. Negli esperimenti di questo capitolo il modello e i suoi ritardi restano identici a quelli usati con la finestra fissa.

\paragraph{Durata fissa e lunghezza variabile}
\label{subsec:lunghezza_variabile}
La rappresentazione a lunghezza variabile (\emph{fixed-duration bins with variable sequence length}) conserva invece il tempo fisico. Tutti i bin di un dataset hanno la stessa durata $\Delta$, ma la sequenza parte dal primo evento e si ferma al bin che contiene l'ultimo, e la sua lunghezza $T_i$ segue per ciascun campione $i$ la regola dell'equazione~\eqref{eq:lunghezza_clock_fisico}, con $\Delta$ al posto di $\Delta t$. Il numero dei passi varia così da un campione all'altro, mentre il loro significato fisico resta lo stesso. Non c'è coda vuota, nessun evento viene scartato, e i silenzi interni alla registrazione restano visibili come bin validi privi di eventi.

Resta da scegliere $\Delta$, e per conservare lo stesso budget medio di passi della finestra fissa la durata dei bin viene calibrata sul solo insieme di addestramento $\mathcal D_{\mathrm{tr}}$, come il più piccolo intero, in microsecondi, per cui la lunghezza media non supera 40,
\begin{equation}
\label{eq:calibrazione_delta}
\Delta=\min\Bigl\{d\in\mathbb N^{+}\ \Big|\ \frac{1}{|\mathcal D_{\mathrm{tr}}|}\sum_{i\in\mathcal D_{\mathrm{tr}}}T_i(d)\leq40\Bigr\},
\end{equation}
dove $T_i(d)$ è la lunghezza del campione $i$ con bin di durata $d$. La tabella~\ref{tab:calibrazione_vc} riporta i valori ottenuti. La stessa $\Delta$ viene applicata senza modifiche alla validazione, dove la lunghezza media resta vicina a 40 ma i singoli campioni si distribuiscono su intervalli ampi, da 16 a 74 passi su DVS128 Gesture e fino a 285 per la registrazione più lunga di THU-EACT-50-CHL.

Sequenze di lunghezza diversa richiedono, all'interno di un batch, passi di padding. Coerentemente con la distinzione discussa a proposito dell'equazione~\eqref{eq:maschera_passi_validi}, questi passi sono esclusi da ogni calcolo. Non contribuiscono alle statistiche della batch normalization, non aggiornano i potenziali di membrana né la memoria del mixer temporale, e non entrano nel readout, che diventa una media sui soli $T_i$ passi validi del campione, nella forma dell'equazione~\eqref{eq:readout_media} con $T_i$ al posto di $T$. Un bin valido senza eventi, invece, resta un'osservazione e viene elaborato come tale.

Cambiando la durata dei bin cambia anche il significato dei ritardi, che per conservare gli orizzonti fisici del mixer temporale vengono riscalati e arrotondati alla nuova griglia, $d'_k=\operatorname{round}\bigl(d_k\,\Delta t/\Delta\bigr)$, dove $\Delta t$ è la durata dei bin della finestra fissa. Su DVS-Lip, per esempio, i ritardi $(1,2,4,8)$ da $50\,\mathrm{ms}$ diventano $(2,4,7,15)$ da $27{,}4\,\mathrm{ms}$, e il più lungo passa da $400$ a $411\,\mathrm{ms}$. La costante $\tau_d$ dei neuroni resta invece pari a 2 passi.

\begin{table}[tb]
\centering
\small
\renewcommand{\arraystretch}{1.2}
\caption[Calibrazione della rappresentazione a lunghezza variabile]{Calibrazione della rappresentazione a lunghezza variabile. La durata $\Delta$ è calibrata sul solo addestramento e applicata invariata alla validazione. I ritardi del mixer temporale sono riscalati per conservare gli orizzonti fisici della finestra fissa, indicati nell'ultima colonna.}
\label{tab:calibrazione_vc}
\begin{adjustbox}{max width=\textwidth}
\begin{tabular}{@{}l c c c c c@{}}
\hline
\textbf{Dataset} & \textbf{$\Delta t$ fissa} & \textbf{$\Delta$ variabile} & \textbf{Passi in validazione} & \textbf{Ritardi} & \textbf{Orizzonti [ms]} \\
\hline
DVS-Lip & $50\,\mathrm{ms}$ & $27{,}42\,\mathrm{ms}$ & da 25 a 67 & $2,4,7,15$ & $55,110,192,411$ \\
DVS128 Gesture & $500\,\mathrm{ms}$ & $164{,}36\,\mathrm{ms}$ & da 16 a 74 & $3,6,12,24$ & $493,986,1972,3945$ \\
CIFAR10-DVS & $50\,\mathrm{ms}$ & $32{,}88\,\mathrm{ms}$ & da 37 a 44 & $2,3,6,12$ & $66,99,197,395$ \\
THU-EACT-50-CHL & $200\,\mathrm{ms}$ & $70{,}00\,\mathrm{ms}$ & da 20 a 285 & $3,6,11,23$ & $210,420,770,1610$ \\
\hline
\end{tabular}
\end{adjustbox}
\end{table}

\subsection{Risultati sui quattro dataset}
\label{subsec:risultati_clock}

La tabella~\ref{tab:risultati_clock} riporta la Macro-F1 ottenuta con i tre clock, e la figura~\ref{fig:risultati_clock} vi aggiunge le ripetizioni disponibili. Il quadro non è uniforme, ed è proprio la sua disomogeneità a essere istruttiva.\footnote{Su DVS-Lip e DVS128 Gesture le esecuzioni con la finestra fissa e con la lunghezza variabile usano frazionamenti in microbatch diversi (sezione~\ref{subsec:microbatch_bn}). Dove esistono esecuzioni che differiscono soltanto per il frazionamento, la Macro-F1 cambia di meno di due decimi di punto, molto meno degli effetti discussi nel testo.}

\begin{table}[tb]
\centering
\small
\renewcommand{\arraystretch}{1.2}
\caption[Macro-F1 con i tre clock]{Macro-F1 [\%] del modello di riferimento con i tre clock, sulla validazione di sviluppo, per l'addestramento principale di ciascuna configurazione. Le ripetizioni disponibili sono riportate nella figura~\ref{fig:risultati_clock}. Su THU-EACT-50-CHL la rappresentazione a eventi costanti non ha superato il test di overfitting (n.d.).}
\label{tab:risultati_clock}
\begin{adjustbox}{max width=\textwidth}
\begin{tabular}{@{}l r r r@{}}
\hline
\textbf{Dataset} & \textbf{Finestra fissa} & \textbf{Eventi costanti} & \textbf{Lunghezza variabile} \\
\hline
DVS-Lip & $55{,}2$ & $47{,}2$ & $56{,}9$ \\
DVS128 Gesture & $86{,}5$ & $94{,}8$ & $94{,}2$ \\
CIFAR10-DVS & $67{,}8$ & $70{,}3$ & $69{,}4$ \\
THU-EACT-50-CHL & $45{,}6$ & n.d. & $54{,}9$ \\
\hline
\end{tabular}
\end{adjustbox}
\end{table}

\begin{figure}[t]
    \centering
    \includegraphics[width=\textwidth]{images/cap4/cap4_risultati_clock.pdf}
    \caption[Macro-F1 del modello di riferimento con i tre clock]{Macro-F1 del modello di riferimento con i tre clock, sulla validazione di sviluppo. Il simbolo pieno indica l'addestramento principale, i simboli vuoti le ripetizioni con inizializzazioni diverse. Ogni pannello ha una propria scala verticale. Su THU-EACT-50-CHL la rappresentazione a eventi costanti non ha superato il test di overfitting.}
    \label{fig:risultati_clock}
\end{figure}

Il caso più netto è DVS128 Gesture, dove entrambe le alternative superano la finestra fissa di circa otto punti e le ripetizioni disponibili confermano il distacco. Lo conferma anche la media delle ultime epoche, meno esposta alla scelta del massimo, che con la finestra fissa resta oltre dieci punti più in basso. È il dataset su cui la sezione~\ref{subsec:aliasing} ha mostrato l'aliasing delle rotazioni, e una parte del guadagno si lascia ricondurre direttamente a quel fenomeno. Dei 19 campioni che la finestra fissa sbaglia in più rispetto alla lunghezza variabile, 13 appartengono alle quattro classi di rotazione.

CIFAR10-DVS si comporta come ci si poteva attendere da un dataset con dinamica imposta e durate quasi identiche. Con l'addestramento principale le due alternative guadagnano tra un punto e mezzo e due punti e mezzo. Con tre ripetizioni per configurazione, però, le medie della finestra fissa e della lunghezza variabile distano poco più di mezzo punto, meno della variabilità tra ripetizioni. Su questo dataset il clock conta poco.

DVS-Lip separa invece le due alternative, con gli eventi costanti che perdono otto punti rispetto alla finestra fissa e la lunghezza variabile che ne guadagna poco più di uno e mezzo. Questo valore supera di poco tutte e tre le ripetizioni della finestra fissa, e va letto come un piccolo vantaggio, non come un salto. Il risultato degli eventi costanti chiama in causa la natura del segnale, perché nella lettura labiale contano il ritmo e la durata relativa dei movimenti, e una suddivisione che dilata i tratti quieti e comprime quelli concitati altera proprio queste proporzioni. La sezione~\ref{sec:cosa_conta_clock} mette alla prova questa lettura.

Su THU-EACT-50-CHL, infine, la lunghezza variabile guadagna oltre nove punti di Macro-F1, e più di dieci di accuratezza. Una seconda esecuzione della finestra fissa riduce il distacco senza annullarlo. La rappresentazione a eventi costanti, invece, non ha superato su questo dataset il test di overfitting controllato descritto nella sezione~\ref{subsec:verifiche_preliminari}, e non è stata addestrata in forma completa. Dopo 500 epoche sui pochi campioni del test non riesce ancora a mantenere la soglia richiesta, che con lo stesso numero di passi la rappresentazione a durata normalizzata della sezione~\ref{sec:cosa_conta_clock} raggiunge in meno di 150.

Su DVS128 Gesture e CIFAR10-DVS la classe si riconosce soprattutto dalla forma e da un movimento ampio o ripetuto, più che dalla durata e dall'ordine delle sue fasi, e su questi due dataset la suddivisione per numero di eventi tiene il passo della lunghezza variabile, con $94{,}8$ contro $94{,}2$ sul primo e $70{,}3$ contro $69{,}4$ sul secondo, e supera la finestra fissa. Deformare il tempo all'interno della registrazione costa poco, perché ritmo e durata contano poco. Su DVS-Lip e THU-EACT-50-CHL, i due dataset più complessi e più legati alla struttura temporale, la stessa scelta crolla, perdendo otto punti rispetto alla finestra fissa sul primo e non superando nemmeno il test di overfitting sul secondo. Dove il ritmo, la durata e la successione dei movimenti fanno parte della classe, un clock che segue l'attività anziché il tempo cancella proprio l'informazione che servirebbe.

Per capire dove nasca il guadagno, la figura~\ref{fig:quartili_durata} divide la validazione in quartili di durata e confronta, quartile per quartile, i due modelli che lavorano in tempo fisico. Su DVS128 Gesture e THU-EACT-50-CHL la lunghezza variabile è migliore in tutti i quartili, e il vantaggio più ampio si concentra sui campioni più brevi, quelli che nella finestra fissa lasciano la coda più lunga. Sul quartile più breve di THU-EACT-50-CHL l'accuratezza passa dal $47\%$ al $68\%$. Su DVS-Lip il quadro è quasi rovesciato, perché i campioni più brevi vengono osservati con più passi utili, in mediana 33 contro i 18 della finestra fissa, ma perdono i passi di assestamento della coda e il loro risultato non cambia, mentre il guadagno si concentra nel quartile più lungo, dove l'accuratezza sale di circa cinque punti. Su CIFAR10-DVS, con durate quasi identiche, le differenze tra quartili rientrano nell'incertezza.

\begin{figure}[t]
    \centering
    \includegraphics[width=\textwidth]{images/cap4/cap4_quartili_durata.pdf}
    \caption[Accuratezza per quartili di durata della registrazione]{Accuratezza del modello di riferimento con la finestra fissa e con la lunghezza variabile, per quartili di durata della registrazione, sulla validazione di sviluppo. Le barre verticali indicano l'intervallo di confidenza al $95\%$ di Wilson, ampio dove i quartili contano pochi campioni, come i 66 di DVS128 Gesture.}
    \label{fig:quartili_durata}
\end{figure}

Nel confronto appaiato, campione per campione, la lunghezza variabile corregge su DVS-Lip 467 errori della finestra fissa, ma ne introduce 431 nuovi. I due modelli hanno un'accuratezza simile e sbagliano in buona parte campioni diversi, un indizio che il clock non sposta soltanto il livello delle prestazioni ma anche il tipo di informazione che la rete sfrutta. Su DVS128 Gesture e THU-EACT-50-CHL lo sbilanciamento è invece netto, con 25 correzioni contro 6 nuovi errori e 72 contro 29.


% ==============================================================================
%  Che cosa conta nel clock
% ==============================================================================

\section{Che cosa conta nel clock}
\label{sec:cosa_conta_clock}

Il passaggio dalla finestra fissa alla lunghezza variabile toglie la coda vuota e il troncamento e porta nella sequenza l'intera registrazione, ma insieme rende variabile il numero dei passi e dà a ciascun passo la stessa durata fisica in tutto il dataset. Per separare questi fattori, un controllo riaddestrato conserva la copertura dell'intera registrazione e il numero fisso di passi, e rinuncia invece a un passo di durata costante.

\label{subsec:durata_normalizzata}
La rappresentazione a durata normalizzata (\emph{duration-normalized bins}) divide l'intervallo tra il primo e l'ultimo evento di ciascuna registrazione in 40 bin di uguale durata, secondo l'equazione~\eqref{eq:delta_t_rapporto}, e corrisponde alla suddivisione per tempo che librerie come SpikingJelly offrono accanto a quella per numero di eventi~\cite{fang2023spikingjelly}. Come gli eventi costanti produce sempre 40 passi e copre l'intera registrazione senza coda, ma divide il tempo in modo uniforme all'interno del campione, senza dilatare i tratti quieti né comprimere quelli concitati. Come la lunghezza variabile segue il tempo fisico, ma la durata del passo cambia da un campione all'altro, perché ogni registrazione viene stirata o compressa fino a occupare 40 passi. È il caso discusso nella sezione~\ref{subsec:discretizzazione_clock}, in cui la stessa distanza in passi corrisponde a tempi fisici diversi senza che la rete possa saperlo. Il controllo è stato eseguito su DVS-Lip e THU-EACT-50-CHL, i due dataset su cui restavano domande aperte, con lo stesso frazionamento in microbatch e gli stessi ritardi riscalati della lunghezza variabile, che in media corrispondono agli stessi orizzonti fisici.

Su DVS-Lip la durata normalizzata raggiunge la lunghezza variabile, con una Macro-F1 di $56{,}2$ contro $56{,}9$. Nel confronto appaiato corregge 420 suoi errori e ne introduce 431, un bilancio compatibile con l'assenza di differenze ($p\approx0{,}73$ con il test di McNemar), e la figura~\ref{fig:durata_normalizzata} mostra curve quasi sovrapposte in tutti i quartili di durata. La lettura proposta per gli eventi costanti trova così una conferma, dato che su DVS-Lip non nuocciono né il numero fisso di passi né la normalizzazione della durata, che qui varia poco, ma la suddivisione per numero di eventi, che deforma il tempo all'interno della registrazione.

Su THU-EACT-50-CHL, invece, la durata normalizzata recupera soltanto una parte del vantaggio, e con una Macro-F1 di $49{,}7$ supera la finestra fissa di circa quattro punti, poco più della distanza tra le sue due ripetizioni, e resta cinque punti sotto la lunghezza variabile, con uno sbilanciamento netto nel confronto appaiato. La lunghezza variabile corregge 59 suoi errori e ne introduce 30 ($p\approx0{,}003$). Il deficit si concentra sulle registrazioni più brevi, le stesse su cui la lunghezza variabile aveva guadagnato di più rispetto alla finestra fissa. Togliere la coda e il troncamento a parità di passi, dunque, non basta a spiegare il vantaggio della lunghezza variabile su questo dataset, e una parte consistente viene dal fatto che un passo abbia la stessa durata in tutte le registrazioni.

\begin{figure}[t]
    \centering
    \includegraphics[width=\textwidth]{images/cap4/cap4_durata_normalizzata.pdf}
    \caption[Accuratezza per quartili di durata con la durata normalizzata]{Accuratezza per quartili di durata della registrazione su DVS-Lip e THU-EACT-50-CHL, con la finestra fissa, la durata normalizzata e la lunghezza variabile, sulla validazione di sviluppo. Sotto ogni quartile sono indicate la durata mediana delle registrazioni e la durata del bin che la rappresentazione a durata normalizzata le assegna. Le barre verticali indicano l'intervallo di confidenza al $95\%$ di Wilson.}
    \label{fig:durata_normalizzata}
\end{figure}

Il dato va letto con cautela, perché stirando ogni registrazione fino a 40 passi la durata normalizzata cancella insieme la scala fisica del passo e la durata assoluta dell'azione, e l'esperimento non dice quale delle due conti di più. Ogni configurazione, inoltre, è stata addestrata una sola volta. Il quadro resta comunque coerente, con un effetto che compare dove le durate variano molto, da poco più di un secondo a quasi venti su THU-EACT-50-CHL, e scompare dove variano poco, come su DVS-Lip, in cui metà dei campioni dura tra uno e 1,2 secondi circa. Su THU-EACT-50-CHL, infine, la durata normalizzata ha superato il test di overfitting che gli eventi costanti avevano fallito. Il fallimento di questi ultimi sembra quindi legato alla suddivisione per numero di eventi più che alla normalizzazione della durata, forse per la forte variabilità del numero di eventi, che su questo dataset spazia su oltre due ordini di grandezza, un'ipotesi che resta da verificare.


% ==============================================================================
%  Conseguenze per i confronti successivi
% ==============================================================================

\section{Conseguenze per i confronti successivi}
\label{sec:conseguenze_clock}

Nessuna delle alternative prevale ovunque, ma la scelta per i capitoli successivi è netta e cade sulla rappresentazione a lunghezza variabile. Il motivo principale è il clock fisico, in cui ogni passo ha la stessa durata in tutto il dataset e ritardi e costanti di tempo conservano un significato preciso, una condizione indispensabile per studiare la memoria della rete e, come mostra la sezione~\ref{sec:cosa_conta_clock}, non soltanto una questione di principio. A questo si aggiungono la copertura dell'intera registrazione senza code artificiali e la possibilità di costruire i bin man mano che gli eventi arrivano, senza conoscere in anticipo né il numero totale degli eventi, come richiedono gli eventi costanti, né l'istante dell'ultimo, come richiede la durata normalizzata. Anche sul piano dei risultati la scelta regge il confronto, e la lunghezza variabile resta vicina agli eventi costanti su DVS128 Gesture e CIFAR10-DVS, li supera nettamente su DVS-Lip e funziona su THU-EACT-50-CHL, dove gli eventi costanti non hanno superato il test di overfitting e la durata normalizzata resta cinque punti più in basso.

Da qui in avanti, il termine di confronto di ogni modifica della memoria è quindi il modello di riferimento addestrato con la rappresentazione a lunghezza variabile, con i valori della colonna a lunghezza variabile della tabella~\ref{tab:risultati_clock}. I ritardi del mixer temporale sono quelli riscalati della tabella~\ref{tab:calibrazione_vc}, mentre $\tau_d$ resta pari a 2 passi.
% Per l'appendice di riproducibilità: su ciascun dataset il riferimento condivide il frazionamento in microbatch
% con le varianti del capitolo successivo (su DVS-Lip e CIFAR10-DVS un unico microbatch da 32 campioni).

Il cambio di clock non rende la rete più economica, ma nemmeno più costosa nel modo che ci si potrebbe aspettare. Poiché $\Delta$ è calibrata per mantenere in media 40 passi, il numero medio dei passi, e con esso quello delle MAC dense, resta praticamente invariato. Crescono invece le operazioni sinaptiche, del $30\%$ su DVS128 Gesture e del $28\%$ su THU-EACT-50-CHL, perché i passi che la finestra fissa dedicava alla coda sono ora occupati da eventi. Lo stesso budget di passi viene speso osservando anziché aspettando. Per una futura implementazione hardware, però, la lunghezza della sequenza non è più limitata a 40 passi e cresce con la durata della registrazione, fino ai 285 passi del campione più lungo di THU-EACT-50-CHL, e con essa crescono latenza ed energia del singolo campione. Inoltre il buffer del mixer temporale, che deve conservare le correnti fino al ritardo massimo, si allunga con i ritardi espressi in passi, da 8 a 15 su DVS-Lip e a 24 su DVS128 Gesture, mentre lo stato dei neuroni resta invariato.

Infine, la rappresentazione a lunghezza variabile non rende la rete invariante al ritmo, perché ritardi e costanti di tempo restano espressi in passi e un gesto eseguito al doppio della velocità occupa la metà dei passi. La robustezza a questo tipo di variazione viene esaminata nella sezione~\ref{sec:ritmo}.

Il capitolo lascia così in eredità un clock in cui un passo ha lo stesso significato fisico dall'inizio alla fine della sequenza e da un campione all'altro. È la condizione che rende sensato chiedere alla memoria della rete di confrontare il presente con il passato a distanze fisiche prestabilite, ed è da qui che parte il capitolo successivo.


% Capitolo 5: memoria temporale relazionale.
\chapter{Memoria temporale relazionale}
\label{chap:memoria_relazionale}

% Riscrittura del 2026-11-19 secondo docs/STRUTTURA_NARRATIVA_E_RISULTATI_CAPITOLI_5_6_7.md (repository di ricerca,
% versione del 19/11). Scaletta approvata: 5.1 dal richiamo al confronto, 5.2 il modulo (ex 5.2 e 5.3 unite),
% 5.3 risultato a 0,5 M e che cosa usa la rete, 5.4 controlli, 5.5 due colli di bottiglia. La compressione passa al
% capitolo 6. Titoletti interni con \titoletto (mai \paragraph).
% Su DVS-Lip l'IRM antisimmetrico di riferimento e' VC_RI_M--interval-antisymmetric-r16__20261112_082513_589585__seed42
% (Macro-F1 57,47), addestrato nelle stesse condizioni del modello di riferimento e dei controlli. Nel testo non si
% menzionano altri addestramenti della stessa configurazione. Solo DVS-Lip e THU-EACT-50-CHL.
% Relazione verificata in src/etsr/models/interval.py: r = (A m)(B z) -/+ (B m)(A z), A,B in R^{16x128} condivise,
% U_k inizializzate a zero, buffer con corrente e proiezioni (C + 2r valori per passo). Stato del sito su DVS-Lip:
% 15 x 64 x 160 = 153.600 valori (600 KiB in FP32, come nei report di b_minus_final_inference).
% Bande senza confronto contro modello di riferimento (validazione, figure_scripts/cap5/risultati5.paired):
% DVS-Lip -2,0 acc. (p = 0,039), THU -0,5 acc. (p = 0,89).
% Wang et al. (arXiv 2610.37823) verificato sulla pagina dell'abstract.
% Figure: figure_scripts/cap5 (fig_richiamo_confronto, fig_bande, fig_disponibilita, fig_relazione).

Con il clock fissato nel capitolo precedente ogni passo ha una durata fisica nota, e i ritardi del mixer temporale corrispondono a distanze precise nel passato. Questo capitolo si chiede come usare meglio quella memoria, e la domanda non riguarda tanto quanto passato conservare, quanto che cosa farne quando lo si mette accanto al presente.

Il punto di partenza è il modo in cui il modello di riferimento tratta il passato, che il mixer a ritardi richiama e somma al presente senza mai metterlo a confronto con esso. Il capitolo propone un modulo che riassume il passato a distanze fisiche fisse, tiene conto di quanta storia esista davvero e confronta ciascun riassunto con il presente attraverso un termine che cambia segno quando i due istanti si scambiano. Dopo averlo descritto, ne misura l'effetto con circa mezzo milione di parametri, verifica che la rete addestrata si appoggi davvero al confronto e cerca, con una serie di controlli, da dove venga il guadagno. I controlli finiscono per raccontare qualcosa anche sui dati, e mostrano che i due compiti studiati sono frenati da limiti diversi, e questa diagnosi orienta le scelte del capitolo~\ref{chap:scalabilita}. Per tutto il capitolo l'ingresso è la rappresentazione a lunghezza variabile, e il termine di confronto è il modello di riferimento addestrato con essa (tabella~\ref{tab:risultati_clock}).

Gli esperimenti si concentrano su DVS-Lip e THU-EACT-50-CHL, perché il capitolo~\ref{chap:clock_input} ha mostrato che su DVS128 Gesture il principale problema temporale, il verso delle rotazioni, viene già risolto dal clock, e su CIFAR10-DVS la dinamica è imposta dall'acquisizione. Su questi due dataset la memoria della rete ha poco da aggiungere, e nelle prove svolte il modulo non si distingue dal riferimento. La lettura labiale e il riconoscimento di azioni sono invece compiti in cui la struttura temporale del segnale conta davvero, ed è su di essi che conviene misurare una memoria.


% ==============================================================================
%  Dal richiamo del passato al confronto tra passato e presente
% ==============================================================================

\section{Dal richiamo del passato al confronto tra passato e presente}
\label{sec:richiamo_confronto}

Nel modello di riferimento il passato entra nella corrente di uno strato attraverso il mixer causale a più ritardi dell'equazione~\eqref{eq:cmlcm}. A ogni passo il mixer rilegge le correnti di alcuni passi precedenti, le trasforma con una matrice per ciascun ritardo e le somma alla corrente presente. È, a tutti gli effetti, un'operazione di richiamo, in cui il passato viene riportato nel presente e il neurone che segue integra il risultato insieme a tutto il resto.

In questa somma ogni termine dipende da un solo istante, $z_t$ dal presente e $W_kz_{t-d_k}$ dal passato a distanza $d_k$, e il contributo di ciascuno resta lo stesso qualunque cosa contengano gli altri, sicché il mixer è additivamente separabile nel tempo. Una relazione tra passato e presente, come il fatto che una certa configurazione delle feature ne segua un'altra, non viene mai calcolata in modo esplicito. Può emergere soltanto più avanti, dalla soglia del neurone, che reagisce quando più contributi si sommano oltre il valore critico, oppure dagli strati successivi. Il mixer mette a disposizione il passato e lascia ad altri il compito di confrontarlo con il presente.

Poiché ogni ritardo ha una propria matrice, la sequenza «prima A, poi B» produce comunque una corrente diversa da «prima B, poi A», e la somma non è quindi cieca all'ordine. L'ordine, però, resta legato alla distanza dal presente, perché la stessa coppia di eventi produce correnti diverse a seconda di quanto indietro si trovi, e per riconoscere che A precede B, ovunque accada, gli strati successivi devono imparare una lettura diversa per ogni distanza.

Combinazioni lineari seguite da una soglia possono rappresentare, in linea di principio, molte relazioni, e una rete abbastanza grande è in grado di costruirle, sicché la separabilità del mixer non è di per sé un limite. In un regime compatto, però, ogni relazione che la rete deve ricavare implicitamente consuma capacità che potrebbe essere spesa altrove, e l'idea di questo capitolo è proprio rendere esplicito il confronto, con un termine che dipenda congiuntamente dal passato e dal presente.

Un confronto di questo tipo ha un precedente illustre nello studio della percezione del movimento, dove Bernhard Hassenstein e Werner Reichardt lo ricavarono negli anni Cinquanta dalla reazione optomotoria di un piccolo coleottero, il curculionide \emph{Chlorophanus viridis}, cioè dalla tendenza dell'animale a girare nel verso in cui si muove la scena che lo circonda~\cite{hassenstein1956systemtheoretische}. Il rilevatore di correlazione che porta il loro nome mette a confronto i segnali di due fotorecettori vicini, ritardando quello del primo e moltiplicandolo per quello istantaneo del secondo. Quando un contrasto si sposta dal primo recettore verso il secondo e impiega a percorrere la distanza un tempo simile al ritardo, i due segnali giungono insieme al moltiplicatore e il prodotto è grande. Lo stesso prodotto risponde però anche al moto nel verso opposto, seppure in misura minore, e a una variazione di luce che raggiunge i due recettori nello stesso istante. Per questo il rilevatore completo affianca al primo ramo il suo speculare, in cui i ruoli dei due recettori sono scambiati, e sottrae le due uscite,
\begin{equation}
\label{eq:reichardt}
R(t)=s_1(t-\delta)\,s_2(t)-s_2(t-\delta)\,s_1(t),
\end{equation}
dove $s_1$ e $s_2$ sono i segnali dei due recettori e $\delta$ il ritardo. Uno stimolo che non si sposta eccita i due rami allo stesso modo e si cancella, mentre un movimento fa prevalere uno dei due, e il segno della risposta ne indica il verso.\footnote{Negli esperimenti originali il coleottero, tenuto fermo, camminava su una sfera leggera fatta di segmenti che si biforcano a Y, che girava sotto le sue zampe, e a ogni biforcazione sceglieva se andare a destra o a sinistra. Stimolando due ommatidi vicini con sequenze controllate di variazioni di luce, Hassenstein e Reichardt osservarono che due aumenti di luminosità in successione facevano girare l'animale nel verso della sequenza, e così pure due diminuzioni, mentre un aumento seguito da una diminuzione lo faceva girare nel verso opposto. È ciò che prevede una moltiplicazione, in cui il prodotto di due termini negativi è positivo. Una conseguenza meno intuitiva del modello, ritrovata poi nel comportamento delle mosche, è che per un reticolo a strisce in movimento la risposta non dipende dalla velocità in quanto tale ma dalla frequenza con cui le strisce passano davanti ai recettori, sicché la velocità preferita cresce con la larghezza delle strisce~\cite{borst2011seeing}.} Il prodotto rende quindi la risposta dipendente dalla coincidenza di passato e presente, la sottrazione del termine speculare la rende sensibile all'ordine, ed è questa coppia di ingredienti a distinguere il rilevatore da un semplice richiamo. Nelle sue varianti elaborate lo schema è equivalente ai modelli di energia del movimento, ed è ancora oggi un riferimento per i circuiti biologici che codificano la direzione~\cite{vansanten1985elaborated,borst2011seeing}. Quanto l'ordine possa contare lo ha mostrato il capitolo precedente, dove sulle rotazioni di DVS128 Gesture perdere il verso del movimento bastava a confondere classi per il resto identiche (sezione~\ref{subsec:aliasing}).

Un lavoro recente, sviluppato in parallelo a questo, formalizza la stessa distinzione per le reti che elaborano eventi, sostenendo che nessuno strato additivo possa rappresentare in modo uniforme le interazioni ordinate tra coppie di istanti, a differenza di uno stato non commutativo, e propone a sua volta una generalizzazione apprendibile del correlatore di Hassenstein e Reichardt~\cite{wang2026learning}. Le due strade condividono l'intuizione di fondo, ma ne danno realizzazioni indipendenti e diverse.

Il modulo proposto in questo capitolo trasferisce lo stesso schema dai pixel alle feature latenti (figura~\ref{fig:richiamo_confronto}). Il ruolo dei due recettori è affidato a due proiezioni apprese delle feature, quello del segnale ritardato a un riassunto del passato, e il confronto avviene in ogni posizione spaziale tra il presente e più riassunti a distanze diverse. L'analogia motiva il calcolo senza garantirne l'interpretazione, perché nulla obbliga i canali della rete a diventare rilevatori di movimento e le proiezioni possono combinare feature dal significato molto diverso.

\begin{figure}[t]
    \centering
    \includegraphics[width=\textwidth]{images/cap5/cap5_richiamo_confronto.pdf}
    \caption[Dal richiamo al confronto]{Dal richiamo al confronto. (a)~Il mixer a ritardi del modello di riferimento proietta le correnti di alcuni passi precedenti e le somma alla corrente presente, sicché passato e presente contribuiscono in modo indipendente. (b)~Il rilevatore di correlazione di Hassenstein e Reichardt moltiplica il segnale ritardato di un recettore per quello istantaneo dell'altro e sottrae il prodotto speculare, ottenendo una risposta il cui segno dipende dal verso del moto. (c)~Il mixer proposto riassume il passato in bande, lo richiama linearmente e lo confronta con il presente attraverso due proiezioni apprese, con un ruolo analogo a quello dei due recettori. Il contributo di ogni banda è pesato dalla storia disponibile.}
    \label{fig:richiamo_confronto}
\end{figure}

Negli Spiking Transformer l'interazione esplicita tra istanti diversi è già stata affrontata, con un'attenzione che correla feature appartenenti a blocchi di passi~\cite{lee2025spiking} o con la combinazione di informazioni provenienti da più ritardi~\cite{shi2026multidelay}, come discusso nella sezione~\ref{sec:strategie_modellazione_temporale}. La strada seguita qui è più locale e più economica, dato che il confronto non passa per una matrice di attenzione ma per un termine bilineare a basso rango, calcolato in ogni posizione e collocato nello stesso punto della rete in cui oggi opera il mixer a ritardi.

Le osservazioni fatte fin qui e i capitoli che precedono fissano i requisiti del modulo, che deve anzitutto essere causale e leggere soltanto il passato già osservato. Deve poi conservare la scala fisica introdotta dalla rappresentazione a lunghezza variabile, e quindi riassumere il passato a distanze fisiche prestabilite. Deve sapere quanta storia esista davvero, perché all'inizio di ogni sequenza il passato è incompleto. Deve mantenere il percorso lineare del mixer a ritardi, in modo che il confronto si aggiunga al richiamo anziché sostituirlo. Deve infine costare poco e partire, all'inizializzazione, dalla stessa funzione del modello di riferimento, così che ogni differenza sia frutto dell'addestramento.

I requisiti si traducono nella forma generale
\begin{equation}
\label{eq:irm_generale}
u_t=z_t+\sum_{k=1}^{K}b_{k,t}\Bigl[W_k\,m_{k,t}+U_k\,\rho\bigl(m_{k,t},z_t\bigr)\Bigr],
\end{equation}
dove $z_t\in\mathbb R^{C}$ è la corrente dello strato al passo $t$ in una posizione spaziale, prima del neurone. Il vettore $m_{k,t}$ riassume il passato alla $k$-esima scala, il coefficiente $b_{k,t}$ tiene conto della storia disponibile, $W_km_{k,t}$ è il richiamo lineare, erede diretto del mixer a ritardi, e $\rho(m_{k,t},z_t)$ è una relazione tra passato e presente, riportata nello spazio delle feature dalla matrice $U_k$. Il modulo che ne risulta viene indicato nel seguito come \emph{interval-relational mixer} (IRM).

La forma tiene distinte tre domande, cioè che cosa del passato venga conservato, nel riassunto $m_{k,t}$, quanto di quel passato esista davvero, nel coefficiente $b_{k,t}$, e come il passato venga confrontato con il presente, nella relazione $\rho$. Ponendo $\rho=0$ si ottiene una memoria puramente lineare, che è a sua volta un controllo informativo, mentre cambiando la relazione a parità di tutto il resto si possono confrontare forme di interazione con proprietà diverse. La sezione successiva affronta le tre domande nell'ordine.

Per attribuire l'effetto con chiarezza, l'IRM sostituisce il mixer a ritardi in un solo punto della rete, nel secondo patch embedding, sulla corrente della convoluzione con passo 2, dopo la batch normalization e prima del neurone. Il punto è scelto perché a quel livello la rappresentazione ha 128 canali a un sedicesimo della risoluzione d'ingresso, le feature sono più astratte di quelle del primo stadio e le posizioni sono poche, 64 per un ingresso di $128\times128$ pixel, sicché il costo del confronto resta contenuto. Il primo patch embedding conserva il mixer a ritardi, mentre SSA, MLP e readout restano invariati.


% ==============================================================================
%  Il modulo
% ==============================================================================

\section{Riassumere il passato e confrontarlo con il presente}
\label{sec:memoria_intervalli}
\label{sec:relazione}

Il modulo si costruisce in due tempi, sostituendo prima i pochi ritardi letti dal mixer con bande di passi contigui, che tengono conto di quanta storia ciascuna contenga davvero, e aggiungendo poi il confronto tra questi riassunti e il presente.

\titoletto{Dai ritardi puntuali alle bande}
\label{subsec:ritardi_bande}
Il mixer a ritardi legge il passato in pochi punti, e con la rappresentazione a lunghezza variabile i ritardi $(2,4,7,15)$ di DVS-Lip campionano 4 dei 15 passi che separano il presente dal ritardo massimo, quelli di THU-EACT-50-CHL 4 su 23 (figura~\ref{fig:bande}). Ciò che accade nei passi intermedi raggiunge il mixer soltanto attraverso le altre vie della rete, in particolare la memoria dei neuroni. Con una lettura così rada, un movimento che avviene un passo prima o un passo dopo cade fuori dal ritardo e sparisce, anche se la sua distanza fisica dal presente è quasi la stessa.

La memoria a intervalli sostituisce i punti con bande contigue di ritardi. Conservando come bordi i ritardi del modello di riferimento, $d_0=0<d_1<\dots<d_K$, la $k$-esima banda raccoglie i ritardi
\begin{equation}
\label{eq:bande_ritardi}
I_k=\{\ell\in\mathbb N\mid d_{k-1}<\ell\le d_k\},
\qquad k=1,\ldots,K.
\end{equation}
Le bande coprono senza sovrapposizioni tutti i ritardi da 1 a $d_K$. Il presente ne resta escluso e continua ad arrivare dal percorso diretto, e l'orizzonte complessivo coincide con quello del mixer a ritardi. Su DVS-Lip le bande sono $[1,2]$, $[3,4]$, $[5,7]$ e $[8,15]$, da circa $27$ a $411\,\mathrm{ms}$ nel passato. Su THU-EACT-50-CHL diventano $[1,3]$, $[4,6]$, $[7,11]$ e $[12,23]$, e arrivano a $1{,}6\,\mathrm{s}$.

\begin{figure}[t]
    \centering
    \includegraphics[width=\textwidth]{images/cap5/cap5_bande.pdf}
    \caption[Ritardi e bande della memoria a intervalli]{Ritardi del mixer a ritardi e bande della memoria a intervalli, in tempo fisico, con la rappresentazione a lunghezza variabile. I cerchi indicano i quattro ritardi letti dal mixer, i trattini grigi tutti i passi del passato fino al ritardo massimo. Le bande, delimitate dagli stessi ritardi, coprono senza lacune l'intero orizzonte e si allargano allontanandosi dal presente. A destra, quanti dei passi del passato vengono letti dal mixer a ritardi.}
    \label{fig:bande}
\end{figure}

Allontanandosi dal presente le bande si allargano, e il passato recente è descritto con la risoluzione di uno o due passi, quello più remoto con medie su otto o dodici passi. È una forma rudimentale di compressione del tempo, in cui la precisione con cui un evento viene collocato diminuisce con la sua distanza, vicina nello spirito alle rappresentazioni a scala logaritmica proposte per la memoria temporale~\cite{shankar2012scale}. Un'alternativa sarebbe lasciare che la rete apprenda la posizione dei ritardi, come avviene in alcune SNN dotate di ritardi sinaptici apprendibili~\cite{hammouamri2024learning}. Qui si è preferito conservare i bordi del modello di riferimento, così che memoria a punti e memoria a bande coprano esattamente lo stesso orizzonte. I bordi sono quindi una scelta tracciabile, non un ottimo trovato dalla rete.

\titoletto{La storia disponibile}
\label{subsec:medie_disponibili}
\label{subsec:storia_disponibile}
All'interno di ciascuna banda il passato viene riassunto dalla media delle correnti. All'inizio della sequenza, però, una banda può essere vuota o soltanto in parte riempita, perché i passi corrispondenti precederebbero il primo evento. Indicando con $J_{k,t}=\{\ell\in I_k\mid \ell\le t\}$ i ritardi della banda effettivamente osservati al passo $t$, contato a partire da zero, e con $n_{k,t}=|J_{k,t}|$ il loro numero, il riassunto è
\begin{equation}
\label{eq:media_intervallo}
m_{k,t}=
\begin{cases}
\dfrac{1}{n_{k,t}}\displaystyle\sum_{\ell\in J_{k,t}}z_{t-\ell}, & n_{k,t}>0,\\[2ex]
0, & n_{k,t}=0.
\end{cases}
\end{equation}
La media divide per la storia realmente disponibile e non per la larghezza nominale della banda, così che un riassunto parziale non venga diluito da zeri che non corrispondono ad alcuna osservazione. La media attenua inoltre il rumore e rende la memoria tollerante a piccoli spostamenti temporali, perché un movimento anticipato o ritardato di un passo resta quasi sempre nella stessa banda. Il prezzo è la perdita dell'ordine all'interno della banda, che la media rende invisibile. Le bande strette vicino al presente e il percorso diretto limitano il danno dove l'ordine fine conta di più, ma non lo annullano.

Resta da dire alla rete quanto di quel passato esista. Nel mixer a ritardi il passato mancante è rappresentato da zeri nel buffer, e poiché le proiezioni non hanno bias uno zero produce una corrente nulla. La scelta è innocua dal punto di vista numerico, ma lascia indistinguibili due situazioni molto diverse, una feature che vale davvero zero e un passato che non esiste ancora.

Con le lunghezze della rappresentazione a lunghezza variabile, la banda più lontana dispone di una storia completa nel $62\%$ dei passi su DVS-Lip, ma soltanto nel $43\%$ su THU-EACT-50-CHL, dove in quasi un terzo dei passi la trova del tutto vuota (figura~\ref{fig:disponibilita}). Su questo dataset il ritardo massimo, 23 passi, supera la metà della lunghezza tipica di una sequenza, e la memoria più lontana lavora più spesso a regime ridotto che a regime pieno.

\begin{figure}[t]
    \centering
    \includegraphics[width=\textwidth]{images/cap5/cap5_disponibilita.pdf}
    \caption[Storia disponibile nella memoria a intervalli]{Storia disponibile nella memoria a intervalli. (a)~Le bande di DVS-Lip al passo $t=5$, vicino all'inizio della sequenza, e al passo $t=30$, a regime. I passi tratteggiati precedono il primo evento e non esistono. Sopra e sotto ciascuna banda, il numero di passi osservati rispetto alla sua larghezza. (b)~Frazione dei passi della validazione di sviluppo in cui ciascuna banda dispone di una storia completa e, in piccolo, frazione in cui è del tutto vuota.}
    \label{fig:disponibilita}
\end{figure}

Per rendere esplicita questa informazione, ogni banda riceve la propria frazione di storia disponibile, $q_{k,t}=n_{k,t}/|I_k|$, che modula il contributo della banda attraverso il coefficiente
\begin{equation}
\label{eq:gate_disponibilita}
b_{k,t}=\mathbf 1\bigl[n_{k,t}>0\bigr]\;2\,\sigma\bigl(\alpha_k\,q_{k,t}+\beta_k\bigr),
\end{equation}
dove $\sigma$ è la sigmoide e $\alpha_k$, $\beta_k$ sono due parametri apprendibili per banda. All'inizializzazione $\alpha_k=\beta_k=0$, e il coefficiente vale 1 per ogni banda che contenga almeno un passo. Durante l'addestramento la rete può imparare ad attenuare o ad amplificare un riassunto in funzione di quanto è completo, al costo di $2K=8$ parametri.

Il coefficiente dipende soltanto dal tempo trascorso dall'inizio della sequenza, non dalle feature, e non va confuso con il gate di contenuto della variante condizionale del mixer a ritardi discussa nel capitolo~\ref{chap:riallocazione_budget}. Rende accessibile alla rete l'età della registrazione, cioè quanti passi siano trascorsi dal primo evento, ma non la sua durata complessiva, che resta ignota fino all'ultimo passo. Con questo meccanismo la memoria tiene separate tre condizioni che sul piano numerico potrebbero confondersi, le stesse discusse a proposito dell'equazione~\eqref{eq:maschera_passi_validi}. Il passato che precede il primo evento non esiste e non entra nel conteggio. Un bin valido privo di eventi è un'osservazione di silenzio, partecipa alla media e contribuisce a $n_{k,t}$. Il padding introdotto dal batching, infine, non entra in alcun calcolo e non fa avanzare il conteggio.

\titoletto{Il confronto antisimmetrico}
Presa da sola, la memoria a intervalli offre un riassunto del passato più completo e più robusto dei punti isolati e sa quanto di quel passato esista davvero, ma lo restituisce al presente come una somma, e resta quindi un richiamo. Come mostreranno i controlli della sezione~\ref{sec:controlli_irm}, questa sola modifica non migliora il modello di riferimento su nessuno dei due dataset, e su DVS-Lip lo peggiora. Riassumere meglio il passato non basta, e il passo successivo consiste nel confrontarlo con il presente.

Si tratta di scegliere la forma della relazione $\rho$ nell'equazione~\eqref{eq:irm_generale}, che deve dipendere da passato e presente insieme, costare poco e distinguere l'ordine, rispondendo in modo diverso a «prima A, poi B» e a «prima B, poi A». La forma più semplice che soddisfa le prime due è un prodotto elemento per elemento tra proiezioni a basso rango, una costruzione già diffusa per combinare due rappresentazioni con pochi parametri~\cite{kim2017hadamard}. Due matrici $A,B\in\mathbb R^{r\times C}$, con $r=16$, proiettano su $r$ componenti tanto il riassunto del passato quanto la corrente presente, e la relazione antisimmetrica è
\begin{equation}
\label{eq:relazione_antisimmetrica}
\rho^{-}(m,z)=(Am)\odot(Bz)-(Bm)\odot(Az),
\end{equation}
dove $\odot$ indica il prodotto elemento per elemento e gli indici $k$ e $t$ sono omessi per leggibilità. Come nel correlatore dell'equazione~\eqref{eq:reichardt}, ogni componente moltiplica il passato letto da una proiezione per il presente letto dall'altra, e sottrae il prodotto speculare. Le proiezioni $A$ e $B$ prendono il posto dei due recettori, la banda quello del ritardo. Scambiando passato e presente la relazione cambia segno, $\rho^{-}(z,m)=-\rho^{-}(m,z)$, ed è questa proprietà a soddisfare la terza richiesta. Le due proiezioni sono condivise da tutte le bande, e ogni banda ha una propria matrice $U_k\in\mathbb R^{C\times r}$ che riporta le $r$ componenti nello spazio dei canali. La relazione simmetrica $\rho^{+}$, ottenuta con il segno più, resta invariata se si scambiano passato e presente, e servirà come controllo.

\titoletto{Che cosa vede la relazione}
Le due forme sembrano differire soltanto per un segno, ma rispondono a cose diverse, come si vede scrivendo il passato come il presente più una variazione, $m=z+\Delta$. Nella forma antisimmetrica il presente confrontato con sé stesso si annulla, e restano
\begin{align}
\label{eq:relazione_variazione}
\rho^{-}(m,z)&=(A\Delta)\odot(Bz)-(B\Delta)\odot(Az),\\
\rho^{+}(m,z)&=(A\Delta)\odot(Bz)+(B\Delta)\odot(Az)+2\,(Az)\odot(Bz).\nonumber
\end{align}
La forma antisimmetrica dipende soltanto dalla variazione e vale zero quando il passato coincide con il presente. Non vede ogni cambiamento, perché una variazione che le proiezioni non colgono, o che si limita a scalare il vettore, dà ancora zero, ma ciò che vede è sempre un cambiamento. La forma simmetrica contiene invece un termine che dipende dal solo presente, la co-attivazione istantanea delle due proiezioni, e risponde anche a configurazioni che non cambiano nel tempo. Una parte della sua capacità è quindi spesa su un'informazione che con il tempo non ha nulla a che fare.

La differenza si può leggere anche in modo geometrico, riunendo per ciascuna componente $j$ le due proiezioni di un vettore $x$ in un numero complesso, $c_j(x)=a_j^{\top}x+\mathrm{i}\,b_j^{\top}x$, dove $a_j$ e $b_j$ sono le righe $j$-esime di $A$ e $B$. Con questa notazione
\begingroup
\addtolength{\abovedisplayskip}{6pt}\addtolength{\abovedisplayshortskip}{6pt}%
\addtolength{\belowdisplayskip}{6pt}\addtolength{\belowdisplayshortskip}{6pt}%
\begin{equation}
\label{eq:relazione_fase}
\rho^{-}_j=\lvert c_j(m)\rvert\,\lvert c_j(z)\rvert\,\sin(\theta_z-\theta_m),
\qquad
\rho^{+}_j=\lvert c_j(m)\rvert\,\lvert c_j(z)\rvert\,\sin(\theta_z+\theta_m),
\end{equation}
\endgroup
dove $\theta_m$ e $\theta_z$ sono le fasi di $c_j(m)$ e di $c_j(z)$. Ogni componente della forma antisimmetrica misura di quanto, e in quale verso, la coppia di proiezioni sia ruotata passando dal passato al presente (figura~\ref{fig:relazione}a). La forma simmetrica dipende invece dalla somma delle fasi, cioè dalla posizione della coppia più che dal suo movimento. È la stessa struttura dei modelli di energia del moto, nei quali il verso emerge dalla differenza di fase tra due filtri in quadratura~\cite{adelson1985spatiotemporal}, e delle rappresentazioni complesse con cui si modellano relazioni asimmetriche tra entità~\cite{trouillon2016complex}.

\begin{figure}[t]
    \centering
    \includegraphics[width=\textwidth]{images/cap5/cap5_relazione.pdf}
    \caption[Relazione antisimmetrica e simmetrica]{Relazione antisimmetrica e simmetrica, su un esempio sintetico. (a)~Le due proiezioni di un vettore formano un punto del piano. La forma antisimmetrica dipende dalla rotazione tra passato e presente, quella simmetrica dalla somma delle due fasi. (b)~Un vettore ruota a velocità costante, si ferma e torna indietro, e il passato è il presente di tre passi prima. La forma antisimmetrica resta costante durante ciascuna rotazione, con il segno del verso, e si annulla durante la pausa. La forma simmetrica segue la posizione del vettore e resta diversa da zero anche quando nulla si muove.}
    \label{fig:relazione}
\end{figure}

Nell'esempio sintetico della figura~\ref{fig:relazione}b un vettore ruota a velocità costante in un verso, si ferma e poi ruota nel verso opposto. La forma antisimmetrica resta costante e positiva durante la prima rotazione, si annulla durante la pausa e diventa negativa durante la seconda, codificando così il verso del movimento e ignorando la posizione. La forma simmetrica oscilla seguendo la posizione del vettore in entrambe le rotazioni, e durante la pausa assume un valore diverso da zero anche se nulla si muove. Le proiezioni apprese dalla rete non sono costrette a descrivere rotazioni, e i loro canali possono avere significati molto diversi, ma la proprietà resta, e la forma antisimmetrica risponde soltanto al modo in cui il presente differisce dal passato.

\titoletto{Costo e punto di partenza}
\label{subsec:costo_memoria_intervalli}
Il calcolo delle medie non richiede di sommare ogni volta tutti i passi di una banda. La memoria conserva le correnti degli ultimi $d_K$ passi, esattamente come il buffer del mixer a ritardi, e un contatore della storia disponibile per ciascun campione. Durante l'addestramento le somme di ogni banda si ottengono come differenze di somme cumulative lungo il tempo, sicché il costo per passo non dipende dalla larghezza delle bande. Rispetto al mixer a ritardi restano le stesse $K$ proiezioni $C\times C$ per posizione e passo, applicate ai riassunti anziché ai singoli punti. Le medie e i coefficienti di disponibilità aggiungono per ogni canale una manciata di operazioni elementari per banda, oltre un ordine di grandezza in meno delle $KC$ moltiplicazioni e accumuli che lo stesso canale richiede alle proiezioni.

La relazione aggiunge le due proiezioni condivise, $2rC=4\,096$ parametri, e le $K$ matrici $U_k$, altri $KCr=8\,192$. Con gli otto coefficienti di disponibilità, il modello passa da $501\,028$ a $513\,324$ parametri su DVS-Lip, il $2{,}5\%$ in più, e da $494\,578$ a $506\,874$ su THU-EACT-50-CHL. Le proiezioni del presente si calcolano una sola volta per passo e si conservano nel buffer accanto alla corrente, così che anche i loro riassunti si ottengano con le stesse somme cumulative. Il buffer passa da $C$ a $C+2r$ valori per posizione e passo, un quarto in più, e su DVS-Lip, con $d_K=15$ e 64 posizioni, conserva $153\,600$ valori per campione. Le moltiplicazioni dell'intero modello crescono di circa il $2\%$. Il capitolo~\ref{chap:scalabilita} mostrerà che, dopo l'addestramento, il modulo si lascia comprimere fino a costare meno del mixer a ritardi che sostituisce (sezione~\ref{sec:compressione}).

All'inizializzazione le matrici $W_k$ e $U_k$ valgono zero e i coefficienti di disponibilità valgono 1, sicché il modulo restituisce esattamente la corrente presente. Anche il mixer a ritardi parte da matrici nulle, e i pesi condivisi con il resto della rete vengono generati nello stesso modo, mentre $A$ e $B$ provengono da un flusso pseudo-casuale separato. La rete con l'IRM parte quindi esattamente dalla stessa funzione del modello di riferimento, e ogni differenza tra i due nasce dall'addestramento. Nel resto del capitolo l'IRM usa la relazione antisimmetrica. La memoria senza relazione e quella con la relazione simmetrica tornano nella sezione~\ref{sec:controlli_irm}, tra i controlli.




% ==============================================================================
%  Il risultato a mezzo milione di parametri
% ==============================================================================
% Fonti (2026-11-19). Validazione: metriche salvate durante l'addestramento (figure_scripts/cap5/risultati5.py).
% Test ufficiale: artifacts/reports/b_minus_final_inference/{lip,thu}_{c0,b}.npz, valutato a posteriori in sola
% inferenza; confronti appaiati calcolati in locale dai logit (DVS-Lip val 434/423 p=0,73, test 691/651 p=0,29;
% THU val 54/22 p=0,0003, test 80/26 p<0,001). IC bootstrap al 95% della differenza di accuratezza da
% docs/STRUTTURA_NARRATIVA_E_RISULTATI_CAPITOLI_5_6_7.md, tabella A (report b_minus_final_inference).
% Media delle ultime 16 epoche (Macro-F1): DVS-Lip 56,1 -> 56,7, THU 52,5 -> 59,8.
% Interventi sul checkpoint: DVS-Lip artifacts/reports/lip_ref_bands (checkpoint di riferimento, 19/11/2026:
% clean 57,13/57,47, relazione spenta 17,73/15,31, bande 55,36/52,42/42,74/33,19 di accuratezza); THU
% artifacts/reports/rt_fbc_comparison/ablations.json (clean 69,17/62,56, relazione spenta 19,05/10,72).
% Ordine dei passi sul test: E2.csv per il riferimento e per THU, artifacts/reports/lip_ref_diagnostics per l'IRM su
% DVS-Lip (19/11/2026: originale 50,85 nella diagnostica e 50,81 nella tabella, invertito 2,95, rimescolato 4,46/5,11/4,96,
% media 4,84). Figura: figure_scripts/cap5/fig_uso_memoria.py.

\section{Il risultato a mezzo milione di parametri}
\label{sec:risultati_irm}

Il primo confronto, il più semplice, mette di fronte il modello di riferimento e la rete con l'IRM antisimmetrico, che ricevono lo stesso ingresso, seguono la stessa ricetta di addestramento e hanno quasi lo stesso numero di parametri, circa mezzo milione. Differiscono soltanto per il modulo del secondo patch embedding. Entrambi sono stati scelti sulla validazione di sviluppo, che resta la misura principale di questo capitolo. Il test ufficiale è stato valutato soltanto dopo, in sola inferenza, sui modelli già selezionati (sezione~\ref{subsec:partizioni_riservate}), e nessuna scelta dipende da esso. Offre però una verifica preziosa, perché su THU-EACT-50-CHL contiene registrazioni di persone diverse da quelle usate nello sviluppo. La tabella~\ref{tab:risultato_principale} riporta le due partizioni insieme.

\begin{table}[t]
\centering
\small
\renewcommand{\arraystretch}{1.2}
\caption[Modello di riferimento e IRM antisimmetrico, validazione e test]{Modello di riferimento e IRM antisimmetrico con circa mezzo milione di parametri, sulla validazione di sviluppo e sul test ufficiale, valutato a posteriori. Valori in percentuale, differenze in punti. L'asterisco indica una differenza di accuratezza significativa nel confronto appaiato, l'intervallo è quello di confidenza al 95\% della differenza di accuratezza.}
\label{tab:risultato_principale}
\begin{tabular}{@{}l r r r r@{}}
\hline
 & \multicolumn{2}{c}{\textbf{Validazione}} & \multicolumn{2}{c}{\textbf{Test ufficiale}} \\
\textbf{Modello} & \textbf{Accuratezza} & \textbf{Macro-F1} & \textbf{Accuratezza} & \textbf{Macro-F1} \\
\hline
\multicolumn{5}{@{}l}{\emph{DVS-Lip}} \\
Modello di riferimento & $56{,}8$ & $56{,}9$ & $50{,}0$ & $49{,}7$ \\
IRM antisimmetrico & $57{,}1$ & $57{,}5$ & $50{,}8$ & $50{,}4$ \\
Differenza & $+0{,}4$ & $+0{,}6$ & $+0{,}8$ & $+0{,}7$ \\
\quad intervallo al 95\% & \footnotesize$[-1{,}5;\,+2{,}3]$ & & \footnotesize$[-0{,}6;\,+2{,}3]$ & \\
\hline
\multicolumn{5}{@{}l}{\emph{THU-EACT-50-CHL}} \\
Modello di riferimento & $61{,}2$ & $54{,}9$ & $56{,}1$ & $48{,}7$ \\
IRM antisimmetrico & $69{,}2$ & $62{,}6$ & $66{,}0$ & $59{,}6$ \\
Differenza & $+8{,}0^{*}$ & $+7{,}6$ & $+9{,}9^{*}$ & $+10{,}8$ \\
\quad intervallo al 95\% & \footnotesize$[+3{,}8;\,+12{,}3]$ & & \footnotesize$[+6{,}2;\,+13{,}4]$ & \\
\hline
\end{tabular}
\end{table}

Su THU-EACT-50-CHL l'IRM supera il riferimento in validazione di 7,6 punti di Macro-F1 e di 8 di accuratezza, un guadagno netto che la media delle ultime sedici epoche conferma con un distacco simile, segno che il risultato non dipende dalla scelta del checkpoint migliore. Sul test ufficiale il vantaggio non si riduce, anzi cresce, fino a quasi dieci punti di accuratezza e quasi undici di Macro-F1. L'IRM corregge 80 errori del riferimento e ne introduce soltanto 26. È il risultato centrale del capitolo, e il resto della sezione e quella successiva cercano di capire da dove venga.

Su DVS-Lip l'IRM resta davanti, ma di poco più di mezzo punto in validazione e di meno di un punto sul test, ed entrambi gli intervalli comprendono lo zero. In validazione i due modelli danno esiti diversi su 857 clip, e su 434 di esse l'IRM corregge un errore del riferimento mentre su 423 ne introduce uno nuovo. Sul test il rapporto è di 691 contro 651. Due reti di accuratezza quasi identica sbagliano dunque clip in buona parte diverse, e il saldo di poche decine di campioni si perde in questo rimescolamento, che spiega perché una differenza così piccola non si distingua dal caso. Con mezzo milione di parametri, insomma, la memoria relazionale non sposta il risultato di DVS-Lip in modo misurabile. Tale risultato rappresenta il punto di partenza della sezione~\ref{sec:colli_bottiglia}, che ne cerca la ragione, e del capitolo~\ref{chap:scalabilita}, che mostra quali leve funzionino invece su questo compito.

\titoletto{Che cosa usa la rete}
\label{sec:uso_memoria}
Un guadagno in tabella non dice ancora come la rete addestrata usi il modulo. Per saperlo si interviene sul checkpoint, come già nel capitolo~\ref{chap:riallocazione_budget} (sezione~\ref{sec:disegno_controlli}). Il modello viene modificato soltanto in inferenza, spegnendo una parte del modulo o alterando l'ingresso, e si misura quanto cambi il risultato. Un calo indica quanto la soluzione appresa dipenda da quella parte, non che cosa otterrebbe una rete addestrata senza di essa.

Il primo intervento azzera tutte le matrici $U_k$, cioè il contributo della relazione, e lascia intatto il richiamo lineare delle bande. La Macro-F1 crolla, di 42 punti su DVS-Lip e di 52 su THU-EACT-50-CHL. La rete non tratta il confronto come un termine accessorio, ma vi si appoggia in modo massiccio, e accade anche su DVS-Lip, dove pure il guadagno sul riferimento è piccolo. Su questo dataset il confronto è dunque una delle strade con cui la rete arriva a un risultato che il riferimento raggiunge per altre vie.

\begin{figure}[t]
    \centering
    \includegraphics[width=\textwidth]{images/cap5/cap5_uso_memoria.pdf}
    \caption[Che cosa usa l'IRM antisimmetrico addestrato]{Che cosa usa l'IRM antisimmetrico addestrato, con interventi sul checkpoint in sola inferenza. (a)~Variazione di accuratezza sulla validazione di sviluppo quando si azzerano le matrici $W_k$ e $U_k$ di una sola banda. Le bande sono ordinate dalla più vicina alla più lontana dal presente. (b)~Accuratezza sul test ufficiale del modello di riferimento e dell'IRM con i passi di ogni clip nell'ordine originale, invertiti oppure rimescolati a caso, con la media di tre permutazioni. Il contenuto di ciascun passo non cambia.}
    \label{fig:uso_memoria}
\end{figure}

Il secondo intervento spegne una banda alla volta, azzerandone le matrici $W_k$ e $U_k$ (figura~\ref{fig:uso_memoria}a). Su entrambi i dataset le bande lontane pesano molto più di quelle vicine, e senza la terza o la quarta banda l'accuratezza perde da 14 a 24 punti e senza una delle prime due da 2 a 7. La rete si appoggia soprattutto a un passato separato dal presente da qualche passo, fra $0{,}14$ e $0{,}41\,\mathrm{s}$ su DVS-Lip e fra $0{,}49$ e $1{,}61\,\mathrm{s}$ su THU-EACT-50-CHL. Una ragione plausibile del peso minore delle bande vicine è che il passato più recente somiglia ancora al presente, e il confronto con esso porta poca informazione nuova.

L'ultimo intervento altera l'ordine dei passi di ogni clip del test, invertendolo oppure rimescolandolo a caso, senza toccare il contenuto di ciascun passo. Dentro un passo gli eventi non vengono riordinati e le polarità non vengono scambiate, sicché la sequenza invertita non equivale a una scena riprodotta all'indietro, ma è piuttosto una sequenza mai vista, fatta degli stessi passi in un altro ordine. Su DVS-Lip entrambi i modelli crollano a pochi punti di accuratezza (figura~\ref{fig:uso_memoria}b), e la lettura labiale dipende quindi interamente dall'ordine, con o senza IRM. Su THU-EACT-50-CHL entrambi perdono molto, ma il vantaggio dell'IRM scompare, tanto che con i passi invertiti i due modelli si equivalgono e con i passi rimescolati l'IRM finisce anzi sotto il riferimento. I quasi dieci punti che l'IRM guadagna sulla sequenza originale dipendono quindi dall'ordine in cui i passi si susseguono, come ci si aspetta da un modulo che confronta il presente con un passato collocato a distanze precise. La prova non dimostra che l'IRM sia l'unico modo di sfruttare l'ordine, perché anche il riferimento ne dipende in larga misura, e le sequenze alterate sono comunque ingressi diversi da quelli visti in addestramento.

La soluzione appresa, in sintesi, si appoggia al confronto, lo applica soprattutto alle bande più lontane e trae il proprio vantaggio dall'ordine dei passi. Resta da capire da dove venga il guadagno, se dal confronto o semplicemente dai parametri e dal passato in più che il modulo porta con sé, e perché lo stesso modulo che su THU-EACT-50-CHL cambia il risultato su DVS-Lip quasi non lo sposti. Le due sezioni che seguono affrontano queste domande nell'ordine.


% ==============================================================================
%  Da dove viene il guadagno
% ==============================================================================
% Fonti (2026-11-19). Macro-F1, differenze, intervalli bootstrap appaiati e p di McNemar da
% figure_scripts/cap5/fig_controlli.py (controlli5.json), che legge le predizioni salvate durante l'addestramento;
% coincidono con artifacts/reports/b_minus_final_inference/validation_controls_vs_bminus.txt. Su DVS-Lip tutti i controlli
% sono appaiati con l'IRM di riferimento (57,47). Meccanismi verificati nel codice (commit dcc054b): tracce in
% src/etsr/models/adaptive_memory.py (lettura prima della scrittura, divisione per la massa, orizzonte iniziale dal
% ritardo medio della banda), orologi in src/etsr/models/change_time.py (inc = s((1+a/a_unit)/2)^gamma da s=1, gamma=0;
% a_unit = mediana pesata dagli eventi sul training), memoria remota in src/etsr/models/remote_context.py. Parametri
% appresi dell'orologio da history.csv: DVS-Lip s=0,491 gamma=0,396; THU s=0,819 gamma=0,203. Orizzonti appresi delle
% tracce dai checkpoint: DVS-Lip 0,13/0,50/1,36/4,98 s, THU 0,17/0,54/1,42/3,92 s. Tracce guidate dal sensore con il
% termine degli eventi annullato in inferenza: THU 60,2 -> 50,8 (artifacts/reports/temporal_memory).
% Variante condizionale (gate di contenuto, nota): THU 58,9 contro 54,9 del riferimento e 62,6 dell'IRM (p=0,064).
% Durata normalizzata: guadagno dell'IRM THU +1,0 F1 (49,7 -> 50,8) contro +7,6 con la lunghezza variabile.

\section{Da dove viene il guadagno}
\label{sec:controlli_irm}

Il guadagno di quasi otto punti su THU-EACT-50-CHL potrebbe avere spiegazioni meno interessanti del confronto. L'IRM aggiunge circa dodicimila parametri, e una rete un po' più grande potrebbe migliorare soltanto per questo. Porta con sé, inoltre, un passato organizzato in modo diverso da quello del mixer a ritardi, in bande anziché in punti, e il merito potrebbe essere di questa organizzazione più che del confronto. I controlli di questa sezione chiedono quindi, nell'ordine, se il guadagno venga dai parametri aggiunti, se venga dal confronto e dal suo verso, e se dipenda dal modo in cui il passato è organizzato.

Ogni controllo cambia un solo elemento dell'IRM antisimmetrico, oppure lo sostituisce con un'alternativa ragionevole, e lascia fermo tutto il resto, compreso il punto della rete in cui opera. Tutti sono addestrati da zero con la stessa ricetta. La tabella~\ref{tab:controlli} ne riporta la Macro-F1 sulla validazione di sviluppo, la figura~\ref{fig:controlli} la differenza rispetto all'IRM con il suo intervallo di confidenza. Ciascuna alternativa è stata provata in una sola versione, e le conclusioni riguardano quella versione, non ogni memoria dello stesso tipo.

\begin{table}[t]
\centering
\small
\renewcommand{\arraystretch}{1.2}
\caption[Controlli dell'IRM antisimmetrico]{Controlli dell'IRM antisimmetrico sulla validazione di sviluppo, con circa mezzo milione di parametri. Macro-F1 in percentuale e differenza rispetto all'IRM in punti. L'asterisco indica una differenza di accuratezza significativa nel confronto appaiato (test esatto di McNemar, $p<0{,}05$).}
\label{tab:controlli}
\begin{tabular}{@{}l r r r r@{}}
\hline
 & \multicolumn{2}{c}{\textbf{DVS-Lip}} & \multicolumn{2}{c}{\textbf{THU-EACT-50-CHL}} \\
\textbf{Variante} & \textbf{Macro-F1} & \textbf{Diff.} & \textbf{Macro-F1} & \textbf{Diff.} \\
\hline
IRM antisimmetrico & $57{,}5$ & & $62{,}6$ & \\
\hline
\multicolumn{5}{@{}l}{\emph{Parametri in più}} \\
\quad modello di riferimento & $56{,}9$ & $-0{,}6$ & $54{,}9$ & $-7{,}6^{*}$ \\
\quad sito vuoto & $56{,}2$ & $-1{,}3$ & $50{,}9$ & $-11{,}7^{*}$ \\
\quad MLP di pari parametri & $56{,}7$ & $-0{,}7$ & $52{,}3$ & $-10{,}3^{*}$ \\
\multicolumn{5}{@{}l}{\emph{Il confronto e il suo verso}} \\
\quad bande senza confronto & $54{,}7$ & $-2{,}8^{*}$ & $54{,}4$ & $-8{,}1^{*}$ \\
\quad confronto simmetrico & $57{,}0$ & $-0{,}5$ & $59{,}0$ & $-3{,}5$ \\
\multicolumn{5}{@{}l}{\emph{L'organizzazione del passato}} \\
\quad ritardi puntuali & $57{,}2$ & $-0{,}2$ & $59{,}7$ & $-2{,}9^{*}$ \\
\quad tracce a decadimento appreso & $54{,}7$ & $-2{,}7^{*}$ & $53{,}4$ & $-9{,}2^{*}$ \\
\quad tracce guidate dal sensore & $56{,}7$ & $-0{,}8$ & $60{,}2$ & $-2{,}4$ \\
\quad orologio appreso & $58{,}3$ & $+0{,}9$ & $58{,}6$ & $-4{,}0^{*}$ \\
\quad passato remoto aggiunto & $58{,}1$ & $+0{,}6$ & $59{,}4$ & $-3{,}2$ \\
\hline
\end{tabular}
\end{table}

\begin{figure}[t]
    \centering
    \includegraphics[width=\textwidth]{images/cap5/cap5_controlli.pdf}
    \caption[Differenza di Macro-F1 dei controlli rispetto all'IRM]{Differenza di Macro-F1 di ciascun controllo rispetto all'IRM antisimmetrico, sulla validazione di sviluppo, con l'intervallo di confidenza bootstrap al 95\%. Valori negativi indicano un controllo peggiore dell'IRM. Il simbolo pieno indica una differenza di accuratezza significativa nel confronto appaiato (test esatto di McNemar, $p<0{,}05$).}
    \label{fig:controlli}
\end{figure}

\titoletto{Parametri in più?}
Il primo controllo svuota il sito, eliminando la memoria nel secondo patch embedding. La rete non resta per questo senza memoria temporale, perché il mixer del primo patch embedding e i potenziali di membrana dei neuroni continuano a portare il passato in avanti. Il secondo controllo colloca nello stesso punto un percettrone multistrato con gli stessi parametri dell'intero IRM, che però elabora ogni passo per conto proprio, senza accesso al passato. Insieme al modello di riferimento, i due controlli delimitano ciò che si ottiene con più parametri ma senza confronto.

Su THU-EACT-50-CHL entrambi restano da dieci a dodici punti di Macro-F1 sotto l'IRM, e anche sotto il modello di riferimento. Parametri in più, senza accesso al passato, non bastano, e il guadagno nasce dal modo in cui il passato viene usato. Il sito vuoto dice qualcosa anche sul modello di riferimento, al quale togliere il mixer a ritardi dal secondo patch embedding costa circa quattro punti, segno che già il semplice richiamo era utile e che il confronto ne aggiunge altri otto. Su DVS-Lip, invece, i tre modelli restano entro poco più di un punto dall'IRM, e nessuna differenza è significativa.

\titoletto{Il confronto, e il suo verso}
Il secondo gruppo interviene sulla relazione con due varianti, le bande senza confronto e il confronto simmetrico. Le prime conservano bande, medie e coefficienti di disponibilità e tolgono soltanto il termine bilineare, ponendo $\rho=0$ nell'equazione~\eqref{eq:irm_generale}, e coincidono quindi con la memoria a intervalli della sezione~\ref{sec:memoria_intervalli}, un richiamo più completo di quello del mixer. Il confronto simmetrico usa invece la relazione $\rho^{+}$, con la stessa forma e lo stesso rango di quella antisimmetrica ma con il segno più. Come mostra l'equazione~\eqref{eq:relazione_variazione}, risponde anche al solo presente e non distingue il verso del cambiamento.

Su THU-EACT-50-CHL le bande senza confronto restano al livello del modello di riferimento, otto punti sotto l'IRM, e su DVS-Lip perdono quasi tre punti, in modo significativo. Senza confronto il guadagno scompare, e riassumere il passato meglio del mixer, da solo, non serve. Il confronto simmetrico, con gli stessi parametri, recupera su THU-EACT-50-CHL circa metà del guadagno e resta tre punti e mezzo sotto quello antisimmetrico. Con 399 campioni il confronto appaiato non basta a rendere significativo questo margine, che va quindi letto come un'indicazione. La direzione è però quella prevista dall'algebra della sezione~\ref{sec:relazione}, secondo cui la forma simmetrica spende una parte della propria capacità sulla co-attivazione istantanea delle proiezioni, quella antisimmetrica soltanto sul cambiamento. Su DVS-Lip le due forme si equivalgono.\footnote{Il capitolo~\ref{chap:riallocazione_budget} aveva messo da parte la variante condizionale del mixer a ritardi (sezione~\ref{sec:ritardi_adattivi}), che decide in base alla corrente presente quanto passato far entrare, senza confrontarlo con il presente. Ripresa con la rappresentazione a lunghezza variabile, migliora il modello di riferimento su THU-EACT-50-CHL di circa quattro punti di Macro-F1, ma resta a sua volta sotto l'IRM antisimmetrico, con un margine vicino alla significatività. Su DVS-Lip coincide con l'IRM (appendice~\ref{app:tabelle_estese}). Decidere quanto passato usare aiuta, confrontarlo con il presente sembra aiutare di più.}

\titoletto{L'organizzazione del passato}
Il terzo gruppo mantiene il confronto antisimmetrico e cambia il passato a cui si applica, con cinque alternative che conviene descrivere una alla volta prima di guardare i risultati.

I \emph{ritardi puntuali} applicano lo stesso confronto ai quattro istanti letti dal mixer a ritardi, anziché alle medie delle bande. Misurano quanto valga riassumere il passato, rispetto a campionarlo.

Le \emph{tracce} sostituiscono ciascuna banda con una media mobile esponenziale delle correnti passate, la forma di memoria tipica delle reti ricorrenti e dei modelli a spazio di stato~\cite{gu2023mamba}. A ogni passo il modulo legge la traccia, che riassume il passato fino al passo precedente, la usa al posto della media di banda $m_{k,t}$ e soltanto dopo la aggiorna con il presente,
\begin{equation}
\label{eq:traccia}
m_{k,t+1}=\lambda_{k,t}\,m_{k,t}+\bigl(1-\lambda_{k,t}\bigr)\,z_t,
\qquad
\lambda_{k,t}=\exp\bigl(-h_{k,t}\bigr).
\end{equation}
Il coefficiente $\lambda_{k,t}$, compreso tra zero e uno, stabilisce quanta parte della traccia sopravviva a ogni passo, e il tasso di oblio $h_{k,t}$ ne fissa l'orizzonte. Dopo $1/h_{k,t}$ passi un contributo pesa poco più di un terzo di quanto pesava appena entrato. All'inizializzazione l'orizzonte di ogni traccia coincide con il ritardo medio della banda che sostituisce, e le proiezioni $Az$ e $Bz$ vengono mediate allo stesso modo, così che la relazione resti quella dell'IRM. Nelle \emph{tracce a decadimento appreso} il tasso di oblio è un parametro per traccia, $h_{k,t}=h_k$, uguale per tutti i passi e tutti i campioni. La rete sceglie quanto lontano guardare, ma non può adattarlo a ciò che accade. Nelle \emph{tracce guidate dal sensore} l'oblio dipende anche dall'attività,
\begin{equation}
\label{eq:traccia_eventi}
h_{k,t}=h_k+\eta_k\,\frac{e_t}{\bar e},
\end{equation}
dove $e_t$ conta gli eventi registrati al passo $t$ nella zona del sensore che corrisponde alla posizione, un quadrato di $16\times16$ pixel, $\bar e$ è il suo valore medio sul training e $\eta_k\ge0$ è appreso a partire da zero. Quando la zona resta in silenzio la traccia conserva il proprio contenuto, salvo il decadimento di fondo, e quando arrivano molti eventi si rinnova in fretta. La memoria avanza così con l'attività del sensore più che con il tempo. Entrambe le versioni costano molto meno delle bande in memoria, perché conservano quattro riassunti per posizione anziché le correnti degli ultimi 15 passi, 23 su THU-EACT-50-CHL, e lo stato del sito si riduce di circa tre quarti su DVS-Lip e di oltre quattro quinti su THU-EACT-50-CHL. La figura~\ref{fig:forme_passato} confronta il peso che bande e tracce attribuiscono ai passi del passato.

Detto in modo meno tecnico, la traccia a decadimento appreso è un ricordo che sbiadisce a un ritmo costante, qualunque cosa accada nella scena. La rete sceglie durante l'addestramento quanto in fretta far sbiadire ciascuna delle quattro tracce, ma una volta scelto il ritmo vale per tutte le clip e per tutti gli istanti, sia che la bocca si muova sia che resti ferma. Nella traccia guidata dal sensore, invece, il ricordo si rinnova soprattutto quando c'è qualcosa di nuovo da ricordare. Una zona dell'immagine che resta immobile dimentica soltanto con il lento decadimento di fondo, e conserva a lungo ciò che vi è accaduto l'ultima volta, mentre una zona molto attiva sostituisce di continuo il proprio contenuto con quello più recente. La prima traccia ricorda quindi gli ultimi istanti trascorsi, la seconda piuttosto gli ultimi cambiamenti avvenuti, anche se fra l'uno e l'altro è passato molto tempo.

\begin{figure}[t]
    \centering
    \includegraphics[width=0.78\textwidth]{images/cap5/cap5_forme_passato.pdf}
    \caption[Due modi di riassumere il passato]{Due modi di riassumere il passato, rappresentati dal peso che ciascuno dei quattro riassunti attribuisce ai passi precedenti. (a)~Le bande di DVS-Lip mediano gruppi di passi entro ritardi fissi. (b)~Le tracce a decadimento appreso, con gli orizzonti raggiunti alla fine dell'addestramento su DVS-Lip, partono tutte dal passo precedente e pesano di più il passato recente.}
    \label{fig:forme_passato}
\end{figure}

L'\emph{orologio appreso} conserva le bande ma ne misura i bordi su un orologio locale $c_t$, proprio di ogni posizione, anziché sul numero di passi. A ogni passo l'orologio avanza di
\begin{equation}
\label{eq:orologio_locale}
c_t-c_{t-1}=s\,\Bigl(\frac{1}{2}+\frac{a_t}{2\,\hat a}\Bigr)^{\gamma},
\end{equation}
dove $a_t$ conta gli eventi registrati al passo $t$ nella zona del sensore che corrisponde alla posizione e $\hat a$ è un'attività tipica, misurata sul training in modo che metà degli eventi cada in zone e passi con un'attività non superiore. La rete parte dal tempo fisico, con $s=1$ e $\gamma=0$, e apprende due soli numeri, condivisi da tutte le bande e le posizioni. La scala $s$ allunga o accorcia tutte le bande, l'esponente $\gamma$ stabilisce quanto l'attività acceleri l'orologio. Se la rete lo trova utile, il passato si misura così in parte in tempo e in parte in cambiamento accumulato.

Per capire che cosa significhi lasciare alla rete il modo di contare il tempo, conviene ripartire dalle bande dell'IRM. Il loro orologio avanza di un passo ogni $27$ o $70\,\mathrm{ms}$, sempre allo stesso ritmo, e una banda che raccoglie i passi da 8 a 15 copre quindi sempre lo stesso intervallo di tempo fisico. Con l'orologio appreso le bande restano le stesse, ma l'orologio su cui se ne misurano i bordi può cambiare velocità. La scala $s$ la cambia per tutte le posizioni allo stesso modo, e un valore minore di uno rallenta l'orologio, sicché ogni banda finisce per coprire un tratto più lungo del passato. L'esponente $\gamma$ lega invece la velocità all'attività locale, e con $\gamma>0$ l'orologio di una posizione corre quando lì accade molto e rallenta quando accade poco. Una banda non copre allora un numero fisso di millisecondi, ma una certa quantità di cambiamento, che in una fase concitata si consuma in poco tempo e in una fase di quiete in molto. È lo stesso principio delle tracce guidate dal sensore, applicato alle bande e governato da due soli numeri, che la rete può anche lasciare al valore iniziale, conservando il tempo fisico.

Il \emph{passato remoto aggiunto}, infine, nasce da un risultato della sezione~\ref{sec:uso_memoria}, che ha mostrato come la rete addestrata si appoggi soprattutto alle bande più lontane dal presente. Se il passato lontano conta tanto, è naturale chiedersi se alle bande ne manchi una parte, oltre il ritardo massimo. La variante conserva quindi le quattro bande e vi aggiunge un quinto riassunto, la media di tutto il passato più vecchio dell'ultima banda, fino all'inizio della sequenza. Ha proiezioni e matrici proprie, che partono da zero, e un coefficiente che dipende dalla sua età media e dalla sua durata.

\medskip
Su THU-EACT-50-CHL tutte e cinque le alternative perdono terreno rispetto alle bande, i ritardi puntuali per quasi tre punti, le tracce a decadimento appreso per nove, quelle guidate dal sensore per poco più di due, l'orologio appreso per quattro e il passato remoto per tre. Quattro differenze su cinque sono significative o al limite della significatività. La lettura comune è che renda di più un passato riassunto a distanze fisiche definite, ed entro un orizzonte limitato. Riassumere vale più che campionare, perché le bande battono i punti, mentre un passato più lungo o misurato in modo più flessibile, come quello offerto dalle tracce, dall'orologio o dal riassunto remoto, non aiuta.

Il motivo si intuisce dalla forma delle tracce, ciascuna delle quali pesa ogni passo meno di quello che lo segue, sicché anche la più lunga dà il peso maggiore al passo appena trascorso (figura~\ref{fig:forme_passato}b). Il passato che riassume resta così mescolato con ciò che è appena accaduto, cioè proprio con la parte di passato che, secondo gli interventi della sezione~\ref{sec:uso_memoria}, conta meno per il confronto. Una banda lontana riassume invece un intervallo a distanza definita, separato dal presente dalle bande che la precedono. Durante l'addestramento la rete allunga gli orizzonti delle tracce, fino a circa cinque secondi su DVS-Lip e quattro su THU-EACT-50-CHL, e fa lo stesso con l'orologio appreso, che rallenta fino a una scala di $0{,}49$ su DVS-Lip e $0{,}82$ su THU-EACT-50-CHL. In entrambi i casi l'addestramento spinge verso un passato più lungo, ma la validazione non ne trae vantaggio.

Le tracce guidate dal sensore, invece, guadagnano sette punti su quelle a decadimento appreso e arrivano a poco più di due punti dalle bande, una distanza non significativa. La rete, inoltre, usa davvero il termine degli eventi, perché annullandolo in inferenza la Macro-F1 scende da $60{,}2$ a $50{,}8$. Con uno stato molto più piccolo, sono l'alternativa più economica fra quelle provate, e un'opzione da tenere presente quando la memoria disponibile è scarsa.

Su DVS-Lip, ancora una volta, il quadro è piatto, e tutte le alternative restano entro un punto dall'IRM, tranne le tracce a decadimento appreso, che perdono quasi tre punti come le bande senza confronto. Sono le due varianti che tolgono qualcosa, il confronto nell'una e il passato a distanza definita nell'altra. Le due che allungano il passato, l'orologio appreso e il passato remoto, finiscono appena sopra l'IRM, senza alcuna differenza significativa.

Altre prove, che spostano il confronto dentro l'attenzione del secondo stadio o nel primo patch embedding, misurano il passato in tempo di attività o lasciano al contenuto il compito di delimitarne le porzioni, non hanno migliorato le bande, e sono riportate nell'appendice~\ref{app:tabelle_estese}. Va nella stessa direzione anche una prova di natura diversa, con la rappresentazione a durata normalizzata della sezione~\ref{subsec:durata_normalizzata}, nella quale le bande restano uguali in passi ma non corrispondono più a intervalli di tempo fissi. In queste condizioni, su THU-EACT-50-CHL il guadagno dell'IRM sul modello di riferimento scende da quasi otto punti di Macro-F1 a uno soltanto.

Su THU-EACT-50-CHL i controlli raccontano dunque una storia coerente, in cui il guadagno non viene dai parametri ma dal confronto, rende di più nella forma antisimmetrica e si appoggia a un passato collocato a distanze fisiche precise. Su DVS-Lip le stesse varianti quasi non si distinguono, e peggiorano soltanto quelle che tolgono qualcosa al modulo. Prima di cercare altre memorie, conviene chiedersi perché.


% ==============================================================================
%  Due dataset, due colli di bottiglia
% ==============================================================================
% Fonti (2026-11-19). Accuratezze di training (ultime 16 epoche) e Macro-F1 di validazione: controlli5.json
% (fig_controlli.py). Correlazione di Pearson fra training e validazione sulle 11 varianti: THU 0,88 (pendenza 0,56),
% DVS-Lip 0,28 (pendenza 0,11). Voto e istogramma: fig_colli_bottiglia.py (DVS-Lip membri 57,13/57,53/57,60/57,40/56,76,
% media 57,28, voto 65,18, riconosciute da tutte 27,7%, da nessuna 16,5%; THU membri 69,17/66,67/66,67/65,41/66,17, media
% 66,82, voto 70,93, da tutte 50,6%, da nessuna 21,3%). Concordanze, ensemble e medie dei tre addestramenti dell'IRM 2 M:
% artifacts/reports/b_minus_final_inference/seed_agreement_analysis.txt (rivalutazione in inferenza, per questo le medie
% differiscono di pochi decimi dai valori salvati durante l'addestramento).

\section{Due dataset, due colli di bottiglia}
\label{sec:colli_bottiglia}

Lo stesso modulo, con la stessa ricetta, cambia il risultato su THU-EACT-50-CHL e quasi non lo sposta su DVS-Lip. Anche i controlli si comportano in modo opposto, separandosi su THU-EACT-50-CHL in un ordine coerente con le proprietà del modulo ed equivalendosi su DVS-Lip. La spiegazione più semplice sarebbe che la lettura labiale abbia poco bisogno di memoria, ma gli interventi della sezione~\ref{sec:uso_memoria} la smentiscono, perché anche su DVS-Lip la rete dipende interamente dall'ordine dei passi e si appoggia in modo massiccio al confronto. Una spiegazione migliore viene da due indicatori che con la memoria non hanno nulla a che fare, cioè quanto bene la rete si adatti ai dati di addestramento e quanto concordino fra loro modelli diversi (tabella~\ref{tab:colli_bottiglia}).

\begin{table}[t]
\centering
\small
\renewcommand{\arraystretch}{1.2}
\caption[Adattamento e varianza nei due dataset]{Adattamento e varianza nei due dataset, sulla validazione di sviluppo salvo dove indicato. Valori in percentuale. Le righe senza concordanza riportano accuratezze, e quella di training è la media delle ultime 16 epoche. La concordanza è la quota di clip su cui due modelli predicono la stessa classe. Le ultime due righe usano tre addestramenti indipendenti dell'IRM con due milioni di parametri (capitolo~\ref{chap:scalabilita}), con la concordanza media fra le tre coppie e l'accuratezza del singolo modello, in media, accanto a quella della media delle loro probabilità.}
\label{tab:colli_bottiglia}
\begin{tabular}{@{}l r r@{}}
\hline
\textbf{Indicatore} & \textbf{DVS-Lip} & \textbf{THU-EACT-50-CHL} \\
\hline
\multicolumn{3}{@{}l}{\emph{Circa mezzo milione di parametri}} \\
\quad training, riferimento / IRM & $97{,}3$ / $98{,}7$ & $74{,}2$ / $83{,}2$ \\
\quad validazione, riferimento / IRM & $56{,}8$ / $57{,}1$ & $61{,}2$ / $69{,}2$ \\
\quad concordanza riferimento e IRM, val. / test & $55{,}2$ / $51{,}8$ & $71{,}2$ / $66{,}5$ \\
\quad cinque varianti, media / voto & $57{,}3$ / $65{,}2$ & $66{,}8$ / $70{,}9$ \\
\multicolumn{3}{@{}l}{\emph{Tre addestramenti dell'IRM con due milioni di parametri}} \\
\quad concordanza, val. / test & $65{,}3$ / $61{,}3$ & $76{,}6$ / $75{,}7$ \\
\quad singolo modello / media dei tre & $64{,}2$ / $69{,}3$ & $73{,}0$ / $74{,}7$ \\
\hline
\end{tabular}
\end{table}

\titoletto{THU-EACT-50-CHL, un problema di adattamento}
Con mezzo milione di parametri il modello di riferimento non riesce a descrivere nemmeno i dati su cui è addestrato, e alla fine dell'addestramento ne riconosce il $74\%$. L'IRM porta questa quota all'$83\%$, e con essa la validazione sale di otto punti. Fra le varianti della sezione precedente, inoltre, quelle che si adattano meglio ai dati di addestramento sono anche quelle che vanno meglio in validazione, con una correlazione di $0{,}88$, e ogni punto di adattamento in più vale circa mezzo punto di Macro-F1 (figura~\ref{fig:colli_bottiglia}a). La rete compatta è a corto di capacità temporale, e la memoria relazionale gliene fornisce.

I modelli tendono anche a sbagliare le stesse clip, tanto che riferimento e IRM danno la stessa risposta sul $71\%$ delle clip di validazione e metà delle clip viene riconosciuta da tutte e cinque le varianti compatte di un piccolo gruppo di confronto (figura~\ref{fig:colli_bottiglia}b).\footnote{Su DVS-Lip le cinque varianti sono l'IRM, il modello di riferimento, i ritardi puntuali e due collocazioni alternative dell'IRM, nel primo patch embedding oppure in entrambi i patch embedding. Su THU-EACT-50-CHL sono l'IRM, i ritardi puntuali, le tracce guidate dal sensore, l'IRM in entrambi i patch embedding e l'IRM con un instradamento per contenuto, nel quale una piccola rete stabilisce, a ogni passo e in ogni posizione, quanto passato usare e come ripartirlo fra le bande. Hanno tutte circa mezzo milione di parametri e accuratezze vicine, e i dettagli sono nell'appendice~\ref{app:tabelle_estese}.} Combinarle con un voto a maggioranza aggiunge quattro punti alla loro media e meno di due alla migliore. Su THU-EACT-50-CHL ciò che una variante sa, in buona parte, lo sanno anche le altre.

\begin{figure}[t]
    \centering
    \includegraphics[width=\textwidth]{images/cap5/cap5_colli_bottiglia.pdf}
    \caption[Adattamento, validazione e accordo fra varianti]{Adattamento, validazione e accordo fra varianti, con circa mezzo milione di parametri. (a)~Accuratezza di training, media delle ultime 16 epoche, e Macro-F1 di validazione dell'IRM antisimmetrico (stella) e dei dieci controlli della sezione~\ref{sec:controlli_irm}. Su THU-EACT-50-CHL le due grandezze crescono insieme, su DVS-Lip le varianti si addensano verso il $99\%$ di training con una validazione quasi piatta. (b)~Quota delle clip di validazione riconosciute da nessuna, da una, \dots, da tutte e cinque le varianti compatte del gruppo di confronto.}
    \label{fig:colli_bottiglia}
\end{figure}

\titoletto{DVS-Lip, un problema di varianza}
Su DVS-Lip la situazione si rovescia, perché tutti i modelli imparano quasi per intero i dati di addestramento, dal 97 al $99\%$, ma riconoscono poco più della metà delle clip di validazione. Le varianti si addensano così in un angolo della figura~\ref{fig:colli_bottiglia}a, e la loro validazione si muove di pochi punti qualunque cosa accada al training, con una correlazione di appena $0{,}28$. Sul training non c'è più molto da imparare, e il limite sta altrove, nella stabilità della soluzione.

Due modelli di accuratezza quasi identica, come il riferimento e l'IRM, danno la stessa risposta soltanto sul $55\%$ delle clip di validazione e sul $52\%$ di quelle del test. Più della metà delle clip viene riconosciuta da alcune varianti e non da altre, e appena il $28\%$ da tutte e cinque. Il voto a maggioranza ne approfitta, e supera di quasi otto punti sia la media delle varianti sia la migliore di esse. Ciascuna variante coglie una parte diversa di ciò che serve, e nessuna lo coglie per intero. Ogni scelta sulla memoria cambia quindi \emph{quali} clip vengono riconosciute molto più di \emph{quante}, spostandone qualche centinaio in un senso e qualche centinaio nell'altro con un saldo che resta dentro la variabilità, ed è per questo che i controlli della sezione precedente si equivalgono.

Il voto mescola però architetture diverse, e potrebbe misurare le loro differenze più che la variabilità dell'addestramento. Una misura più pulita viene dall'addestramento ripetuto della stessa architettura, disponibile per l'IRM con due milioni di parametri che il capitolo~\ref{chap:scalabilita} introdurrà. Tre addestramenti indipendenti, che differiscono soltanto per la componente casuale dell'addestramento, danno la stessa risposta sul $65\%$ delle clip di validazione, contro il $77\%$ di THU-EACT-50-CHL, e mediarne le probabilità fa guadagnare oltre cinque punti al singolo modello, contro meno di due. Anche a parità di architettura e di ricetta, su DVS-Lip l'addestramento approda ogni volta a una soluzione diversa.

\titoletto{Una diagnosi del regime compatto}
La diagnosi, da leggere con misura, riguarda il regime compatto e non afferma che DVS-Lip sia per natura un problema di varianza e THU-EACT-50-CHL un problema di adattamento. Con più parametri anche la rete per THU-EACT-50-CHL arriva a riconoscere quasi tutto il training e mostra a sua volta un divario di generalizzazione, mentre DVS-Lip migliora con la capacità, come mostrerà il capitolo~\ref{chap:scalabilita}. Con circa mezzo milione di parametri, però, i due dataset si trovano in regimi opposti, e questo basta a spiegare perché la stessa memoria abbia effetti così diversi. Dove la rete fatica ad adattarsi ai dati e il passato offre qualcosa da confrontare, la memoria relazionale cambia il risultato. Dove la rete si adatta già e il limite è la stabilità della soluzione, nessuna memoria collocata in quel punto della rete lo sposta in modo misurabile.

Su DVS-Lip resta però una domanda che la diagnosi non scioglie, perché l'IRM, come tutte le varianti di questo capitolo, confronta passato e presente nel secondo patch embedding, dove la mappa delle feature è già stata ridotta dal pooling a $8\times8$ posizioni. Le parole di DVS-Lip si distinguono spesso per movimenti piccoli e brevi della bocca, e le coppie più simili differiscono per un dettaglio di poche decine di millisecondi. Un movimento di questo tipo, che passa fra due celle vicine della mappa a piena risoluzione, può sparire nel pooling prima che una memoria lo veda. È un'ipotesi e non una misura, ma indica per questo dataset una seconda direzione, accanto al rimedio alla varianza.

L'architettura resta dunque comune ma le diagnosi sono due, e il capitolo~\ref{chap:scalabilita} tratta ciascun dataset secondo il proprio collo di bottiglia. Per DVS-Lip affianca un rimedio alla varianza e un confronto fra passato e presente calcolato prima del pooling, fra celle vicine, mentre per THU-EACT-50-CHL conserva la memoria relazionale così com'è. Verifica poi su entrambi come la soluzione scali con la capacità e quanto se ne possa comprimere il costo, prima di confrontare i modelli finali con lo stato dell'arte.




% Capitolo 6: scalabilita', generalizzazione e configurazioni finali.
\chapter{Dalla diagnosi ai modelli finali}
\label{chap:scalabilita}

% Riscrittura del 2026-11-19 secondo docs/STRUTTURA_NARRATIVA_E_RISULTATI_CAPITOLI_5_6_7.md (versione del 19/11).
% Scaletta approvata: 6.1 distillazione (DVS-Lip, controllo su THU), 6.2 capacita', 6.3 compressione, 6.4 relazioni
% locali prima del pooling (condizionale allo scan Legion 2123991-2123995), 6.5 modelli finali e stato dell'arte,
% 6.6 dove sbagliano ancora i modelli finali. Titoletti con \titoletto (mai \paragraph).
% Stato a fine 19/11: capitolo completo. Unico segnaposto: lo scan delle relazioni locali (sezione 6.4, Legion
% 2123991-2123995); a esito noto aggiornare 6.4, la tabella tab:configurazioni_finali e tab:stato_arte_lip.
% Test ufficiale: figure_scripts/cap6/finali6.py (logit di artifacts/reports/b_minus_final_inference). Validazione:
% figure_scripts/cap6/risultati6.py. Confronti appaiati sul test: report b_minus_final_inference e
% docs/STRUTTURA_NARRATIVA..., tabelle D, E, F. Terzo addestramento dell'IRM 2 M distillato (Legion 2117616) e riferimento
% 2 M distillato (Legion 2117619) conclusi il 19/11; test e compressione valutati lo stesso giorno sui SMILIES
% (configs/evaluation/b_minus_final_inference_20261119.yaml): IRM distillato terzo addestramento 58,13 denso, 58,49 q64;
% riferimento distillato 57,99/57,50. Appaiati sul test: riferimento -> distillato +2,77 [+1,57; +4,00] p=2e-5; riferimento
% distillato -> IRM distillato +0,50/+0,58/+0,14, p 0,43/0,36/0,84. Concordanza fra i tre distillati: val 73,0, test 69,8.

Il capitolo precedente si è chiuso con una diagnosi, secondo cui con circa mezzo milione di parametri la stessa architettura incontra su THU-EACT-50-CHL un limite di adattamento, che la memoria relazionale allenta, e su DVS-Lip un limite di varianza, che nessuna memoria sposta. Questo capitolo ne trae le conseguenze, trattando prima ciascun dataset secondo il proprio collo di bottiglia. Su DVS-Lip cerca di ridurre la varianza senza toccare l'architettura, raccogliendo in un solo modello ciò che varianti diverse sanno separatamente, mentre su THU-EACT-50-CHL verifica che la memoria relazionale basti così com'è. Esamina poi due assi comuni ai due dataset, la capacità, che dice come la soluzione cambi quando la rete cresce, e la compressione, che dice quanto della memoria si possa togliere dopo l'addestramento senza perdere accuratezza. Il capitolo si chiude con i modelli finali, con il loro confronto con lo stato dell'arte e con gli errori che restano.

Tutte le scelte, comprese quelle sugli insegnanti della distillazione e sul rango della compressione, sono fatte sulla validazione di sviluppo. Il test ufficiale viene valutato a posteriori, sui modelli già scelti e sui loro controlli diretti, e serve a verificare e a confrontare, non a scegliere.


% ==============================================================================
%  Ridurre la varianza: la distillazione
% ==============================================================================
% Fonti: src/etsr/training/distillation.py (perdita (1-w) CE + w T^2 KL(p_T || softmax(z/T)), p_T media delle
% distribuzioni temperate degli insegnanti congelati sullo stesso batch dello studente; T=4, w=0,9; studente da zero) e
% src/etsr/training/self_ensemble.py (media mobile per campione delle proprie predizioni temperate, momento 0,8, peso
% crescente linearmente fino a 0,9, bersaglio usato solo se la sua classe piu' probabile coincide con l'etichetta).
% Insegnanti (config dei run kd5): DVS-Lip IRM, IRM nel primo e in entrambi i patch embedding, ritardi puntuali,
% riferimento; THU IRM, IRM in entrambi i patch embedding, IRM con instradamento, ritardi puntuali, tracce guidate dal
% sensore. Validazione: risultati6.py (Lip IRM 57,13/57,47, +SE 58,86/58,90, +KD5 62,74/62,76; THU 69,17/62,56,
% 66,92/60,98, 63,16/56,02; training u16 Lip 98,7/98,0/95,3, THU 83,2/86,6/75,7; appaiati Lip KD 455/287 p=7e-10,
% SE 411/359 p=0,066; THU KD 15/39 p=0,0015, SE 14/23 p=0,19). Test Lip KD5 55,66/55,04, +4,84 [+3,5; +6,2].
% Concordanza fra addestramenti dell'IRM 2 M, con e senza distillazione: fig_distillazione.py (65,3/61,2 contro
% 72,5/69,3). Durate dai run.log: IRM 7,7 h, distillazione da cinque insegnanti 19,3 h (circa 2,5 volte).

\section{Ridurre la varianza: la distillazione}
\label{sec:distillazione}

Su DVS-Lip varianti di accuratezza quasi identica riconoscono clip diverse, e il loro voto supera di quasi otto punti ciascuna di esse. La conoscenza che serve, dunque, in buona parte esiste già, ma è sparsa fra modelli diversi. Viene naturale chiedersi se la si possa raccogliere in un solo modello compatto, senza pagare in inferenza il costo di cinque reti.

È il compito della distillazione, che addestra un modello, detto studente, a imitare le predizioni di altri modelli già addestrati, detti insegnanti~\cite{hinton2015distilling}. Gli insegnanti sono qui le cinque varianti compatte del gruppo di confronto della sezione~\ref{sec:colli_bottiglia}, e restano congelati per tutto l'addestramento. A ogni passo ricevono lo stesso batch dello studente, e le loro distribuzioni di probabilità sulle classi, ammorbidite da una temperatura $T$, vengono mediate in un'unica distribuzione $\bar p_T$. Lo studente, che ha l'architettura dell'IRM compatto e parte da pesi casuali, minimizza
\begin{equation}
\label{eq:distillazione}
\mathcal L=(1-w)\,\mathrm{CE}(z,y)+w\,T^{2}\,\mathrm{KL}\bigl(\bar p_T\,\big\|\,\mathrm{softmax}(z/T)\bigr),
\end{equation}
dove $z$ sono i suoi logit, $y$ l'etichetta, CE l'entropia incrociata, KL la divergenza di Kullback-Leibler, $T=4$ e $w=0{,}9$, fissati a priori. Con la temperatura diventano visibili anche le classi a cui gli insegnanti assegnano una probabilità piccola, cioè quali parole o azioni considerino simili a quella corretta. È un'informazione che l'etichetta da sola non contiene.

Un bersaglio più morbido dell'etichetta potrebbe però aiutare di per sé, indipendentemente da chi lo fornisce. Per distinguere i due effetti serve un controllo che ammorbidisca i bersagli senza portare informazione da altri modelli, come l'\emph{auto-ensemble}. Per ogni clip di addestramento la rete conserva una media mobile, lungo le epoche, delle proprie predizioni ammorbidite~\cite{laine2017temporal}, e la usa come bersaglio al posto di $\bar p_T$ nell'equazione~\eqref{eq:distillazione}. Il peso del bersaglio cresce linearmente durante l'addestramento fino a $0{,}9$~\cite{kim2021selfknowledge}, e una clip lo usa soltanto quando la classe più probabile della media coincide con l'etichetta, così che la rete non venga spinta verso i propri errori.

\begin{table}[t]
\centering
\small
\renewcommand{\arraystretch}{1.2}
\caption[Interventi sull'addestramento dell'IRM compatto]{Interventi sull'addestramento dell'IRM antisimmetrico compatto, a parità di architettura e di costo in inferenza. Valori in percentuale. L'asterisco indica una differenza di accuratezza significativa nel confronto appaiato con lo stesso modello addestrato in modo standard. L'auto-ensemble non è stato valutato sul test.}
\label{tab:addestramento}
\begin{tabular}{@{}l r r r r@{}}
\hline
 & \multicolumn{2}{c}{\textbf{Validazione}} & \multicolumn{2}{c}{\textbf{Test ufficiale}} \\
\textbf{Modello} & \textbf{Accuratezza} & \textbf{Macro-F1} & \textbf{Accuratezza} & \textbf{Macro-F1} \\
\hline
\multicolumn{5}{@{}l}{\emph{DVS-Lip}} \\
IRM antisimmetrico & $57{,}1$ & $57{,}5$ & $50{,}8$ & $50{,}4$ \\
\quad con auto-ensemble & $58{,}9$ & $58{,}9$ & -- & -- \\
\quad con distillazione & $62{,}7^{*}$ & $62{,}8$ & $55{,}7^{*}$ & $55{,}0$ \\
\hline
\multicolumn{5}{@{}l}{\emph{THU-EACT-50-CHL}} \\
IRM antisimmetrico & $69{,}2$ & $62{,}6$ & $66{,}0$ & $59{,}6$ \\
\quad con auto-ensemble & $66{,}9$ & $61{,}0$ & -- & -- \\
\quad con distillazione & $63{,}2^{*}$ & $56{,}0$ & -- & -- \\
\hline
\end{tabular}
\end{table}

\titoletto{Che cosa cambia su DVS-Lip}
Su DVS-Lip la distillazione è la leva più forte trovata nel regime compatto (tabella~\ref{tab:addestramento}). Porta l'IRM da $57{,}5$ a $62{,}8$ di Macro-F1 in validazione, correggendo 455 errori e introducendone 287, e sul test ufficiale ne alza l'accuratezza di quasi cinque punti, con un intervallo di confidenza che va da $+3{,}5$ a $+6{,}2$. Lo studente recupera circa due terzi della distanza fra le singole varianti e il loro voto, pur restando un solo modello con il costo in inferenza dell'IRM compatto. Ci arriva adattandosi un po' meno ai dati di addestramento, che riconosce al $95\%$ anziché al $99\%$, e generalizzando molto meglio (figura~\ref{fig:distillazione}a). È il comportamento atteso da un rimedio alla varianza, che rinuncia a un po' di aderenza al training in cambio di una soluzione più stabile.

L'auto-ensemble, invece, si ferma a un punto e mezzo, e il confronto appaiato non lo distingue dal caso. Ammorbidire i bersagli non basta, e il grosso del guadagno viene dall'informazione portata da modelli diversi.

\begin{figure}[t]
    \centering
    \includegraphics[width=\textwidth]{images/cap6/cap6_distillazione.pdf}
    \caption[Effetto della distillazione]{Effetto della distillazione. (a)~Accuratezza di training, media delle ultime 16 epoche, e Macro-F1 di validazione dell'IRM antisimmetrico compatto (stella), addestrato con auto-ensemble (cerchio vuoto) e con distillazione da cinque insegnanti (cerchio pieno). Su DVS-Lip la distillazione riduce l'adattamento al training e alza la validazione, su THU-EACT-50-CHL li riduce entrambi. (b)~Quota delle clip su cui due addestramenti indipendenti dell'IRM con due milioni di parametri danno la stessa risposta, senza e con distillazione, in media sulle tre coppie di tre addestramenti.}
    \label{fig:distillazione}
\end{figure}

La distillazione riduce proprio la varianza che la sezione~\ref{sec:colli_bottiglia} aveva individuato, e lo si può verificare in modo diretto. Con due milioni di parametri, dove l'addestramento è stato ripetuto (sezione~\ref{sec:capacita_irm}), tre modelli distillati danno in media la stessa risposta sul $73\%$ delle clip di validazione e sul $70\%$ di quelle del test, contro il $65\%$ e il $61\%$ di tre modelli addestrati in modo standard (figura~\ref{fig:distillazione}b). Lo studente non diventa soltanto più accurato. Approda, da un addestramento all'altro, a soluzioni più simili fra loro.

In inferenza lo studente è identico all'IRM compatto, e il prezzo si paga durante l'addestramento, per il quale servono cinque insegnanti già addestrati e, a ogni passo, anche il loro calcolo in avanti. Nelle prove svolte l'addestramento dello studente è durato circa due volte e mezzo quello dell'IRM, senza contare quello degli insegnanti.

\titoletto{Il controllo su THU-EACT-50-CHL}
Se la diagnosi è corretta, la stessa ricetta non dovrebbe aiutare THU-EACT-50-CHL, dove il limite è l'adattamento e non la varianza, ed è proprio quello che accade. Con cinque insegnanti compatti dello stesso tipo, la distillazione abbassa la Macro-F1 dell'IRM da $62{,}6$ a $56{,}0$, in modo significativo, e l'auto-ensemble la abbassa di un punto e mezzo. Il risultato è comprensibile, perché lo studente ha già l'architettura del migliore degli insegnanti e il loro voto supera l'IRM di meno di due punti, sicché imitarne la media significa in parte imitarne gli errori. Inoltre una rete che già fatica ad adattarsi ai dati di addestramento, con il bersaglio degli insegnanti vi si adatta ancora meno, e ne riconosce il $76\%$ anziché l'$83\%$ (figura~\ref{fig:distillazione}a).

Su THU-EACT-50-CHL, dunque, la memoria relazionale è già la risposta al collo di bottiglia, e la configurazione per questo dataset resta l'IRM antisimmetrico con la ricetta standard. Né la distillazione né l'auto-ensemble lo migliorano, e nessuna delle memorie alternative della sezione~\ref{sec:controlli_irm} supera le bande. La leva che resta è la capacità, che di fronte a un limite di adattamento è anche la più naturale.


% ==============================================================================
%  Capacita'
% ==============================================================================
% Fonti: run capacity-1m/2m (embed_dim 192 e 256, rango 16) di IRM e riferimento VC; validazione da risultati6.py e
% validation_controls_vs_bminus.txt; test da finali6.py. Distillazione da quattro insegnanti (config dei run kd4):
% IRM 2 M, IRM 1 M, riferimento 2 M, IRM compatto addestrato con VC stocastico (scala temporale log-uniforme in
% [0,9; 1,1] e fase della griglia casuale, solo in addestramento); scelti sulla validazione fra 14 modelli come il gruppo
% di quattro con la migliore accuratezza del voto ammorbidito (69,65%); bersagli precalcolati, durata come
% l'addestramento standard (2 M: 12,9 h contro 13,2 h). Confronti appaiati sul test (documento di struttura, tab. E e F):
% THU IRM-rif 0,5 M +9,93 (p<0,001), 1 M +3,12 [-0,2; +6,6] (p=0,086), 2 M media tre coppie +1,96 [-0,18; +4,17] acc,
% +2,17 [+0,25; +4,11] F1, singole coppie +2,02/+2,21/+1,65; validazione 2 M +0,08 acc, +0,37 F1. DVS-Lip IRM-rif 0,5 M
% +0,80 (p=0,29), 1 M +3,18 [+1,7; +4,6] (p<0,001), 2 M +1,09 (p=0,12, primo addestramento); IRM 2 M - 0,5 M +5,49;
% distillazione 1 M +1,49 (p=0,029), 2 M +2,19 e +2,91 (p<0,001), media +2,55 [+1,66; +3,46]; KD5 contro KD4 a 0,5 M
% 55,66 contro 54,25 (p=0,035). Accuratezza di training del riferimento THU: 1 M 98,4%, 2 M 100,0%.

\section{La soluzione al crescere del budget}
\label{sec:capacita_irm}

Fin qui tutti i modelli avevano circa mezzo milione di parametri, e la prova di questa sezione allarga la rete da 128 a 192 e a 256 canali, come per il modello di riferimento nella sezione~\ref{sec:solidita_riferimento}, e porta i parametri a circa uno e due milioni. Ogni modello è addestrato da zero con la stessa ricetta. Il rango della relazione resta $r=16$, ma il modulo cresce insieme alla rete, perché le matrici $W_k$ crescono con il quadrato del numero di canali e le proiezioni $A$, $B$ e $U_k$ in proporzione a esso. Il termine di confronto è il modello di riferimento della stessa larghezza, addestrato con la stessa rappresentazione a lunghezza variabile. Le versioni più larghe del capitolo~\ref{chap:riallocazione_budget} usano la finestra fissa e non possono fare da controllo. Con due milioni di parametri l'addestramento dell'IRM è stato ripetuto tre volte su entrambi i dataset, e così quello del modello di riferimento su THU-EACT-50-CHL, e i valori sono riportati come media e deviazione standard.

Su DVS-Lip la sezione precedente suggerisce di affiancare alla capacità la distillazione. Per uno studente da uno o due milioni di parametri, però, le varianti compatte sono insegnanti più piccoli dello studente stesso. Si usano allora quattro insegnanti scelti sulla validazione fra quattordici modelli già addestrati, come il gruppo il cui voto sulle probabilità ammorbidite raggiunge l'accuratezza più alta. Sono l'IRM da uno e da due milioni di parametri, il modello di riferimento da due milioni e un IRM compatto addestrato con piccole perturbazioni casuali della scala temporale dell'ingresso. Le loro predizioni vengono calcolate una volta sola prima dell'addestramento, sicché distillare costa quanto addestrare in modo standard. Con lo studente compatto restano preferibili i cinque insegnanti compatti della sezione precedente, che sul test danno $55{,}7\%$ di accuratezza contro $54{,}3\%$, e sono quindi la ricetta del modello compatto, mentre i quattro insegnanti sono quella dei modelli più grandi.

\begin{table}[t]
\centering
\small
\renewcommand{\arraystretch}{1.2}
\caption[Modelli al crescere della larghezza]{Modello di riferimento, IRM antisimmetrico e IRM distillato al crescere della larghezza. Macro-F1 di validazione e accuratezza sul test ufficiale, in percentuale. I parametri sono quelli su DVS-Lip, su THU-EACT-50-CHL sono di poco inferiori. Con 256 canali, media e deviazione standard di tre addestramenti, salvo i due modelli di riferimento su DVS-Lip, addestrati una volta sola.}
\label{tab:capacita_irm}
\setlength{\tabcolsep}{4.5pt}
\begin{tabular}{@{}l r r r r r@{}}
\hline
 & & \multicolumn{2}{c}{\textbf{DVS-Lip}} & \multicolumn{2}{c}{\textbf{THU-EACT-50-CHL}} \\
\textbf{Modello} & \textbf{Parametri} & \textbf{Valid.} & \textbf{Test} & \textbf{Valid.} & \textbf{Test} \\
\hline
\multicolumn{6}{@{}l}{\emph{128 canali}} \\
\quad riferimento & $0{,}50$\,M & $56{,}9$ & $50{,}0$ & $54{,}9$ & $56{,}1$ \\
\quad IRM & $0{,}51$\,M & $57{,}5$ & $50{,}8$ & $62{,}6$ & $66{,}0$ \\
\quad IRM distillato & $0{,}51$\,M & $62{,}8$ & $55{,}7$ & & \\
\multicolumn{6}{@{}l}{\emph{192 canali}} \\
\quad riferimento & $1{,}11$\,M & $60{,}2$ & $51{,}6$ & $64{,}9$ & $68{,}8$ \\
\quad IRM & $1{,}13$\,M & $61{,}8$ & $54{,}8$ & $65{,}6$ & $71{,}9$ \\
\quad IRM distillato & $1{,}13$\,M & $64{,}3$ & $56{,}3$ & & \\
\multicolumn{6}{@{}l}{\emph{256 canali}} \\
\quad riferimento & $1{,}97$\,M & $63{,}4$ & $55{,}2$ & $68{,}0\pm1{,}8$ & $72{,}3\pm0{,}9$ \\
\quad IRM & $1{,}99$\,M & $64{,}0\pm0{,}9$ & $55{,}8\pm0{,}5$ & $67{,}7\pm1{,}5$ & $74{,}3\pm1{,}0$ \\
\quad IRM distillato & $1{,}99$\,M & $67{,}1\pm0{,}6$ & $58{,}4\pm0{,}2$ & & \\
\quad riferimento distillato & $1{,}97$\,M & $66{,}1$ & $58{,}0$ & & \\
\hline
\end{tabular}
\end{table}

\begin{figure}[t]
    \centering
    \includegraphics[width=\textwidth]{images/cap6/cap6_capacita.pdf}
    \caption[Accuratezza sul test al crescere dei parametri]{Accuratezza sul test ufficiale al crescere dei parametri, in scala logaritmica. Con due milioni di parametri, media e deviazione standard degli addestramenti ripetuti. I numeri indicano il margine dell'IRM antisimmetrico sul modello di riferimento della stessa larghezza, in punti. Su DVS-Lip la distillazione usa cinque insegnanti compatti per il modello da mezzo milione di parametri e quattro insegnanti per i modelli più grandi, e la stella vuota indica il modello di riferimento distillato con gli stessi quattro insegnanti, disegnato poco a sinistra della sua posizione per non coprire l'IRM distillato.}
    \label{fig:capacita}
\end{figure}

\titoletto{THU-EACT-50-CHL}
Entrambe le architetture migliorano con il budget (tabella~\ref{tab:capacita_irm}), e l'IRM passa sul test dal $66{,}0$ al $74{,}3\%$ di accuratezza. Il suo margine sul modello di riferimento, però, si riduce al crescere della capacità (figura~\ref{fig:capacita}). Vale quasi dieci punti con mezzo milione di parametri, tre con un milione e due con due milioni. A due milioni il margine resta positivo in tutti e tre gli addestramenti ripetuti, e la differenza di Macro-F1 sul test, mediata sui tre, è significativa, con un intervallo che va da $+0{,}3$ a $+4{,}1$ punti. In validazione, invece, il margine si annulla.

Il quadro è coerente con la diagnosi, perché con abbastanza capacità anche il modello di riferimento riesce ad adattarsi ai dati di addestramento, che a un milione di parametri riconosce già al $98\%$, e la differenza nel modo di rappresentare il tempo pesa meno. La memoria relazionale conta soprattutto dove il budget è scarso, che è proprio il regime per cui questa tesi l'ha pensata. Un milione di parametri, d'altra parte, è anche la soglia oltre la quale, sul test, la rete con l'IRM supera i metodi pubblicati che non usano dati esterni, mentre il modello di riferimento della stessa larghezza resta sotto (sezione~\ref{sec:configurazioni_finali}).

\titoletto{DVS-Lip}
Su DVS-Lip la capacità sposta il risultato più di qualunque memoria, e dal modello compatto a quello con due milioni di parametri l'IRM guadagna cinque punti e mezzo di accuratezza sul test. A ogni larghezza resta davanti al modello di riferimento, ma con un margine variabile, piccolo con mezzo milione e con due milioni di parametri, netto e significativo soltanto con un milione, dove vale tre punti.

La distillazione si somma alla capacità, e aggiunge sul test un punto e mezzo a un milione di parametri e più di due punti e mezzo a due milioni, con guadagni significativi in tutti gli addestramenti ripetuti, anche se più piccoli dei quasi cinque punti del modello compatto. Il modello migliore su DVS-Lip è così l'IRM con due milioni di parametri e distillazione, che arriva in media al $58{,}4\%$ sul test, con tre addestramenti che differiscono di pochi decimi. Lo stesso guadagno, però, va anche al modello di riferimento, che con gli stessi quattro insegnanti passa dal $55{,}2$ al $58{,}0\%$. A parità di ricetta l'IRM distillato resta davanti di meno di mezzo punto, e nessuno dei tre confronti appaiati distingue questa differenza dal caso. Su DVS-Lip, dunque, il risultato finale si deve soprattutto alla capacità e alla distillazione, e la memoria relazionale vi contribuisce poco, come già nel regime compatto. Per avere un'idea di quanto valga la ricetta basta notare che il modello compatto distillato, con $55{,}7\%$, raggiunge sul test l'IRM da due milioni di parametri addestrato in modo standard, con un quarto dei parametri.


% ==============================================================================
%  Compressione dopo l'addestramento
% ==============================================================================
% Fonti (2026-11-19): figure_scripts/cap6/compressione6.py sui report di artifacts/reports/b_minus_final_inference
% (sola inferenza; rango scelto sulla validazione con la regola del manifest, tolleranza 0,5 punti di Macro-F1; costi per
% passo su 32 clip di validazione, funzione cost() di notebooks/b_minus_thesis.py; stato in valori da 32 bit).
% Ranghi scelti: THU 24 / 32 / 24, 24, 16; DVS-Lip 64 / 64 / 64, 32, 48; distillati DVS-Lip: cinque insegnanti nessuno
% (migliore q48, -0,71), quattro insegnanti 0,5 M q96, 1 M q96, 2 M q128, 64, 64. Test denso -> compresso, nessuna
% differenza significativa: THU 65,99->65,99, 71,88->70,40 (p=0,17), 2 M media 74,26->74,20; DVS-Lip 50,81->50,79,
% 54,81->54,73, 2 M media 55,79->55,74. Risparmi del sito (MAC): THU -64,5 / -70,4 / -80,7..-84,3%; DVS-Lip -31,6 /
% -51,9 / -62,9..-77,1%; rete: MAC THU -8,6..-10,8%, DVS-Lip -4,1..-9,6%; stato THU -14,9..-17,6%, DVS-Lip -7,3..-13,0%.
% Moltiplicazioni elemento per elemento della relazione, a 0,5 M: circa 41 mila (DVS-Lip) e 239 mila (THU) per passo.
% Figura: figure_scripts/cap6/fig_compressione.py.

\section{Comprimere la memoria}
\label{sec:compressione}

Nella forma descritta finora l'IRM antisimmetrico costa un po' più del mixer a ritardi che sostituisce, con il $2{,}5\%$ di parametri in più nell'intera rete compatta e un quarto di stato in più nel suo sito (sezione~\ref{sec:relazione}). Quando la rete si allarga, inoltre, il sito cresce più in fretta del resto, perché le matrici $W_k$ crescono con il quadrato dei canali. Questa sezione mostra che il costo si può ridurre dopo l'addestramento, senza riaddestrare la rete, fino a scendere sotto quello del mixer a ritardi a ogni larghezza.

\titoletto{Una base comune per le quattro bande}
La maggior parte dei parametri e delle moltiplicazioni del modulo sta nelle quattro matrici $W_k$, di $C\times C$ elementi, che riportano nei canali i riassunti delle bande. Le matrici sono apprese separatamente, ma leggono tutte lo stesso segnale, la corrente del secondo patch embedding, ed è plausibile che usino in buona parte le stesse direzioni dello spazio dei canali. In tal caso basta un numero ridotto di direzioni comuni per descriverle tutte.

Per trovarle si impilano le quattro matrici una sopra l'altra e se ne calcola la decomposizione ai valori singolari. Le $q$ direzioni più importanti dello spazio dei canali, cioè i primi $q$ vettori singolari destri, formano una base condivisa $P\in\mathbb R^{q\times C}$, e ogni matrice viene sostituita dalla sua approssimazione in quella base,
\begin{equation}
\label{eq:base_condivisa}
W_k\approx D_k\,P,
\qquad
D_k=W_k\,P^{\top}\in\mathbb R^{C\times q}.
\end{equation}
Poiché la media di banda e la proiezione sono entrambe lineari, il loro ordine si può scambiare, e $W_k\,m_{k,t}\approx D_k\,\overline{(Pz)}_{k,t}$, dove la barra indica la media nella banda. Basta allora proiettare la corrente presente una sola volta, nel momento in cui la si scrive nella memoria, e conservare $Pz$, di $q$ valori, al posto di $z$, che ne ha $C$. Il ramo relazionale non cambia, e le proiezioni $Az$ e $Bz$ restano quelle addestrate. Il rango $r=16$ della relazione, del resto, non va confuso con il rango $q$ della base. Il resto della rete resta identico, e la compressione non richiede alcun addestramento.

\titoletto{Una regola fissata in anticipo}
Resta da scegliere $q$, e la scelta va fatta senza guardare il test. La regola, fissata prima di qualunque valutazione e applicata allo stesso modo a tutti i modelli, prende il rango più piccolo, fra 8, 16, 24, 32, 48, 64, 96, 128 e 192, con cui la Macro-F1 di validazione perde al massimo mezzo punto rispetto al modello non compresso. Se nessun rango soddisfa la regola, il modello non viene compresso. Il test ufficiale serve soltanto a verificare a posteriori che la compressione non abbia tolto nulla. Poiché il rango viene scelto sulla stessa validazione su cui si misura la tolleranza, i piccoli guadagni che a volte vi compaiono non vanno letti come un beneficio della compressione.

\begin{figure}[t]
    \centering
    \includegraphics[width=\textwidth]{images/cap6/cap6_compressione.pdf}
    \caption[Compressione dell'IRM dopo l'addestramento]{Compressione dell'IRM antisimmetrico dopo l'addestramento, senza riaddestramento. (a)~Differenza di Macro-F1 di validazione rispetto al modello non compresso al variare del rango $q$ della base condivisa. La banda grigia indica la tolleranza di mezzo punto, i cerchi grandi i ranghi scelti. Con due milioni di parametri è mostrato il primo dei tre addestramenti. (b)~Moltiplicazioni con accumulo per passo nel sito del secondo patch embedding, per il mixer a ritardi del modello di riferimento e per l'IRM, prima e dopo la compressione.}
    \label{fig:compressione}
\end{figure}

\titoletto{Quanto si può togliere}
Su THU-EACT-50-CHL le matrici sono fortemente ridondanti, e bastano 24 direzioni su 128 nel modello compatto, e da 16 a 24 su 256 in quello da due milioni di parametri, mentre su DVS-Lip ne servono 64, e da 32 a 64 con due milioni di parametri (figura~\ref{fig:compressione}a). Le quattro bande di DVS-Lip sembrano usare direzioni più diverse fra loro. Il rango scelto, inoltre, cambia da un addestramento all'altro della stessa architettura, e dipende quindi dal singolo modello più che dal dataset.

Il risparmio cresce con la larghezza, perché il rango necessario cresce molto più lentamente del numero di canali. Su THU-EACT-50-CHL le moltiplicazioni del sito si riducono del $65\%$ nel modello compatto, del $70\%$ con un milione di parametri e di oltre l'$80\%$ con due, e lo stato del sito in misura simile (figura~\ref{fig:compressione}b). Su DVS-Lip la riduzione va dal $32\%$ del modello compatto a valori fra il $63$ e il $77\%$ con due milioni di parametri. Sull'intera rete il risparmio è più piccolo, perché il sito ne è soltanto una parte, e vale fra il 4 e l'$11\%$ delle moltiplicazioni e fra il 7 e il $18\%$ dello stato.

La compressione, d'altra parte, non costa accuratezza, e sul test nessuna differenza fra modello compresso e non compresso è significativa. La più ampia è una perdita di un punto e mezzo su THU-EACT-50-CHL con un milione di parametri, che il confronto appaiato non distingue dal caso, e con due milioni di parametri la media dei tre addestramenti passa dal $74{,}26$ al $74{,}20\%$.

\titoletto{Meno costosa del mixer che sostituisce}
La tabella~\ref{tab:compressione} mette a confronto, a ogni larghezza, il modello di riferimento con la rete che porta l'IRM compresso. Il secondo ha meno parametri, meno moltiplicazioni e meno stato nel sito, e quindi anche nell'intera rete, su entrambi i dataset e a ogni budget, ed è più accurato sul test in ogni caso. Su THU-EACT-50-CHL, con due milioni di parametri, il sito compresso esegue un quinto delle moltiplicazioni del mixer a ritardi e conserva un quinto del suo stato. Le moltiplicazioni elemento per elemento della relazione, non incluse nel conteggio, sono circa un millesimo di quelle con accumulo dell'intera rete. Il confronto fra passato e presente, insomma, non richiede un budget più grande, ma un uso diverso dello stesso budget, nello spirito della riallocazione del capitolo~\ref{chap:riallocazione_budget}.

\begin{table}[t]
\centering
\small
\renewcommand{\arraystretch}{1.2}
\caption[Modello di riferimento e IRM compresso a ogni larghezza]{Modello di riferimento e rete con l'IRM compresso a ogni larghezza. Parametri dell'intera rete, moltiplicazioni con accumulo per passo e stato persistente del sito del secondo patch embedding, accuratezza sul test ufficiale. Con due milioni di parametri il rango varia fra i tre addestramenti, i costi sono quelli del rango più grande e l'accuratezza è la media dei tre. Il modello di riferimento su DVS-Lip con due milioni di parametri è stato addestrato una volta sola.}
\label{tab:compressione}
\begin{adjustbox}{max width=\textwidth}
\begin{tabular}{@{}l r r r r c r r r r@{}}
\hline
 & \multicolumn{4}{c}{\textbf{Modello di riferimento}} & \multicolumn{5}{c}{\textbf{IRM compresso}} \\
\textbf{Budget} & \textbf{Parametri} & \textbf{MAC} & \textbf{Stato} & \textbf{Test} & \textbf{Rango} & \textbf{Parametri} & \textbf{MAC} & \textbf{Stato} & \textbf{Test} \\
 & & [M] & [migliaia] & [\%] & & & [M] & [migliaia] & [\%] \\
\hline
\multicolumn{10}{@{}l}{\emph{THU-EACT-50-CHL}} \\
\quad $0{,}5$\,M & $0{,}49$\,M & $24{,}5$ & $1\,101$ & $56{,}1$ & 24 & $0{,}46$\,M & $10{,}3$ & $482$ & $66{,}0$ \\
\quad 1\,M & $1{,}10$\,M & $55{,}1$ & $1\,652$ & $68{,}8$ & 32 & $1{,}01$\,M & $18{,}4$ & $551$ & $70{,}4$ \\
\quad 2\,M & $1{,}96$\,M & $98{,}0$ & $2\,202$ & $72{,}3$ & 16--24 & $1{,}75$\,M & $20{,}7$ & $482$ & $74{,}2$ \\
\multicolumn{10}{@{}l}{\emph{DVS-Lip}} \\
\quad $0{,}5$\,M & $0{,}50$\,M & $4{,}2$ & $123$ & $50{,}0$ & 64 & $0{,}49$\,M & $3{,}4$ & $92$ & $50{,}8$ \\
\quad 1\,M & $1{,}11$\,M & $9{,}4$ & $184$ & $51{,}6$ & 64 & $1{,}05$\,M & $5{,}1$ & $92$ & $54{,}7$ \\
\quad 2\,M & $1{,}97$\,M & $16{,}8$ & $246$ & $55{,}2$ & 32--64 & $1{,}81$\,M & $6{,}8$ & $92$ & $55{,}7$ \\
\hline
\end{tabular}
\end{adjustbox}
\end{table}

\titoletto{I modelli distillati}
I modelli distillati si lasciano comprimere meno, tanto che per lo studente compatto con cinque insegnanti nessun rango rispetta la tolleranza, e il migliore perde sette decimi di punto. Quello con quattro insegnanti scende soltanto a $q=96$, con un risparmio marginale. Con uno e due milioni di parametri la compressione torna utile, con ranghi da 64 a 128 e un risparmio sul sito fra il 33 e il $63\%$, sufficiente a restare sotto il mixer a ritardi della stessa larghezza. Una spiegazione compatibile è che la distillazione distribuisca l'informazione su più direzioni delle matrici $W_k$, ma la prova non basta a dimostrarlo. In pratica, su DVS-Lip il modello compatto più accurato, quello distillato, resta non compresso, con un sito che costa circa un quinto più del mixer a ritardi, mentre il modello compatto compresso è quello addestrato in modo standard. Il dettaglio di tutti i ranghi provati è nell'appendice~\ref{app:tabelle_estese}.


% ==============================================================================
%  Relazioni locali prima del pooling
% ==============================================================================
% Fonti (2026-11-19): src/etsr/models/local_prepool.py (R[d,delta,t,p] = u[t-d,p] u[t,p+delta] - u[t-d,p+delta] u[t,p],
% u = proiezione 1x1 a rango 16 con init ortogonale, spostamenti (0,1), (1,0), (1,1), (1,-1), ritardi 1, 2, 4, media 3x3
% con passo 2 sulla griglia ridotta, uscita 1x1 per ritardo inizializzata a zero, sommata alla corrente dopo il max
% pooling; stato = proiezioni degli ultimi 4 passi). Sito: patch_embed1.main3, 32 canali a 64x64. Parametri +6.656
% (519.980 su DVS-Lip, 513.530 su THU). Run 0,5 M: dvslip ...local-prepool-r16__20261118_133930_643475 (59,23/59,20),
% THU ...__20261118_141606_059420 (69,17/63,08). Confronti con l'IRM (report validation_controls_vs_bminus.txt e documento
% di struttura, sezione 5.5): DVS-Lip val +2,10 acc [+0,1; +3,9] p=0,037, F1 +1,73 [-0,2; +3,6]; test 53,41/52,62, +2,59
% [+1,2; +4,0] p<0,001, F1 +2,23; THU val 0,00 p=1, F1 +0,52. Contro la distillazione compatta sul test: 53,41 contro
% 55,66, -2,25 (p=0,002). Durata: 16,8 h contro 7,7 h dell'IRM (run.log). Scan in corso su Legion: 2123991 (2 M + KD4),
% 2123992 (2 M), 2123993 (1 M + KD4), 2123994 (1 M), 2123995 (0,5 M + KD4). Il test a 0,5 M e' stato valutato il 19/11,
% dopo la decisione di estendere la variante, presa sulla validazione. Costo dal profilo (lip_prepool.json): ramo circa
% 8,0 M MAC per passo (rete 38,4 -> 46,5 M), stato +262.144 valori (rete 842 -> 1.104 mila), moltiplicazioni elemento per
% elemento 1,6 M per passo.

\section{Relazioni locali prima del pooling}
\label{sec:relazioni_locali}

Le sezioni precedenti hanno trovato per DVS-Lip due leve efficaci, la distillazione e la capacità, ma nessuna che riguardi la memoria. Resta da chiedersi se il confronto fra passato e presente possa aiutare anche questo compito, a una scala diversa da quella su cui lavora l'IRM.

L'IRM confronta passato e presente nel secondo patch embedding, su feature che hanno già attraversato diversi pooling e che su DVS-Lip occupano una griglia di $8\times8$ posizioni, e a quel punto parte del dettaglio spaziale è ormai andata perduta. Le parole di DVS-Lip, però, si distinguono spesso per movimenti piccoli e brevi della bocca, e le coppie di parole più simili differiscono per poco. L'ipotesi è che un confronto calcolato prima, fra celle vicine e con ritardi brevi, possa conservare proprio quelle transizioni fini che il pooling cancella.

\titoletto{Un correlatore fra celle vicine}
La variante aggiunge un ramo all'ultimo strato del frontend piramidale che lavora a risoluzione piena, con 32 canali su una griglia di $64\times64$ celle, subito prima del pooling che la dimezza. Una proiezione appresa riduce la corrente $x_t$ di ogni cella a $r=16$ componenti, $u_t=Px_t$. Per ogni coppia di celle vicine $p$ e $p+\delta$, lungo quattro direzioni, orizzontale, verticale e le due diagonali, e per tre ritardi brevi $d\in\{1,2,4\}$ passi, il ramo calcola
\begin{equation}
\label{eq:relazione_locale}
R_{d,\delta,t}(p)=u_{t-d}(p)\odot u_t(p+\delta)-u_{t-d}(p+\delta)\odot u_t(p).
\end{equation}
È il correlatore dell'equazione~\eqref{eq:reichardt} quasi alla lettera, con due celle vicine dell'immagine al posto dei due recettori e un ritardo di pochi passi, da $27$ a $110\,\mathrm{ms}$ su DVS-Lip, sicché il segno della risposta indica in quale verso qualcosa si è spostato da una cella all'altra. Le relazioni vengono poi mediate su piccoli gruppi di celle, per portarle alla risoluzione dopo il pooling, e riportate nei canali da una matrice appresa per ciascun ritardo, inizializzata a zero, che le somma alla corrente come una correzione prima del neurone. All'inizializzazione il ramo non contribuisce, e la rete parte di nuovo dalla funzione dell'IRM. L'IRM nel secondo patch embedding e il mixer a ritardi del primo restano al loro posto.

Il ramo costa poco in parametri, $6\,656$ in più su circa mezzo milione, ma non in calcolo e in memoria. La sua memoria conserva soltanto le proiezioni a 16 componenti degli ultimi quattro passi, ma per ciascuna delle 4\,096 celle della griglia, cioè $262\,144$ valori, quasi un terzo in più dello stato dell'intera rete compatta. Il ramo aggiunge inoltre circa otto milioni di moltiplicazioni con accumulo per passo, un quinto in più della rete, e l'addestramento richiede quasi il doppio del tempo.

\begin{table}[t]
\centering
\small
\renewcommand{\arraystretch}{1.2}
\caption[Relazioni locali prima del pooling, circa mezzo milione di parametri]{Relazioni locali prima del pooling, con circa mezzo milione di parametri. Valori in percentuale, differenze in punti rispetto all'IRM antisimmetrico. L'asterisco indica una differenza di accuratezza significativa nel confronto appaiato.}
\label{tab:relazioni_locali}
\begin{adjustbox}{max width=\textwidth}
\begin{tabular}{@{}l r r r r@{}}
\hline
 & \multicolumn{2}{c}{\textbf{IRM}} & \multicolumn{2}{c}{\textbf{IRM e relazioni locali}} \\
\textbf{Dataset e partizione} & \textbf{Accuratezza} & \textbf{Macro-F1} & \textbf{Accuratezza} & \textbf{Macro-F1} \\
\hline
DVS-Lip, validazione & $57{,}1$ & $57{,}5$ & $59{,}2$ \;($+2{,}1^{*}$) & $59{,}2$ \;($+1{,}7$) \\
DVS-Lip, test ufficiale & $50{,}8$ & $50{,}4$ & $53{,}4$ \;($+2{,}6^{*}$) & $52{,}6$ \;($+2{,}2$) \\
THU-EACT-50-CHL, validazione & $69{,}2$ & $62{,}6$ & $69{,}2$ \;($0{,}0$) & $63{,}1$ \;($+0{,}5$) \\
\hline
\end{tabular}
\end{adjustbox}
\end{table}

\titoletto{Il risultato con mezzo milione di parametri}
Su DVS-Lip il ramo migliora l'IRM di due punti di accuratezza in validazione e di due e mezzo sul test ufficiale, in entrambi i casi in modo significativo (tabella~\ref{tab:relazioni_locali}). Sul test corregge 731 errori dell'IRM e ne introduce 602. È l'unica variante della memoria, fra quelle provate, che sposti DVS-Lip in modo misurabile nel regime compatto. Su THU-EACT-50-CHL, invece, non cambia nulla, e l'accuratezza di validazione resta identica. Il contrasto è coerente con l'ipotesi del dettaglio fine, anche se, come mostrerà la sezione~\ref{sec:errori}, il guadagno non viene dagli scambi fra le parole più simili. Una lettura plausibile è che le azioni di THU-EACT-50-CHL dipendano soprattutto da movimenti ampi e da distanze di secondi, già alla portata dell'IRM, mentre la lettura labiale dipende anche da transizioni locali di poche decine di millisecondi.

Il ramo resta comunque sotto la distillazione, che sul test porta il modello compatto al $55{,}7\%$ contro il $53{,}4\%$, con una differenza significativa. Le due leve, d'altra parte, agiscono su cose diverse, perché le relazioni locali aggiungono informazione temporale che la rete prima non vedeva, mentre la distillazione riduce la varianza della soluzione. Se i loro effetti si sommino, e se il guadagno resista al crescere della capacità, è la domanda di una verifica ancora in corso.

\titoletto{La verifica in corso}
La verifica ripete la variante con uno e due milioni di parametri, con e senza la distillazione da quattro insegnanti, e anche con la distillazione nel modello compatto, sempre con la stessa ricetta e gli stessi insegnanti dell'IRM. Secondo una regola fissata in anticipo, le relazioni locali diventano parte della configurazione finale di DVS-Lip soltanto se migliorano l'IRM in validazione in modo coerente, con lo stesso segno a uno e a due milioni di parametri, con e senza distillazione. La conferma decisiva è il confronto con due milioni di parametri e distillazione. Il test ufficiale verrà valutato su tutte le configurazioni concluse, come per gli altri modelli, ma non entrerà nella decisione. Su THU-EACT-50-CHL la variante non viene estesa, perché nel modello compatto non dà alcun guadagno. \textbf{[In corso: risultati della verifica a uno e due milioni di parametri, con e senza distillazione. Se la regola sarà soddisfatta, la sezione diventerà la base della configurazione finale di DVS-Lip. Altrimenti resterà come prospettiva, con i numeri ottenuti.]}


% ==============================================================================
%  Modelli finali e stato dell'arte
% ==============================================================================
% Fonti (2026-11-19). Modelli di questa tesi: finali6.py e compressione6.py (test ufficiale, logit di
% artifacts/reports/b_minus_final_inference). THU compresso: 0,5 M q24 65,99/59,68; 1 M q32 70,40/64,38; 2 M (q24, 24,
% 16) 74,20+-1,29 / 68,21+-1,35. DVS-Lip: distillato compatto 55,66/55,04 (non comprimibile); distillato 2 M
% 58,40+-0,24 / 57,88+-0,21 (compresso q128, 64, 64: 58,31 / 57,79); Acc1/Acc2 dai gruppi di
% configs/dvslip_class_groups.json (IRM 0,5 M 42,1/59,6; distillato 0,5 M 46,0/65,4; IRM 2 M media 45,5/66,2; IRM 2 M
% distillato media 48,2/68,7; relazioni locali 0,5 M 45,3/61,6). Letteratura verificata sulle versioni arXiv il
% 19/11/2026: SMV-EAR tab. 2 (anche 62,9 senza augmentation), HoloEv-Net tab. I, Norrby et al. tab. 1 (le tre
% configurazioni pre-addestrate su DailyDVS-200 con ViT-Tiny raggiungono 75,9, 78,5 e 78,7; addestrato da zero il ViT-Tiny
% si ferma a 19,5-21,3), TVTA tab. 1, RN-Net tab. 4. Il numero di parametri di MSTP non e' verificabile e non e' riportato.
% Split: THU training ufficiale 1.786 clip, di cui 399 usate per la validazione (1.387 di training); DVS-Lip 14.896, di
% cui 2.995 di validazione (11.901 di training). Figura: figure_scripts/cap6/fig_stato_arte.py.

\section{I modelli finali e lo stato dell'arte}
\label{sec:configurazioni_finali}

Le sezioni precedenti hanno fissato, sempre sulla validazione di sviluppo, le configurazioni a cui arriva il lavoro (tabella~\ref{tab:configurazioni_finali}). Su THU-EACT-50-CHL sono l'IRM antisimmetrico alle tre larghezze, addestrato con la ricetta standard, e la sua versione compressa. Su DVS-Lip sono l'IRM distillato, con cinque insegnanti compatti nel modello da mezzo milione di parametri e con quattro insegnanti in quello da due milioni. Con le configurazioni fissate, il test ufficiale permette di collocare i risultati rispetto ai lavori pubblicati.

\begin{table}[t]
\centering
\small
\renewcommand{\arraystretch}{1.2}
\caption[Configurazioni finali sul test ufficiale]{Configurazioni finali sul test ufficiale, valutato a posteriori. Valori in percentuale. Con due milioni di parametri, media e deviazione standard di tre addestramenti. Su DVS-Lip il modello compatto distillato non si lascia comprimere entro la tolleranza (sezione~\ref{sec:compressione}).}
\label{tab:configurazioni_finali}
\begin{tabular}{@{}l l r r r@{}}
\hline
\textbf{Budget} & \textbf{Configurazione} & \textbf{Parametri} & \textbf{Accuratezza} & \textbf{Macro-F1} \\
\hline
\multicolumn{5}{@{}l}{\emph{THU-EACT-50-CHL}} \\
\quad $0{,}5$\,M & IRM compresso & $456\,698$ & $66{,}0$ & $59{,}7$ \\
\quad 1\,M & IRM & $1\,123\,258$ & $71{,}9$ & $65{,}8$ \\
\quad & IRM compresso & $1\,006\,522$ & $70{,}4$ & $64{,}4$ \\
\quad 2\,M & IRM & $1\,981\,626$ & $74{,}3\pm1{,}0$ & $68{,}2\pm1{,}1$ \\
\quad & IRM compresso & $1\,750\,202$ & $74{,}2\pm1{,}3$ & $68{,}2\pm1{,}4$ \\
\multicolumn{5}{@{}l}{\emph{DVS-Lip}} \\
\quad $0{,}5$\,M & IRM distillato & $513\,324$ & $55{,}7$ & $55{,}0$ \\
\quad 2\,M & IRM distillato & $1\,994\,476$ & $58{,}4\pm0{,}2$ & $57{,}9\pm0{,}2$ \\
\quad & IRM distillato e compresso & $1\,814\,252$ & $58{,}3\pm0{,}2$ & $57{,}8\pm0{,}2$ \\
\hline
\end{tabular}
\end{table}

\titoletto{THU-EACT-50-CHL}
La tabella~\ref{tab:stato_arte_thu} raccoglie i risultati pubblicati sul test ufficiale di THU-EACT-50-CHL che abbiamo trovato, e la figura~\ref{fig:stato_arte} li dispone in funzione del numero di parametri. I lavori riportano soltanto l'accuratezza, e il confronto usa quindi questa misura.

\begin{table}[t]
\centering
\small
\renewcommand{\arraystretch}{1.15}
\caption[Stato dell'arte su THU-EACT-50-CHL]{Accuratezza sul test ufficiale di THU-EACT-50-CHL. I valori della letteratura sono quelli riportati dalle fonti indicate, non riprodotti. Il pre-training è indicato soltanto quando usa dati diversi dal training di THU-EACT-50-CHL. Per i modelli di questa tesi con due milioni di parametri, media e deviazione standard di tre addestramenti.}
\label{tab:stato_arte_thu}
\begin{adjustbox}{max width=\textwidth}
\begin{tabular}{@{}l l r l r c@{}}
\hline
\textbf{Metodo} & \textbf{Tipo} & \textbf{Parametri} & \textbf{Pre-training} & \textbf{Acc. [\%]} & \textbf{Fonte} \\
\hline
HMAX SNN & SNN & n.d. & no & $32{,}7$ & \cite{fan2026smvear} \\
Motion SNN & SNN & n.d. & no & $47{,}3$ & \cite{fan2026smvear} \\
EV-ACT & ANN & $21{,}3$\,M & no & $58{,}5$ & \cite{gao2023action,fan2026smvear} \\
EventMG & ANN & $0{,}77$\,M & no & $58{,}8$ & \cite{fan2026smvear} \\
EventMamba & ANN & $0{,}91$\,M & no & $59{,}4$ & \cite{fan2026smvear} \\
MVFNet & ANN & $23{,}6$\,M & no & $59{,}6$ & \cite{hao2026holoev} \\
EventBind+POKER & ANN, eventi e testo & $635$\,M & sì, eventi e testo & $63{,}7$ & \cite{norrby2026event} \\
SMV-EAR & ANN & $23{,}5$\,M & no & $66{,}7$ & \cite{fan2026smvear} \\
HoloEv-Net-Small & ANN & $4{,}4$\,M & no & $67{,}1$ & \cite{hao2026holoev} \\
ViT-Base, MAE sul solo THU & ANN & $85{,}8$\,M & no & $68{,}6$ & \cite{norrby2026event} \\
HoloEv-Net-Base & ANN & $12{,}5$\,M & no & $69{,}9$ & \cite{hao2026holoev} \\
ViT-Tiny & ANN & $5{,}5$\,M & sì, DailyDVS-200 & $78{,}7$ & \cite{norrby2026event} \\
\hline
\multicolumn{6}{@{}l}{\emph{Questo lavoro, SNN, nessun pre-training}} \\
\quad modello di riferimento & SNN & $0{,}49$\,M & no & $56{,}1$ & \\
\quad IRM compresso & SNN & $0{,}46$\,M & no & $66{,}0$ & \\
\quad modello di riferimento & SNN & $1{,}10$\,M & no & $68{,}8$ & \\
\quad IRM compresso & SNN & $1{,}01$\,M & no & $70{,}4$ & \\
\quad IRM & SNN & $1{,}12$\,M & no & $71{,}9$ & \\
\quad modello di riferimento & SNN & $1{,}96$\,M & no & $72{,}3\pm0{,}9$ & \\
\quad IRM compresso & SNN & $1{,}75$\,M & no & $74{,}2\pm1{,}3$ & \\
\quad IRM & SNN & $1{,}98$\,M & no & $74{,}3\pm1{,}0$ & \\
\hline
\end{tabular}
\end{adjustbox}
\end{table}

\begin{figure}[t]
    \centering
    \includegraphics[width=\textwidth]{images/cap6/cap6_stato_arte.pdf}
    \caption[Stato dell'arte su THU-EACT-50-CHL]{Accuratezza sul test ufficiale di THU-EACT-50-CHL in funzione dei parametri, in scala logaritmica, per i modelli di questa tesi e per i metodi pubblicati della tabella~\ref{tab:stato_arte_thu}. I triangoli indicano un pre-training su dati esterni a THU-EACT-50-CHL. Le due SNN pubblicate non dichiarano il numero di parametri, e la migliore è indicata dalla linea punteggiata.}
    \label{fig:stato_arte}
\end{figure}

Fra le reti spiking, la migliore pubblicata su questo dataset che abbiamo trovato arriva al $47{,}3\%$, e la rete con l'IRM compresso, con $0{,}46$ milioni di parametri, la supera di quasi diciannove punti.

Fra tutti i metodi che non usano dati esterni il migliore è HoloEv-Net-Base~\cite{hao2026holoev}, che con $12{,}5$ milioni di parametri arriva al $69{,}85\%$. La rete con l'IRM lo supera già con un milione di parametri, sia nella versione piena sia in quella compressa, e con due milioni lo supera in tutti e tre gli addestramenti ripetuti, di oltre quattro punti in media, con circa sei volte meno parametri. Nella fascia compatta, sotto il milione di parametri, il vantaggio è ancora più netto, perché le due reti ANN più piccole, EventMamba ed EventMG, restano sotto il $60\%$ e il modello compresso da $0{,}46$ milioni le supera di oltre sei punti, avvicinandosi a SMV-EAR~\cite{fan2026smvear}, che ne usa cinquanta volte tanti.\footnote{Il $66{,}7\%$ di SMV-EAR è ottenuto con augmentation. Senza, lo stesso lavoro riporta il $62{,}9\%$. Le configurazioni finali di questa tesi su THU-EACT-50-CHL non usano augmentation.}

Un solo lavoro va oltre i nostri modelli, quello di Norrby e coautori~\cite{norrby2026event}, che pre-addestra in modo auto-supervisionato un ViT-Tiny su DailyDVS-200, un dataset di azioni a eventi molto più grande, e lo rifinisce su THU-EACT-50-CHL, arrivando al $78{,}7\%$. Lo stesso modello addestrato da zero si ferma intorno al $20\%$, e un ViT-Base pre-addestrato sul solo training di THU-EACT-50-CHL al $68{,}6\%$. Il pre-training su dati esterni è quindi la leva che fa la differenza, e il confronto va dichiarato per quello che è..

Ai risultati contribuiscono la rete compatta riallocata del capitolo~\ref{chap:riallocazione_budget}, il clock a passi di durata fisica costante del capitolo~\ref{chap:clock_input}, la capacità e, infine, la memoria relazionale, e sono quindi una proprietà dell'intera catena, non del solo modulo. Il ruolo dell'IRM è decisivo soprattutto a basso budget, dove con mezzo milione di parametri porta la rete dal $56{,}1$ al $66{,}0\%$. Con un milione è ciò che la porta sopra il $69{,}85\%$, perché il modello di riferimento della stessa larghezza resta al $68{,}8\%$. Con due milioni anche il modello di riferimento supera lo stato dell'arte senza dati esterni, con il $72{,}3\%$, e l'IRM vi aggiunge circa due punti, con lo stesso segno in tutti e tre gli addestramenti.

\titoletto{DVS-Lip}
Su DVS-Lip il quadro è diverso, e il confronto con la letteratura serve da contesto più che da gara (tabella~\ref{tab:stato_arte_lip}). I valori assoluti restano lontani dai migliori lavori pubblicati, anche fra le reti spiking, dove SpikingGRU2+ arriva al $75{,}3\%$. Parte della distanza ha ragioni di protocollo, perché i lavori più recenti addestrano con un'augmentation ricca, fatta di ritagli casuali, cancellazioni, zoom e mascheramenti temporali, che qui non è stata usata. Queste differenze non bastano però a spiegare tutto il divario, e non vanno presentate come una giustificazione. In questa tesi DVS-Lip ha avuto soprattutto il ruolo di un caso diagnostico, che ha mostrato quando una memoria non basta e quale leva invece funziona. Come in tutta la letteratura, il gruppo delle parole visivamente simili resta il punto più debole, con un'accuratezza di circa venti punti inferiore a quella delle altre parole.

\begin{table}[t]
\centering
\small
\renewcommand{\arraystretch}{1.15}
\caption[DVS-Lip, test ufficiale]{Accuratezza sul test ufficiale di DVS-Lip, complessiva e sui due gruppi di parole, quelle visivamente simili (Acc1) e le altre (Acc2). I valori della letteratura sono quelli riportati dalle fonti indicate. Per i modelli di questa tesi con due milioni di parametri, media di tre addestramenti.}
\label{tab:stato_arte_lip}
\begin{adjustbox}{max width=\textwidth}
\begin{tabular}{@{}l l r r r c@{}}
\hline
\textbf{Metodo} & \textbf{Tipo} & \textbf{Parametri} & \textbf{Acc. [\%]} & \textbf{Acc1 / Acc2 [\%]} & \textbf{Fonte} \\
\hline
Spiking MSTP & SNN & n.d. & $60{,}2$ & -- & \cite{zheng2026tvta} \\
RN-Net & ibrida & $6{,}0$\,M & $67{,}5$ & -- & \cite{yoo2023rnnet} \\
MSTP & ANN & n.d. & $72{,}1$ & $62{,}2$ / $82{,}1$ & \cite{tan2022dvslip,zheng2026tvta} \\
MTGA & ANN & n.d. & $75{,}1$ & $63{,}9$ / $86{,}4$ & \cite{zheng2026tvta} \\
SpikingGRU2+ & SNN & n.d. & $75{,}3$ & -- & \cite{zheng2026tvta} \\
HFR-Lip & ANN & n.d. & $77{,}1$ & $66{,}5$ / $87{,}7$ & \cite{zheng2026tvta} \\
TVTA & ANN & n.d. & $77{,}5$ & $67{,}2$ / $87{,}8$ & \cite{zheng2026tvta} \\
\hline
\multicolumn{6}{@{}l}{\emph{Questo lavoro, SNN}} \\
\quad IRM & SNN & $0{,}51$\,M & $50{,}8$ & $42{,}1$ / $59{,}6$ & \\
\quad IRM e relazioni locali & SNN & $0{,}52$\,M & $53{,}4$ & $45{,}3$ / $61{,}6$ & \\
\quad IRM distillato & SNN & $0{,}51$\,M & $55{,}7$ & $46{,}0$ / $65{,}4$ & \\
\quad IRM & SNN & $1{,}99$\,M & $55{,}8$ & $45{,}5$ / $66{,}2$ & \\
\quad IRM distillato & SNN & $1{,}99$\,M & $58{,}4$ & $48{,}2$ / $68{,}7$ & \\
\hline
\end{tabular}
\end{adjustbox}
\end{table}

\titoletto{Cautele di protocollo}
Il confronto sul test ufficiale richiede alcune cautele, a cominciare dal fatto che i modelli di questa tesi sono addestrati sullo split di sviluppo, che su THU-EACT-50-CHL contiene $1\,387$ clip invece delle $1\,786$ del training ufficiale, e nessuno dei lavori citati dichiara di aver ridotto il proprio. Il confronto è quindi, semmai, sfavorevole a noi. Le configurazioni finali su THU-EACT-50-CHL non usano augmentation, pre-training o dati esterni. I risultati con mezzo e con un milione di parametri vengono da un solo addestramento, quelli con due milioni sono la media di tre. Il test è stato valutato a posteriori, su modelli scelti sulla validazione, e non ha guidato alcuna scelta. Le cifre della letteratura, infine, sono quelle riportate dai rispettivi lavori, che non abbiamo riprodotto.


% ==============================================================================
%  Dove sbagliano ancora i modelli finali
% ==============================================================================
% Fonti (2026-11-19): figure_scripts/cap6/fig_errori_lip.py (test ufficiale, logit di b_minus_final_inference; gruppi e
% coppie da configs/dvslip_class_groups.json). Errori sul test (scambi nella coppia / altri errori sulle parole delle
% coppie / errori sulle altre parole): riferimento 442/1010/1035, IRM 478/965/1004, IRM distillato 515/832/859, IRM 2 M
% 481/879/840, IRM 2 M distillato 500/792/778 (media di tre addestramenti), relazioni locali 488/877/953. Scambi nelle
% coppie piu' difficili per l'IRM 2 M distillato (tre addestramenti): price/press 39,7%, spend/spent 36,7%,
% benefit/benefits 36,4%, happen/happened 33,3%; terms/times 4,0%, election/action 3,7%. Durata: quartili del numero di
% passi validi (validazione: soglie 36/40/44; test: 38/41/44, passi da c0_ri_diagnostics/workers/dvslip_c0); la parola
% spiega il 39% della varianza della durata sul test; parole delle coppie nel quartile piu' breve 58,2%, nel piu' lungo
% 37,1%; il 95% delle clip del test ha almeno 33 passi (minimo 24). THU: classi da risultati6.classi_thu (validazione),
% durate da E1.csv. Ritmo: artifacts/reports/c0_e0_vc/<dataset>/<E0|VC>/{clean,time_scale_0.5,time_scale_2}_summary.json
% (modello di riferimento, sola inferenza); addestramento con perturbazioni della scala temporale: run
% VC_RI_M--interval-antisymmetric-r16-svc (DVS-Lip 58,60 F1, p=0,155; THU 60,30 F1, p=0,085).

\section{Dove sbagliano ancora i modelli finali}
\label{sec:errori}

Le sezioni precedenti hanno mostrato quanto si possa guadagnare con la memoria relazionale, la capacità e la distillazione. Questa sezione guarda invece a ciò che resta, cioè agli errori che nessuna di queste leve corregge. Su DVS-Lip riguardano soprattutto due aspetti, le coppie di parole simili e le clip più brevi. Su entrambi i dataset ne riguardano un terzo, il ritmo con cui i movimenti vengono eseguiti.

\titoletto{Le coppie di parole simili}
Metà delle parole di DVS-Lip forma 25 coppie visivamente simili, come \emph{spend} e \emph{spent} o \emph{price} e \emph{press}, sulle quali si calcola Acc1 (sezione~\ref{subsec:dvslip}). La figura~\ref{fig:confusione_lip}a mostra gli errori del miglior modello finale sul test ufficiale, con le parole disposte in modo che le due di ciascuna coppia siano adiacenti. Gli errori fra parole simili si raccolgono in piccoli blocchi lungo la diagonale, dove una parola viene presa per la sua compagna. Nelle coppie più difficili accade a più di un terzo delle clip, per esempio con \emph{price} e \emph{press}, \emph{spend} e \emph{spent}, \emph{benefit} e \emph{benefits} o \emph{happen} e \emph{happened}, mentre altre coppie, come \emph{terms} e \emph{times}, quasi non si confondono. Molte delle coppie difficili differiscono soltanto per una desinenza, cioè per un movimento breve e poco visibile alla fine della parola.

\begin{figure}[t]
    \centering
    \includegraphics[width=\textwidth]{images/cap6/cap6_confusione_lip.pdf}
    \caption[Errori sulle coppie di parole simili di DVS-Lip]{Errori sulle coppie di parole simili di DVS-Lip, sul test ufficiale. (a)~Errori dell'IRM con due milioni di parametri e distillazione, sommati sui tre addestramenti. Ogni riga è una parola pronunciata, e il colore indica la percentuale delle sue clip assegnate a ciascuna parola predetta. La diagonale, cioè le risposte corrette, è esclusa e in grigio. Le due parole di ciascuna coppia sono adiacenti, e gli scambi fra loro formano i blocchi lungo la diagonale nel riquadro in alto a sinistra. (b)~Clip sbagliate lungo il percorso dei modelli, divise in scambi fra le due parole di una coppia, altri errori sulle parole delle coppie ed errori sulle altre parole. Per i modelli con due milioni di parametri, media di tre addestramenti.}
    \label{fig:confusione_lip}
\end{figure}

Lungo il percorso dei modelli (figura~\ref{fig:confusione_lip}b), distillazione e capacità riducono di circa un quarto gli errori verso parole diverse dalla compagna, sia fra le parole delle coppie sia fra le altre. Gli scambi dentro la coppia, invece, non scendono, e passano anzi da 442 nel modello di riferimento a circa 500 nei modelli migliori. Via via che la rete impara a riconoscere a quale coppia appartenga una parola, una quota crescente dei suoi errori residui si concentra proprio sulla scelta fra le due parole della coppia. Nemmeno le relazioni locali della sezione~\ref{sec:relazioni_locali}, pensate per i dettagli fini, riducono gli scambi, e anche il loro guadagno viene dagli altri errori. Il conteggio stabile non va letto come un limite invalicabile, perché i migliori lavori pubblicati superano su queste parole l'accuratezza dei nostri modelli di circa venti punti (tabella~\ref{tab:stato_arte_lip}). Indica però che le leve provate in questa tesi non toccano la distinzione più fine che il compito richiede.

\titoletto{Le clip più brevi}
La seconda difficoltà riguarda la durata, e dividendo le clip in quattro gruppi secondo il numero di passi, dalle più brevi alle più lunghe, le prime risultano le più difficili per tutti i modelli (figura~\ref{fig:durata}). Sul test il modello di riferimento passa dal $40\%$ di accuratezza nel gruppo più breve al $64\%$ in quello più lungo. Ogni leva alza tutti i gruppi di una quantità simile, e il modello finale arriva al $50$ e al $71\%$, ma il divario resta di circa venti punti sul test e di quattordici in validazione.

\begin{figure}[t]
    \centering
    \includegraphics[width=\textwidth]{images/cap6/cap6_durata.pdf}
    \caption[Accuratezza su DVS-Lip per durata delle clip]{Accuratezza su DVS-Lip per durata delle clip, in quattro gruppi delimitati dai quartili del numero di passi, dal più breve al più lungo, lungo il percorso dei modelli. Il grigio indica il modello di riferimento, le tinte chiare e scure i modelli con mezzo milione e con due milioni di parametri, il tratteggio la distillazione. Per i modelli con due milioni di parametri, media di tre addestramenti.}
    \label{fig:durata}
\end{figure}

Il divario non nasce da una mancanza di passato, perché quasi tutte le clip, anche nel gruppo più breve, durano più del doppio dell'orizzonte delle bande, e la memoria vi dispone presto di una storia completa. Durata e parola, inoltre, sono legate, tanto che la parola pronunciata spiega circa il $40\%$ della variabilità della durata, e le parole delle coppie simili sono il $58\%$ delle clip nel gruppo più breve contro il $37\%$ in quello più lungo. Le clip brevi sono quindi difficili anche per le parole che contengono, e la prova non permette di separare i due effetti.

\titoletto{THU-EACT-50-CHL}
Su THU-EACT-50-CHL gli errori sono più dispersi, e lo è anche il guadagno della memoria relazionale. In validazione l'IRM migliora l'F1 di 35 delle 46 classi presenti, e il guadagno non è confinato a una fascia di durate. Cresce con la lunghezza delle clip in validazione e si concentra su quelle brevi nel test, e con poche decine di clip per gruppo la sola conclusione che regge su entrambe le partizioni è che non dipende dalla durata.

\titoletto{Il ritmo}
\label{sec:ritmo}
La rappresentazione a lunghezza variabile fissa la durata fisica di un passo, e con essa il significato dei ritardi. Lo stesso movimento eseguito più in fretta o più lentamente occupa però un numero diverso di passi, e il capitolo~\ref{chap:clock_input} ha lasciato aperta la domanda su quanto la rete tolleri questa variazione. La prova moltiplica i tempi degli eventi di ogni clip di validazione per $0{,}5$, come se il movimento fosse eseguito al doppio della velocità, oppure per 2, come se fosse eseguito a metà velocità, e valuta in sola inferenza il modello di riferimento addestrato con la finestra fissa e con la lunghezza variabile.

\begin{table}[t]
\centering
\small
\renewcommand{\arraystretch}{1.2}
\caption[Accuratezza con i tempi degli eventi riscalati]{Accuratezza [\%] del modello di riferimento sulla validazione di sviluppo, con i tempi degli eventi moltiplicati per $0{,}5$ (movimento due volte più veloce) o per 2 (due volte più lento), in sola inferenza. Con la finestra fissa la clip rallentata perde circa il $6\%$ degli eventi, che cadono oltre la finestra, con la lunghezza variabile nessuno.}
\label{tab:ritmo}
\begin{tabular}{@{}l r r r r r r@{}}
\hline
 & \multicolumn{3}{c}{\textbf{DVS-Lip}} & \multicolumn{3}{c}{\textbf{THU-EACT-50-CHL}} \\
\textbf{Rappresentazione} & \textbf{originale} & $\boldsymbol{\times 0{,}5}$ & $\boldsymbol{\times 2}$ & \textbf{originale} & $\boldsymbol{\times 0{,}5}$ & $\boldsymbol{\times 2}$ \\
\hline
finestra fissa & $55{,}6$ & $13{,}3$ & $11{,}9$ & $50{,}4$ & $15{,}0$ & $19{,}0$ \\
lunghezza variabile & $56{,}8$ & $9{,}4$ & $11{,}0$ & $61{,}2$ & $15{,}3$ & $13{,}8$ \\
\hline
\end{tabular}
\end{table}

Con entrambe le rappresentazioni l'accuratezza crolla fra il $9$ e il $19\%$, qualunque sia il verso del cambiamento (tabella~\ref{tab:ritmo}). Con la finestra fissa la clip rallentata viene anche troncata, ma con la lunghezza variabile tutti gli eventi restano, e il crollo è lo stesso. La rete riconosce i movimenti alle velocità viste in addestramento, e un fattore due basta a portarla fuori da ciò che conosce, sicché non è invariante al ritmo. Il limite riguarda il modo in cui la rete ha appreso più che il clock, che rende costante la durata fisica del passo ma non insegna a riconoscere la stessa azione a velocità diverse. Addestrare l'IRM con piccole perturbazioni casuali della scala temporale, entro il $10\%$, non cambia in modo significativo la validazione standard, e l'effetto di questo addestramento su riscalamenti più ampi non è stato misurato.

Il percorso del capitolo lascia così tre limiti che nessuna delle leve provate ha scalfito, gli scambi fra le parole più simili, la difficoltà delle clip più brevi su DVS-Lip e l'assenza di invarianza al ritmo. Insieme ai risultati raccolti finora, sono il punto da cui partono le conclusioni.


% Conclusioni, non numerate come l'introduzione.
\chapter*{Conclusioni}
\label{chap:conclusioni}
\phantomsection
\addcontentsline{toc}{chapter}{Conclusioni}
\markboth{Conclusioni}{Conclusioni}

% Riscrittura del 2026-11-19 sulla nuova struttura dei capitoli 5 e 6 (docs/STRUTTURA_NARRATIVA_E_RISULTATI_CAPITOLI_5_6_7.md,
% sezione 6). Testo unico senza suddivisioni, non numerato come l'introduzione (richiesta dell'utente). Numeri ripresi dai
% capitoli 3-6, stesse fonti. DA AGGIORNARE all'esito dello scan delle relazioni locali prima del pooling (Legion
% 2123991-2123995): il paragrafo sulle prospettive del modello lo cita come verifica in corso.

Questa tesi è partita da una domanda semplice da porre e meno semplice da affrontare, cioè come convenga spendere il budget di uno Spiking Transformer con circa mezzo milione di parametri perché colga le relazioni temporali che contano nei dati a eventi, senza perdere di vista quanto costi ricordare. La risposta non è venuta da un singolo modulo, ma da una sequenza di scelte, ciascuna misurata cambiando un solo elemento alla volta, e conviene ripercorrerla nell'ordine in cui si è costruita.

La prima scelta ha riguardato la baseline, una versione compatta di QKFormer in cui quasi tutto il costo si concentrava nei due patch embedding, mentre la memoria temporale era affidata ai soli potenziali di membrana. Spostare risorse dal frontend a una memoria esplicita, con un frontend piramidale, un mixer causale a più ritardi e una convoluzione depthwise al posto dell'attenzione del primo stadio, ha portato la Macro-F1 di DVS-Lip da 44 a 55 punti con appena 320 parametri in più e con il $42\%$ di stato persistente in meno. Il vantaggio si è ripetuto sui quattro dataset, e la baseline non lo raggiunge nemmeno con un budget quattro volte più grande. In un modello compatto conta dove la capacità viene spesa, prima ancora di quanta se ne aggiunga, ed è questa la prima indicazione del lavoro, da cui in qualche misura discendono tutte le altre.

Il modello così ottenuto contava il tempo in passi, e la seconda scelta ha riguardato proprio il significato di un passo. Sostituire la finestra fissa con una rappresentazione a lunghezza variabile, in cui ogni passo ha la stessa durata fisica in tutto il dataset, ha fatto guadagnare oltre nove punti su THU-EACT-50-CHL e circa otto su DVS128 Gesture, un effetto paragonabile a quello di una modifica dell'architettura. Un controllo con un numero fisso di passi ha mostrato che eliminare le code vuote della finestra fissa non basta, e che dove le durate delle registrazioni variano molto conta la costanza fisica del passo. Il clock dell'input, spesso trattato come un dettaglio di preparazione dei dati, si è rivelato una variabile di progetto a tutti gli effetti, e ha reso possibile ciò che è venuto dopo, perché con un passo di durata nota i ritardi della rete diventano distanze precise nel passato.

Su quelle distanze poggia il contributo centrale, nato dall'osservazione che il mixer a ritardi del modello di riferimento richiamava il passato e lo sommava al presente, senza mai confrontarli. L'IRM antisimmetrico riassume il passato in bande a distanze fisiche fisse, tiene conto di quanta storia esista davvero e confronta ciascun riassunto con il presente attraverso un termine bilineare a basso rango, che cambia segno quando i due istanti si scambiano, sulla falsariga del correlatore di Hassenstein e Reichardt. Su THU-EACT-50-CHL, con mezzo milione di parametri, supera il modello di riferimento di quasi otto punti di Macro-F1 in validazione e di quasi undici sul test ufficiale. Né un sito vuoto né un percettrone con gli stessi parametri riproducono il guadagno, e le stesse bande senza confronto restano al livello del riferimento. Il merito è del confronto, che nella forma antisimmetrica rende più che in quella simmetrica, applicato a un passato collocato a distanze fisiche precise. Tracce esponenziali, un orologio appreso e un passato più remoto, provati al posto delle bande, perdono tutti terreno. Gli interventi sul modello addestrato raccontano la stessa storia dall'interno, perché la rete si appoggia in modo massiccio al confronto, soprattutto con le bande più lontane dal presente, e il suo vantaggio svanisce non appena l'ordine dei passi viene alterato.

Su DVS-Lip, invece, lo stesso modulo sposta il risultato di poco più di mezzo punto, e i controlli quasi non si distinguono fra loro, non per un difetto della memoria ma perché con mezzo milione di parametri i due dataset si sono rivelati frenati da limiti opposti. La rete per THU-EACT-50-CHL non riesce nemmeno a descrivere i propri dati di addestramento, e ogni variante che la aiuta ad adattarsi migliora anche la validazione. Quella per DVS-Lip impara quasi per intero il training ma ne generalizza poco più della metà, e due modelli di accuratezza quasi identica danno la stessa risposta su poco più di metà delle clip. Anche tre addestramenti della stessa architettura, con due milioni di parametri, concordano soltanto su due clip su tre. Su THU-EACT-50-CHL la memoria relazionale fornisce alla rete la capacità temporale che le manca. Su DVS-Lip il limite è la stabilità della soluzione, che nessuna memoria del secondo patch embedding cambia.

Su DVS-Lip, come la diagnosi suggeriva, la prima leva efficace è stata la distillazione, che raccoglie in un solo modello compatto ciò che varianti diverse sanno separatamente. Porta l'IRM da $57{,}5$ a $60{,}9$ di Macro-F1 in validazione e ne alza di quasi tre punti e mezzo l'accuratezza sul test, senza alcun costo in inferenza, e rende più simili fra loro le soluzioni di addestramenti ripetuti. Un bersaglio soltanto più morbido, senza l'informazione di altri modelli, non basta a ottenere lo stesso effetto. La seconda leva è nata da un dubbio che la diagnosi lasciava aperto, cioè che su DVS-Lip una parte dell'informazione utile sparisse nel pooling prima di raggiungere la memoria. Lo stesso confronto antisimmetrico dell'IRM, calcolato prima del pooling fra celle vicine e con ritardi di poche decine di millisecondi, è il correlatore di Hassenstein e Reichardt quasi alla lettera, e migliora l'IRM di due punti e mezzo sul test senza alcun guadagno su THU-EACT-50-CHL. Si somma quasi per intero alla distillazione, perché agisce su una causa diversa, e insieme le due leve portano il modello compatto dal $50{,}8$ al $56{,}8\%$ sul test. Su THU-EACT-50-CHL, dove la diagnosi indica un limite di adattamento e non di varianza, la configurazione resta l'IRM così com'è. In un modello compatto architettura e addestramento vanno scelti insieme al compito.

Oltre il regime compatto entrambi i dataset migliorano con la capacità, e su THU-EACT-50-CHL il margine dell'IRM sul riferimento della stessa larghezza scende da quasi dieci punti con mezzo milione di parametri a tre con un milione e a due con due milioni, perché con abbastanza capacità anche il riferimento riesce ad adattarsi ai dati. La memoria relazionale conta soprattutto dove il budget è scarso, che è proprio il regime per cui è stata pensata. Su DVS-Lip capacità, relazioni locali e distillazione si sommano, e il modello migliore, con due milioni di parametri, arriva al $60{,}2\%$ sul test e al $67{,}8\%$ di Macro-F1 in validazione, i valori più alti ottenuti su questo dataset. La distillazione aiuta allo stesso modo anche il riferimento, ma a parità di ricetta le relazioni locali restano avanti di oltre due punti, mentre l'IRM da solo non se ne distingue più. Il costo della memoria, infine, si lascia ridurre dopo l'addestramento, poiché le quattro matrici con cui le bande vengono riportate nei canali si esprimono in una base comune di rango ridotto, scelto con una regola fissata in anticipo, senza riaddestrare la rete e senza perdite misurabili sul test dei modelli finali. Così compresso, l'IRM costa meno del mixer a ritardi che sostituisce a ogni larghezza su THU-EACT-50-CHL, fino a circa un quinto del suo costo con due milioni di parametri, e su DVS-Lip con uno e due milioni, mentre il modello compatto distillato ne resta di poco sopra. È, in piccolo, la stessa lezione da cui il lavoro era partito. Il ramo locale, invece, si paga, e la configurazione finale di DVS-Lip esegue fra l'8 e il $23\%$ di moltiplicazioni in più del modello di riferimento, in cambio di cinque o sette punti sul test.

Su THU-EACT-50-CHL questa catena di scelte porta oltre lo stato dell'arte, e già con $0{,}46$ milioni di parametri la rete con l'IRM compresso supera di quasi diciannove punti la migliore rete spiking pubblicata. Con uno e con due milioni supera tutti i metodi che non usano dati esterni, fino al $74{,}3\%$ di accuratezza media sul test, contro il $69{,}85\%$ del migliore di essi, che usa sei volte più parametri. Soltanto un modello pre-addestrato su un dataset esterno molto più ampio va ancora oltre. Il risultato appartiene all'intera catena, dalla rete riallocata al clock, alla capacità e alla memoria, ma l'IRM vi è decisivo proprio a basso budget, dove porta la rete compatta dal $56$ al $66\%$ e, con un milione di parametri, oltre i metodi pubblicati che il riferimento della stessa larghezza non raggiunge. Su DVS-Lip il modello finale arriva allo stesso $60{,}2\%$ di Spiking MSTP, una delle due reti spiking pubblicate su questo dataset, ma i valori assoluti restano lontani dai migliori lavori, e il dataset ha avuto soprattutto il ruolo di un caso diagnostico, che ha mostrato quando una memoria a lungo termine non basta e quali leve funzionano al suo posto.

I confini di questi risultati vanno dichiarati con chiarezza, a cominciare dal fatto che quasi tutti i confronti sull'architettura si basano su un solo addestramento per configurazione, e soltanto con due milioni di parametri l'addestramento è stato ripetuto. Il confronto appaiato tiene conto di quali campioni la validazione contenga, ma non della variabilità tra un addestramento e l'altro, che su DVS-Lip è ampia, e i valori di $p$ non sono corretti per i molti confronti eseguiti. Le differenze nette osservate su THU-EACT-50-CHL reggono a questa cautela, quelle di un punto su DVS-Lip no. Modelli, varianti e ranghi di compressione sono stati scelti sulla validazione di sviluppo. Il test ufficiale, aperto soltanto a posteriori, prima per diagnosticare modelli già scelti e poi per confrontare quelli finali con la letteratura, offre una conferma importante, ma non la misura su dati mai osservati che darebbe un protocollo più rigido. I modelli, inoltre, sono addestrati su una parte del training ufficiale, e le cifre della letteratura sono quelle riportate dai rispettivi lavori, non riprodotte. Il beneficio della memoria relazionale dipende dal compito e dalla scala, netto dove la rete compatta fatica ad adattarsi ai dati e il passato offre qualcosa da confrontare, piccolo su DVS-Lip nel secondo patch embedding, dove serve invece il confronto locale prima del pooling, assente su DVS128 Gesture e CIFAR10-DVS nelle prove svolte. Il ramo locale, a sua volta, è stato provato su THU-EACT-50-CHL soltanto nel modello compatto. Anche le alternative scartate lo sono soltanto nelle versioni provate, e nulla esclude che una memoria più lunga o un orologio adattivo, costruiti in un altro modo, possano funzionare. Quanto al costo, i valori riportati sono conteggi analitici di parametri, operazioni e stato, accompagnati da una stima aritmetica dell'energia. Indicano dove il costo si colloca, non quanto varrà una volta che il modello sarà realizzato.

Nessuna delle leve esaminate, poi, ha reso la rete invariante al ritmo, e un movimento eseguito al doppio o alla metà della velocità la porta ancora fuori da ciò che conosce. Su DVS-Lip, allo stesso modo, gli scambi fra le due parole di una coppia simile, come \emph{spend} e \emph{spent}, non diminuiscono con la memoria, le relazioni locali, la capacità o la distillazione, che riducono invece tutti gli altri errori. Le clip più brevi, infine, restano le più difficili, anche se in parte per le parole che contengono più che per la loro durata.

La realizzazione su hardware è la direzione più naturale per proseguire, ed è quella a cui il capitolo~\ref{chap:fondamenti_stato_arte} rimandava. Il modello compatto si presta per alcuni aspetti a un'implementazione su FPGA, dato che la memoria relazionale opera nel secondo patch embedding, dove le posizioni sono poche, e il suo stato è un buffer circolare degli ultimi passi, con ritardi fissi e quindi indirizzi noti in anticipo. Le medie di banda si aggiornano con poche somme e sottrazioni per passo, il confronto aggiunge una manciata di prodotti elemento per elemento, e la compressione porta il buffer al livello di quello del mixer a ritardi originale o al di sotto. Il ramo locale di DVS-Lip è la parte meno favorevole, perché conserva le proiezioni degli ultimi passi su tutta la mappa a piena risoluzione e aggiunge un quinto delle moltiplicazioni nel modello compatto. Restano però difficoltà che i conteggi non vedono, perché la memoria lavora su correnti a più bit e non su spike, e richiede quindi moltiplicazioni vere, e lo stato persistente va letto e riscritto a ogni passo, con accessi in memoria che possono pesare più dell'aritmetica. La rappresentazione a lunghezza variabile, infine, rende la latenza proporzionale alla durata della registrazione. Una sintesi effettiva, con la misura di risorse, latenza e potenza, e uno studio della precisione numerica con cui rappresentare pesi, correnti e buffer sono i passi che separano questi risultati da un sistema reale.

Anche sul piano del modello restano domande aperte, a cominciare dal costo del ramo locale, che oggi lavora sull'intera mappa a piena risoluzione e potrebbe forse conservare meno celle o meno componenti senza perdere il suo guadagno. Resta da vedere, inoltre, se il ramo aiuti THU-EACT-50-CHL con una rete più grande, dove il modello è stato provato soltanto in forma compatta. I modelli distillati, poi, si lasciano comprimere meno degli altri, e resta da capire se convenga distillare in un modello già fattorizzato o comprimere prima di distillare. L'assenza di invarianza al ritmo suggerisce di esporre la rete, durante l'addestramento, a velocità molto diverse, oppure di darle una memoria capace di leggere il passato su più scale insieme. Le scale stesse, oggi ereditate dai ritardi del mixer e riscalate per ciascun dataset, meriterebbero infine un protocollo che ne verifichi la trasferibilità da un compito all'altro.

Il percorso lascia un'indicazione che va oltre i singoli numeri, cioè che in uno Spiking Transformer compatto la capacità temporale non si ottiene allargando la rete né allungando la memoria, ma decidendo con cura che cosa sia un passo, quale passato conservare e come metterlo a confronto con il presente. Lo stesso confronto orientato, ispirato al correlatore di Hassenstein e Reichardt, si è rivelato utile in entrambi i compiti, ma alla scala che ciascuno richiede, su distanze di secondi per le azioni e fra celle vicine, su poche decine di millisecondi, per le labbra. Fatte queste scelte, una memoria piccola e locale basta a cambiare il risultato, e può costare meno di quella che sostituisce. Serve però anche saper riconoscere quando il limite non è la memoria ma la stabilità della soluzione, perché in quel caso la leva giusta sta nell'addestramento più che nell'architettura.

